% preamble.tex
\documentclass[11pt, twoside]{book}

% ------------------------------------------------------------
% Formato: A5, tascabile
% ------------------------------------------------------------
\usepackage{geometry}
\geometry{
  paperwidth=148mm,
  paperheight=210mm,
  top=16mm,
  bottom=18mm,
  inner=16mm,
  outer=13mm,
  headheight=14pt,
  headsep=8mm,
  footskip=12mm
}

% ------------------------------------------------------------
% Font: coppia pensata per la lettura prolungata su formato piccolo.
% TeX Gyre Pagella (clone della Palatino) come testo: occhio alto,
% pensato per restare leggibile a corpi piccoli, più caldo della
% Termes usata nella versione A4. TeX Gyre Heros per i titoli, come
% contrappunto lineare. Entrambe fanno parte della collezione
% TeX Gyre, inclusa in qualsiasi installazione TeX Live standard:
% nessun font da installare a parte.
% ------------------------------------------------------------
\usepackage{fontspec}
\setmainfont{TeX Gyre Pagella}[
  Numbers = OldStyle,
  Ligatures = TeX
]
\setsansfont{TeX Gyre Heros}
\setmonofont{TeX Gyre Cursor}[Scale=0.9]

\usepackage{polyglossia}
\setmainlanguage{italian}

\usepackage{amsmath, amssymb, mathtools}
\usepackage{siunitx}
\usepackage{graphicx}
\usepackage{xcolor}
\usepackage{booktabs}
\usepackage{longtable}
\usepackage{microtype}
% Tolleranza più ampia per la giustificazione: utile su una colonna
% stretta come quella di un formato A5, per assorbire le righe che
% altrimenti sporgerebbero leggermente oltre il margine.
\emergencystretch=3em
\tolerance=1000
\hbadness=3000
\usepackage{setspace}
\onehalfspacing

% Colore discreto per titoli e dettagli grafici (bruno-ocra, richiamo
% cromatico alla roccia e ai pigmenti citati nel testo, non decorativo)
\definecolor{ocra}{RGB}{123,62,32}

% ------------------------------------------------------------
% Intestazioni di pagina: nome parte sulle pagine pari,
% titolo capitolo sulle pagine dispari, numero di pagina in basso
% ------------------------------------------------------------
\usepackage{fancyhdr}
\pagestyle{fancy}
\fancyhf{}
\renewcommand{\headrulewidth}{0.4pt}
\renewcommand{\headrule}{\color{ocra}\hrule height\headrulewidth width\headwidth}
\fancyhead[LE]{\small\itshape\leftmark}
\fancyhead[RO]{\small\itshape\rightmark}
\fancyfoot[LE,RO]{\small\thepage}

\fancypagestyle{plain}{%
  \fancyhf{}
  \fancyfoot[LE,RO]{\small\thepage}
  \renewcommand{\headrulewidth}{0pt}
}

% Marcatori usati da fancyhdr: nome della parte a sinistra (pagine pari),
% titolo del capitolo a destra (pagine dispari). La classe book non
% definisce di suo un comando per marcare le parti nell'intestazione:
% lo si aggancia qui direttamente al comando \part.
\let\oldpart\part
\renewcommand{\part}[1]{\oldpart{#1}\markboth{#1}{}}
\renewcommand{\chaptermark}[1]{\markright{#1}{}}

% ------------------------------------------------------------
% Bibliografia
% ------------------------------------------------------------
\usepackage[sort&compress]{natbib}
\bibliographystyle{plainnat}
% Il testo del libro non contiene citazioni puntuali (\citep, \citet):
% e' un'opera divulgativa in cui le fonti sono richiamate discorsivamente
% nel testo, non con rimandi numerici. Percio' bibtex, senza istruzioni
% aggiuntive, non stamperebbe alcuna voce: nessuna fonte risulterebbe
% "citata" nel senso tecnico del comando \cite. \nocite{*} forza la
% stampa dell'intera bibliografia in references.bib, cosi' come una
% sezione "Fonti" a fine libro elenca tutto il materiale consultato,
% indipendentemente da singoli richiami puntuali nel testo.
\usepackage{etoolbox}
\AtBeginDocument{\nocite{*}}

\usepackage[colorlinks=true,
            linkcolor=ocra,
            citecolor=ocra,
            urlcolor=ocra,
            filecolor=ocra]{hyperref}

% Glossario semplice: comando per termini tecnici con rimando al glossario finale
\newcommand{\termine}[1]{\textit{#1}}

% ------------------------------------------------------------
% Stile dei capitoli e delle parti
% ------------------------------------------------------------
\usepackage{titlesec}

\titleformat{\chapter}[display]
  {\normalfont\sffamily\Large\bfseries\color{ocra}}
  {\chaptertitlename\ \thechapter}
  {6pt}
  {\normalfont\sffamily\LARGE\bfseries\color{black}}
\titlespacing*{\chapter}{0pt}{10pt}{24pt}

\titleformat{\part}[display]
  {\normalfont\sffamily\huge\bfseries\color{ocra}\centering}
  {\partname\ \thepart}
  {12pt}
  {\normalfont\sffamily\Huge\bfseries\color{black}\centering}

\titleformat{\section}
  {\normalfont\sffamily\large\bfseries}
  {\thesection}{0.8em}{}

\titleformat{\subsection}
  {\normalfont\sffamily\normalsize\bfseries}
  {\thesubsection}{0.8em}{}

% ------------------------------------------------------------
% Capolettera per l'apertura di ogni capitolo (tocco tipografico
% da libro tascabile curato, non un semplice muro di testo)
% ------------------------------------------------------------
\usepackage{tikz}
\usepackage{lettrine}
\setcounter{DefaultLines}{3}
\renewcommand{\DefaultLoversize}{0.1}
\renewcommand{\DefaultLraise}{0.0}
\newcommand{\aprecapitolo}[1]{\lettrine[lines=3, findent=2pt, nindent=0pt]{\color{ocra}#1}{}}

% Evita la lettera vietata "—": non caricare pacchetti che la introducano
% automaticamente (nessuna azione necessaria, promemoria per chi edita
% il sorgente a mano)


\title{Il silenzio della roccia\\[0.3em]\large Perché le montagne venete sembrano prive di arte rupestre}
\author{Mauro Barbieri}
\date{\today}

\begin{document}

\frontmatter

% ------------------------------------------------------------
% Copertina: disegnata con TikZ, coerente col formato tascabile.
% Fascia superiore e inferiore color ocra, titolo in maiuscoletto
% sans, un motivo lineare essenziale a richiamo di un'incisione
% rupestre stilizzata (nessuna immagine esterna necessaria).
% (il pacchetto tikz è caricato nel preambolo)
% ------------------------------------------------------------
\begin{titlepage}
  \thispagestyle{empty}
  \begin{tikzpicture}[remember picture, overlay]
    % fascia superiore
    \fill[ocra] (current page.north west) rectangle
      ([yshift=-45mm]current page.north east);
    % fascia inferiore
    \fill[ocra] ([yshift=45mm]current page.south west) rectangle
      (current page.south east);
    % figura antropomorfa stilizzata ispirata all'arte rupestre:
    % corpo distorto con braccia alzate e gambe divaricate,
    % centrata nella fascia inferiore ocra
    \begin{scope}[shift={([yshift=22mm]current page.south)}, white, line width=1.4pt, line cap=round, line join=round]
      % testa (cerchio leggermente ellittico, inclinato)
      \draw (0.4mm,28mm) ellipse (4.5mm and 5mm);
      % corpo (tronco corto, asimmetrico)
      \draw (0,23mm) -- (1mm,14mm);
      % braccio sinistro: alzato e piegato, con dita accennate
      \draw (0,23mm) -- (-9mm,30mm) -- (-12mm,27mm);
      \draw (-9mm,30mm) -- (-11mm,33mm);
      \draw (-9mm,30mm) -- (-7mm,33mm);
      % braccio destro: alzato e leggermente diverso
      \draw (0,23mm) -- (8mm,29mm) -- (11mm,26mm);
      \draw (8mm,29mm) -- (10mm,32mm);
      % gamba sinistra: divaricata, piede a sinistra
      \draw (1mm,14mm) -- (-6mm,4mm) -- (-9mm,5mm);
      % gamba destra: divaricata, piede a destra, lunghezza diversa
      \draw (1mm,14mm) -- (7mm,3mm) -- (11mm,6mm);
    \end{scope}
  \end{tikzpicture}

  \begin{center}
    \vspace*{58mm}
    {\sffamily\Huge\bfseries\color{ocra} Il silenzio\\[2pt] della roccia\par}
    \vspace{10mm}
    {\sffamily\normalsize Perch\'e le montagne venete sembrano prive\\
    di arte rupestre}
    \vspace{18mm}
    {\sffamily\large Mauro Barbieri\par}
  \end{center}

  \vfill
  \begin{center}
    {\sffamily\small\color{white} mauro.barbieri@pm.me \par}
  \end{center}
  \vspace*{8mm}
\end{titlepage}

\chapter*{Prefazione}
\addcontentsline{toc}{chapter}{Prefazione}

Questo libro nasce da una domanda semplice, posta guardando due fotografie affiancate: perché le rocce delle Ande sembrano parlare, coperte di segni lasciati da chi le ha abitate migliaia di anni fa, mentre le rocce delle Alpi venete sembrano tacere?

La domanda è ingannevole, e lo si scopre subito: non è chiaro se il silenzio delle Alpi venete sia reale o se sia il silenzio di chi non ha ancora ascoltato bene. Rispondere ha richiesto di allontanarsi più volte dalla domanda di partenza, per poi tornarci sopra con qualcosa in più.

Il metodo seguito in queste pagine è stato semplice da enunciare e non sempre facile da rispettare: ogni volta che una fonte offriva una spiegazione già pronta, si è cercato di risalire al lavoro originale su cui quella spiegazione si basava, invece di accontentarsi della sintesi. Non è un esercizio di pedanteria. Una sintesi, per sua natura, toglie i numeri, le eccezioni, i casi che non tornano: proprio le cose che permettono a chi legge di giudicare da sé quanto un'affermazione sia solida. Un manuale che dice ``i cacciatori preistorici sceglievano le selle per la loro posizione strategica'' dice qualcosa di vero, ma non permette di capire se quella scelta fosse una regola quasi assoluta o una tendenza fra tante, quanti siti la confermano, quanti la smentiscono, e con quali soglie di distanza, quota, visibilità. Solo tornando alla ricerca originale quei numeri riemergono, e con essi la possibilità di usarli per un problema nuovo, come quello posto in queste pagine per il Veneto.

Questo modo di procedere ha anche un'altra conseguenza, meno comoda ma più onesta: molte delle strade imboccate in questo lavoro si sono rivelate sbagliate, o comunque troppo semplici per il problema che si voleva risolvere. Si è pensato, in un primo momento, che il movimento degli animali potesse essere ricostruito con un semplice calcolo di costo energetico legato alla pendenza del terreno. Si è pensato che la vegetazione dell'epoca potesse essere un vincolo utile per restringere la ricerca. Si è pensato che i valichi di montagna fossero, di per sé, il segno più affidabile di un luogo frequentato in epoca preistorica. Ognuna di queste idee, messa alla prova dei dati reali e della letteratura scientifica, ha mostrato crepe che ne consigliavano l'abbandono o la correzione. Questo libro non nasconde quel percorso: lo racconta, perché il valore di un metodo si vede tanto nei vicoli ciechi che sa riconoscere quanto nelle strade che alla fine imbocca.

Il risultato non è un manuale definitivo, né un algoritmo pronto per trovare incisioni rupestri sulle Alpi venete. È piuttosto un tentativo di mostrare come si costruisce un'ipotesi di ricerca ragionevole quando i dati sono scarsi, le fonti divulgative sono generiche, e l'unico modo per procedere è tornare, un passo alla volta, alle domande più semplici: chi viveva qui, quando, come si muoveva, e perché.

Una nota pratica, per chi legge: le date in questo libro non sono espresse con le sigle convenzionali della letteratura archeologica (BP, cal BP e simili), la cui definizione tecnica \`e poco intuitiva per chi non lavora nel settore. Si \`e scelto invece di misurare il tempo rispetto all'anno zero del calendario comune, quello che tutti hanno incontrato a scuola. La conversione, con le sue approssimazioni dichiarate, \`e spiegata nel capitolo dedicato al tempo e al clima. Allo stesso modo, i termini tecnici della preistoria (nomi di culture, di fasi climatiche, di industrie litiche) sono stati evitati nel corpo del testo ogni volta che un'espressione più semplice poteva bastare; dove il termine tecnico resta necessario, compare fra parentesi alla prima occorrenza ed è raccolto, con la sua definizione, nel glossario finale.


\tableofcontents

\mainmatter

% ============================================================
% PARTE I — La domanda
% ============================================================
\part{La domanda}

% Cap. 1 — PLACEHOLDER: panoramica globale arte rupestre + domanda veneta
% (file: chapters/01_apertura.tex)
% =============================================================
% PLACEHOLDER — Capitolo 1: apertura
% =============================================================
% Da scrivere. Contenuto previsto:
%   - Panoramica globale dell'arte rupestre (grotte e superfici aperte)
%   - Esempi principali: Chauvet, Lascaux, Altamira (cave art paleolitica)
%     Valcamonica, Monte Bego, Atacama, Kimberley, Bhimbetka (aperto)
%   - Tipi di società associati alla produzione di arte rupestre
%   - In Veneto ce n'è pochissima documentata: la domanda è posta
%   - Breve, diretto, senza enfasi su nessuna regione specifica
% Titolo definitivo: da definire
% =============================================================

\chapter{[Titolo da definire]}
\label{cap:apertura}

\emph{Capitolo in preparazione.}


% Cap. 2 — Il territorio e le sue storie (invariato)
% Catena causale: clima → ecologia → occupazione → produzione simbolica
\chapter{Il territorio e le sue storie}
\label{cap:storia_territorio}

Per capire cosa si potrebbe trovare nelle Alpi venete e perché il quadro attuale è quello che è, è necessario avere una base fatta di tre strati sovrapposti che si condizionano a vicenda in modo causale: il clima, che determina l'ecologia; l'ecologia, che determina il tipo di occupazione umana possibile; e il tipo di occupazione, che determina cosa si poteva produrre, inclusa l'arte. Trattare questi livelli come descrizioni parallele indipendenti è comodo, ma nasconde le connessioni che contano. Questo capitolo li intreccia, organizzando il materiale in una sezione geografica e geologica, una climatica e una ecologica, seguite da una storia umana suddivisa in quattro periodi di peso equivalente: Paleolitico, Mesolitico, Neolitico-Eneolitico, e Bronzo-Ferro. Il peso narrativo più elevato ricade sul Tardoglaciale e sul Mesolitico, perché è in quei millenni che le oscillazioni climatiche determinavano direttamente se le Prealpi erano accessibili o no, e il tipo di occupazione umana che ne conseguiva. Il quadro d'insieme è riassunto nelle due tavole che seguono; le sezioni successive ne sviluppano il contenuto.

% ─────────────────────────────────────────────────────────────────
% Tavola 1 di 2: Pleistocene medio → Tardoglaciale
% ─────────────────────────────────────────────────────────────────
\begin{table}[p]
\small
\setlength{\extrarowheight}{2pt}
\caption{Quadro cronologico (parte prima: Pleistocene e Tardoglaciale). Date in anni prima dell'anno zero astronomico (v.\ glossario). \textit{Continua nella tavola seguente.}}
\label{tab:cronologia_a}
\begin{tabularx}{\textwidth}{>{\raggedright\arraybackslash}X
                              >{\centering\arraybackslash}p{3.2cm}
                              >{\raggedright\arraybackslash}X}
\toprule
\textbf{Geologia / Clima} & \textbf{Anni prima anno zero} & \textbf{Presenza umana in Veneto} \\
\midrule
% ── Pleistocene medio ──────────────────────────────────────────
\rowcolor{blue!8}
Pleistocene medio. Alternanza di cicli glaciali e interglaciali. Clima mediamente più freddo dell'attuale. Lessini e Berici come isole di microclima mite nella pianura. &
500.000--130.000 &
Paleolitico inferiore. \textit{Homo heidelbergensis} e forme arcaiche. Strumenti litici in selce su Lessini e Berici. Frequentazione stagionale di grotte e ripari. \\
% ── confine Pleistocene medio / Eemiano ──────────────────────
\specialrule{0.8pt}{2pt}{2pt}
% ── Eemiano ──────────────────────────────────────────────────
\rowcolor{orange!10}
Eemiano (ultimo interglaciale). Temperature di 1--2\textdegree C superiori alle attuali. Foreste temperate sull'arco alpino. &
130.000--115.000 &
Presenza umana documentata sporadicamente. Transizione verso forme neanderthaliane. \\
% ── confine Eemiano / Würm ───────────────────────────────────
\specialrule{0.8pt}{2pt}{2pt}
% ── Würm: Paleolitico medio ──────────────────────────────────
\rowcolor{cyan!8}
Würm. Avanzata progressiva dei ghiacciai alpini. Steppa-tundra in pianura; megafauna glaciale abbondante. &
115.000--38.000 &
Paleolitico medio. \textit{Homo neanderthalensis}. Riparo Tagliente, Grotta di San Bernardino, Ponte di Veia. Caccia a cervo, cavallo, orso delle caverne; mammut nelle fasi più fredde. \\
\cmidrule(lr){3-3}
% ── Würm / UMG: Paleolitico superiore ───────────────────────
\rowcolor{cyan!8}
Würm. Ultimo Massimo Glaciale (UMG, 21.000--19.000): ghiacciai fino alle pianure. Adriatico 120 m sotto l'attuale; pianura veneta estesa verso est. &
38.000--19.000 &
Paleolitico superiore. \textit{Homo sapiens}. Fumane: arte parietale in grotta (41.000 a.p.a.z.). Gruppi Epigravettiani sulla pianura e ai margini prealpini. \\
% ── confine Würm / Tardoglaciale ─────────────────────────────
\specialrule{0.8pt}{2pt}{2pt}
% ── Tardoglaciale: Dryas pre-Bølling ────────────────────────
\rowcolor{yellow!15}
Tardoglaciale iniziale. Primo ritiro dei ghiacciai, irregolare. Prealpi ancora in gran parte inaccessibili. &
19.000--12.700 &
Epigravettiano evoluto. Cacciatori sulla pianura; siti prealpini come avamposti marginali. \\
\cmidrule(lr){3-3}
% ── Bølling-Allerød ──────────────────────────────────────────
\rowcolor{yellow!15}
Bølling-Allerød: riscaldamento rapido (+5--7\textdegree C in poche centinaia di anni). Betulla e pino avanzano in quota. Prealpi accessibili. &
12.700--10.950 &
Prima colonizzazione delle quote. Riparo Dalmeri (area veneto-trentina): lastre dipinte con ocra. Val Lastaro (Asiago). Arte su supporto portatile. \\
\cmidrule(lr){3-3}
% ── Dryas Recente ────────────────────────────────────────────
\rowcolor{yellow!15}
Dryas Recente: raffreddamento brusco (10.950--9.700). I ghiacciai riavanzano parzialmente. Le quote medie diventano di nuovo ostili. &
10.950--9.700 &
I siti alpini vengono abbandonati. Il registro archeologico si svuota alle quote medie e alte. \\
\bottomrule
\end{tabularx}
\end{table}

% ─────────────────────────────────────────────────────────────────
% Tavola 2 di 2: Olocene
% ─────────────────────────────────────────────────────────────────
\begin{table}[p]
\small
\setlength{\extrarowheight}{2pt}
\caption{Quadro cronologico (parte seconda: Olocene). Date in anni prima dell'anno zero astronomico (v.\ glossario). \textit{Continua dalla tavola precedente.}}
\label{tab:cronologia_b}
\begin{tabularx}{\textwidth}{>{\raggedright\arraybackslash}X
                              >{\centering\arraybackslash}p{3.2cm}
                              >{\raggedright\arraybackslash}X}
\toprule
\textbf{Geologia / Clima} & \textbf{Anni prima anno zero} & \textbf{Presenza umana in Veneto} \\
\midrule
% ── Olocene antico: Mesolitico ───────────────────────────────
\rowcolor{green!8}
Olocene antico. Riscaldamento rapido e definitivo. Foreste avanzano verso l'alto; megafauna glaciale scompare. Lupo, lince e orso abbondanti. &
9.700--7.500 &
Mesolitico sauveterriano. Riorganizzazione forzata dall'ecologia. Riparo Battaglia (Altopiano dei Sette Comuni), Val Lastaro. Siti fino a 2.000 m. \\
\cmidrule(lr){3-3}
% ── Olocene medio: Castelnoviano / Neolitico ─────────────────
\rowcolor{green!8}
Olocene medio. Optimum climatico (7.000--6.000 a.p.a.z.): +1--2\textdegree C. Adriatico raggiunge la configurazione attuale; laguna veneziana in formazione (6.000--5.000 a.p.a.z.). &
7.500--5.250 &
Mesolitico castelnoviano; poi Neolitico (da 7.500 a.p.a.z.). Mondeval de Sora: sepoltura con ocra e conchiglie marine. Prime comunità agricole su Lessini, Berici, Euganei. \\
\cmidrule(lr){3-3}
% ── Eneolitico ───────────────────────────────────────────────
\rowcolor{green!8}
Olocene medio-recente. Clima ancora mite. Disboscamento antropico inizia a comparire nelle sequenze polliniche. &
5.250--4.250 &
Eneolitico: primi oggetti in rame. Sepolture collettive nelle grotte prealpine. Palafitte lacustri (Fimon, Lago della Costa di Arquà Petrarca). Perifericità rispetto alla Valcamonica. \\
\cmidrule(lr){3-3}
% ── Età del Bronzo ───────────────────────────────────────────
\rowcolor{green!8}
Olocene recente. Lieve raffreddamento. Pressione antropica intensa sulle foreste; espansione di prati e pascoli d'altura. &
4.250--2.750 &
Età del Bronzo. Palafitte del Garda (UNESCO). Bostel di Rotzo (Altopiano dei Sette Comuni). Prime incisioni rupestri databili in Veneto (Val d'Assa, Garda). Piroghe fluviali. \\
\cmidrule(lr){3-3}
% ── Età del Ferro ────────────────────────────────────────────
\rowcolor{green!8}
Olocene recente. Clima vicino all'attuale. &
2.750--0 &
Età del Ferro. Veneti antichi: alfabeto venetico, situle figurate, centri urbani (Padova, Este, Altino). Conquista romana (I sec.\ a.p.a.z.). \\
\bottomrule
\end{tabularx}
\end{table}

% ─────────────────────────────────────────────────────────────────
\section{Il territorio: geomorfologia e geologia di base}

Il Veneto occupa una posizione geografica di transizione tra le catene alpine e la Pianura Padana. Questa posizione si riflette in una varietà di ambienti che raramente si trova in una regione di dimensioni comparabili: le Dolomiti e le Prealpi calcaree a nord, i Colli Euganei e i Berici nel cuore della pianura, il delta del Po a est.

Dal punto di vista geologico, la distinzione più rilevante per la storia dell'arte rupestre è quella litologica. La grande maggioranza delle Prealpi venete è calcarea: calcare massiccio cretacico nelle Dolomiti, calcari marnosi e selciferi nel settore prealpino, con intercalazioni di dolomia. Queste rocce sono erodibili, solubili nelle precipitazioni acide, e producono superfici che si levigano nel tempo anziché mantenersi aspre. L'incisione è possibile, come dimostrano i siti noti, ma la durata della leggibilità è molto inferiore rispetto alle superfici vulcaniche o metamorfiche.

Fanno eccezione i Colli Euganei, un massiccio di rocce magmatiche intrusive ed effusive (trachiti, fonoliti, rioliti) di età eocenica, e alcune aree della Lessinia con affioramenti di basalto e tufo di origine vulcanica. Queste rocce dure e resistenti costituiscono superfici più favorevoli all'incisione duratura. I Colli Euganei emergono isolati dalla pianura fino a 601 metri (Monte Venda), e i Colli Berici, di natura calcarea, raggiungono i 444 metri (Monte Alonge). Entrambi i gruppi montuosi rappresentano isole geomorfologiche nel piatto paesaggio padano: nei millenni in cui la pianura era steppa glaciale o palude, queste colline erano zone di microclima più mite, con esposizioni meridionali calde, grotte nei Berici, selce di buona qualità nella Lessinia, e acqua affidabile. Hanno funzionato, per tutta la preistoria, come rifugi ecologici e concentratori di risorse, attraendo sia la fauna sia i gruppi umani che la cacciavano.

La pianura veneta, nella sua configurazione attuale, è il prodotto di millenni di sedimentazione alluvionale. Sotto gli strati storici giacciono dossi fluviali, canali abbandonati e depositi di greto che testimoniano un paesaggio molto più dinamico e acquitrinoso di quello odierno. In questo paesaggio di pianura soggetta a inondazioni stagionali, i pochissimi punti rialzati erano di estrema importanza: piccole alture naturali, creste di antico meandro, bordi di terrazza fluviale. Quegli stessi punti hanno attirato gli insediamenti preistorici per millenni, e il reticolo dei toponimi veneti conserva ancora, nella parola \textit{motta} (piccolo rilievo nel territorio piatto), la memoria di quella geografia.

La presenza di metalli e minerali utili non è irrilevante. La Lessinia e le Prealpi vicentine hanno restituito selce di ottima qualità, estratta e lavorata almeno dal Paleolitico medio. I Colli Euganei contengono rame e altri metalli che diventano rilevanti a partire dall'Eneolitico. La disponibilità di ocre e ossidi di ferro per i pigmenti rossi è attestata in diverse località prealpine.

% ─────────────────────────────────────────────────────────────────
\section{Storia climatica: dal Pleistocene all'età del Ferro}

L'arco cronologico trattato in questo capitolo copre mezzo milione di anni. In quel lasso di tempo il clima non era stabile: il Pleistocene medio e superiore sono caratterizzati da una serie di cicli glaciale-interglaciale di durata variabile tra 40.000 e 100.000 anni, scanditi da oscillazioni della temperatura media globale di 5--10\textdegree C. Nell'Italia nordorientale questo si traduceva in alternanze tra fasi di avanzata glaciale (con ghiacciai che scendevano fino alle quote pedemontane, steppa-tundra in pianura, fauna di ambiente aperto) e fasi interglaciali (con foreste che salivano in quota, temperature simili o superiori alle attuali, fauna temperata). I siti del Paleolitico inferiore e medio della Lessinia e dei Berici si distribuiscono in questo contesto oscillante: certi livelli stratigrafici, ricchi di resti di orso delle caverne e mammut, corrispondono alle fasi più fredde; altri, con cervo e rinoceronte, alle fasi meno fredde.

L'ultimo interglaciale, detto Eemiano, è collocato tra circa 130.000 e 115.000 anni prima dell'anno zero. Le temperature estive erano di 1--2\textdegree C superiori alle attuali; le foreste temperano risalivano in quota; il livello del mare era leggermente più alto di oggi. Al termine dell'Eemiano inizia la glaciazione del Würm, lunga e complessa, con fasi di intensificazione e attenuazione. Il culmine si raggiunge durante l'Ultimo Massimo Glaciale (UMG), collocato tra circa 21.000 e 19.000 anni prima dell'anno zero: in quel momento i ghiacciai alpini scendevano fino agli sbocchi delle valli in pianura; le temperature estive erano inferiori di 8--10\textdegree C rispetto a oggi; la linea degli equilibri glaciali si trovava circa 1.100 metri più in basso dell'attuale. Il livello del mare Adriatico era inferiore di circa 120 metri: non esisteva né Venezia né la sua laguna, e il Po sfociava molto più a est, verso la costa croata attuale.

Il Tardoglaciale (19.000--9.700 anni prima dell'anno zero) è il periodo di maggiore rilevanza per la preistoria veneta, e la sua struttura interna merita attenzione. Non è un semplice scioglimento graduale: è una sequenza di oscillazioni rapide. Una prima fase di riscaldamento porta al Bølling-Allerød (12.700--10.950 anni prima dell'anno zero), un periodo caldo in cui le temperature salirono di 5--7\textdegree C in poche centinaia di anni, i ghiacciai si ritirarono significativamente, e le Prealpi diventarono per la prima volta accessibili a gruppi umani. Poi, circa 10.950 anni prima dell'anno zero, il Dryas Recente: un raffreddamento brusco di analoga intensità in pochi decenni, con ghiacciai che riavanzano parzialmente e quote medie che diventano di nuovo inospitali. Il Dryas Recente dura circa 1.300 anni e si conclude circa 9.700 anni prima dell'anno zero con il passaggio all'Olocene. Il capitolo~\ref{cap:tempo_ghiacci_clima} tratta queste oscillazioni nel dettaglio; qui basta sottolineare che la finestra del Bølling-Allerød e la chiusura del Dryas Recente hanno avuto conseguenze dirette e documentabili sulla distribuzione dei siti umani nel Veneto.

Con l'inizio dell'Olocene (9.700 anni prima dell'anno zero) il riscaldamento diventa definitivo. Le temperature raggiunsero e superarono i valori attuali durante l'optimum climatico, collocato tra circa 7.000 e 6.000 anni prima dell'anno zero. Il limite superiore delle foreste si trovava alcune centinaia di metri più in alto di oggi. Il livello dell'Adriatico settentrional risalì progressivamente, raggiungendo la configurazione attuale tra 6.000 e 5.000 anni prima dell'anno zero; in quel processo si formò la laguna veneziana. I laghi prealpini, incluso il Garda, raggiunsero i livelli attuali tra la fine del Tardoglaciale e l'inizio dell'Olocene. Le torbiere dell'altopiano del Cansiglio e del Montello conservano sequenze polliniche continue che permettono di ricostruire con grande dettaglio la storia della vegetazione locale.

Nell'Olocene recente (da circa 5.000 anni prima dell'anno zero) le temperature calarono leggermente; l'età del Bronzo fu mediamente più calda dell'attuale, con il treeline più elevato; un raffreddamento progressivo tra 3.500 e 2.500 anni prima dell'anno zero contribuì all'abbandono di alcuni insediamenti di alta quota. In questo periodo si sovrappone all'effetto climatico quello della pressione antropica: le foreste erano già in ritirata per il disboscamento.

% ─────────────────────────────────────────────────────────────────
\section{Storia ecologica: paesaggi e fauna in trasformazione}

Al culmine dell'Ultimo Massimo Glaciale, il paesaggio veneto era radicalmente diverso dall'attuale. Le Prealpi e le Dolomiti erano in gran parte coperte di ghiaccio o di nuda roccia paraglaciale. La pianura sottostante era una steppa-tundra attraversata da fiumi a canali intrecciati, con scarsa copertura vegetale arborea e forte erosione eolica. La megafauna glaciale era ancora abbondante: mammut lanoso (\textit{Mammuthus primigenius}), rinoceronte lanoso (\textit{Coelodonta antiquitatis}), uro (\textit{Bos primigenius}), bisonte delle steppe (\textit{Bison priscus}), cavallo selvatico, cervo gigante (\textit{Megaloceros giganteus}) e renna. I grandi erbivori di ambiente aperto si muovevano su vasti areali stagionali nella pianura e nei fondovalle.

In questo contesto, i Colli Euganei e i Colli Berici erano anomalie significative. Emergendo di 400--600 metri dalla pianura glaciale, con esposizioni meridionali calde, grotte nei Berici, selce disponibile nella Lessinia, e fonti d'acqua affidabili, questi rilievi concentravano risorse biologiche scarse nel paesaggio circostante: vegetazione più diversificata, rifugio dal vento, microclima mite. Funzionavano come rifugi ecologici nel senso preciso del termine: aree che mantengono popolazioni di specie durante i periodi avversi per il resto del territorio. Quella stessa funzione di concentratore biologico li rendeva attraenti per i gruppi umani che cacciavano in quel paesaggio.

Con il Bølling-Allerød, il paesaggio cambiò rapidamente. La betulla e il pino avanzarono in quota, le foreste di pino si stabilirono sui versanti prealpini, e le Prealpi divennero accessibili per la prima volta. La megafauna glaciale cominciò a cedere: il mammut scomparve dall'arco alpino prima della fine del Tardoglaciale; il cervo gigante lasciò tracce sporadiche fino a circa 9.700 anni prima dell'anno zero. Il Dryas Recente provocò un parziale rollback di questi processi, con le foreste che scendevano di nuovo e la fauna glaciale che resisteva ancora in alcune aree. All'inizio dell'Olocene, con la definitiva stabilizzazione del clima, la composizione faunistica si trasformò completamente: i grandi erbivori di steppa scomparvero dal registro faunistico locale, e gli ungulati di media taglia (stambecco, camoscio, cervo, capriolo, poi alce nelle zone boschive) divennero le specie dominanti. Le foreste di quercia mista, poi di faggio, avanzarono fino alle quote alpine, spostando il treeline di 200--300 metri verso l'alto rispetto alla posizione attuale.

Un aspetto spesso trascurato è il regime di predazione. Nell'Olocene antico e medio, il lupo era abbondante sull'intero arco alpino, la lince era presente, l'orso bruno frequentava i versanti boscosi. Una popolazione di camosci o stambecchi sottoposta a predazione intensa da parte di carnivori naturali e di cacciatori umani specializzati sviluppa pattern di movimento, dimensioni del territorio e comportamenti di allarme radicalmente diversi da quelli osservabili nelle popolazioni moderne, private dei grandi carnivori per secoli. Questo ha conseguenze metodologiche per l'uso dei dati ecologici attuali nella ricostruzione del comportamento delle prede preistoriche, discusse nel capitolo~\ref{cap:cosa_cacciavano}.

Con il Neolitico e poi con il Bronzo, la deforestazione antropica aggiunse al cambiamento climatico un secondo vettore di trasformazione del paesaggio. Le sequenze polliniche mostrano aumenti marcati dei pollini di cereali e di piante ruderali a partire da circa 5.000 anni prima dell'anno zero. Il treeline scese gradualmente per effetto combinato del leggero raffreddamento e dell'espansione del pascolo d'alta quota. I prati e le brughiere che oggi caratterizzano l'Altopiano dei Sette Comuni e le zone sommitali delle Prealpi sono in parte il prodotto di questa lunga pressione antropica, non un elemento naturale del paesaggio alpino.

% ─────────────────────────────────────────────────────────────────
\section{Storia umana}

\subsection{Paleolitico (500.000--11.000 anni prima dell'anno zero)}

Le tracce più antiche di presenza umana nel Veneto risalgono al Paleolitico inferiore, tra 500.000 e 300.000 anni prima dell'anno zero, e si concentrano sui Monti Lessini e sui Colli Berici. Si tratta di strumenti in pietra scheggiata attribuibili a gruppi di \textit{Homo heidelbergensis} o forme arcaiche simili: piccoli nuclei nomadi, senza arte riconoscibile, che frequentavano grotte e ripari come basi temporanee di caccia. La concentrazione geografica su questi due rilievi non è casuale: la selce di qualità della Lessinia forniva il materiale per gli strumenti, le grotte dei Berici offrivano riparo, e il microclima più mite di entrambi i gruppi collinari era un vantaggio reale in un paesaggio di pianura glaciale. La funzione di rifugio ecologico che Euganei e Berici avranno per tutta la preistoria si afferma già in questo primo contatto documentato tra essere umano e territorio veneto.

Con il Paleolitico medio (da circa 120.000 a 38.000 anni prima dell'anno zero) la presenza umana diventa più articolata. I Neanderthal occupano Lessinia e Berici lasciando stratigrafie potenti e leggibili. I siti principali sono il Riparo Tagliente nella Valle Squaranto, la Grotta di San Bernardino nei Berici e numerosi ripari minori nella Lessinia. Tra i siti della Lessinia, Ponte di Veia merita menzione separata: si tratta di uno dei più grandi archi naturali in roccia calcarea d'Europa, nella Valle dell'Alpone, con evidenti tracce di frequentazione musteriana. Un arco naturale di quelle proporzioni non era solo un rifugio fisico: era un punto di riferimento nel paesaggio percepibile da distanze considerevoli, e probabilmente aveva un peso nella geografia mentale dei gruppi che lo frequentavano. La fauna presente nei livelli stratigrafici cambia con il clima: nei livelli corrispondenti alle fasi più fredde del Würm compaiono mammut, rinoceronte lanoso e orso delle caverne; nei livelli meno freddi predominano cervo, cavallo e camoscio. I siti della Lessinia funzionano così anche come termometri del passato, non solo come archivi di comportamento umano.

L'arrivo di \textit{Homo sapiens} in Europa è datato a circa 43.000 anni prima dell'anno zero. I Neanderthal si estinsero circa 38.000 anni prima dell'anno zero, dopo un periodo di sovrapposizione geografica discusso. La grotta di Fumane, in Lessinia (Marano di Valpolicella), è il sito dove questa transizione lascia la traccia più straordinaria del Veneto: frammenti di intonaco dipinto con figure di animali e figure ibride umano-animali, databili a circa 41.000 anni prima dell'anno zero, le testimonianze di arte parietale più antiche d'Italia. Non si tratta di arte rupestre all'aperto nel senso discusso in questo libro, ma di arte in ambiente di grotta, associata a una fase di utilizzo intensivo della caverna come base stagionale di caccia. La sua importanza sta nell'attestare che la capacità di produrre immagini simboliche era già pienamente sviluppata in quest'area alla comparsa di \textit{Homo sapiens}.

Le fasi centrali del Paleolitico superiore (Aurignaziano, Gravettiano) sono scarsamente documentate nel Veneto rispetto ad altre regioni italiane, e questa lacuna riflette probabilmente la distribuzione reale dell'occupazione: durante l'UMG (21.000--19.000 anni prima dell'anno zero), con i ghiacciai che bloccavano le valli alpine e la pianura esposta e ventosa, i cacciatori Epigravettiani vivevano principalmente sulla grande pianura padano-adriatica ora in parte sepolta sotto l'alluvio e in parte sommersa dall'Adriatico. I siti prealpini erano la periferia di quel sistema, non il suo centro. Questo spiega la scarsità di evidenze venete in quelle fasi: non c'erano meno persone, erano altrove.

Con il Bølling-Allerød (12.700--10.950 anni prima dell'anno zero) la situazione cambia in modo documentabile. Per la prima volta i ghiacciai si ritirano abbastanza da rendere le Prealpi accessibili anche alle quote medie e alte durante la stagione calda. I cacciatori seguono le prede verso l'alto. Il sito più emblematico di questa prima colonizzazione delle quote è il Riparo Dalmeri, nell'area veneto-trentina sull'Altopiano di Marcesina, datato all'Allerød e collocato a circa 1.100 metri di quota: vi sono state trovate oltre trenta lastre di calcare decorate con figure animali in ocra rossa (cervi, stambecchi, orsi, uccelli), depositate con cura nel contesto di un accampamento ricorrente. Nella stessa finestra temporale si collocano i ritrovamenti di Val Lastaro sull'Altopiano di Asiago e la sepoltura di Villabruna nel Bellunese, datata a circa 12.000 anni prima dell'anno zero. Queste non sono soste occasionali: sono campi base stagionali a cui i gruppi tornavano anno dopo anno. L'arte del Dalmeri è su lastre portatili, non su superfici fisse: un gruppo ancora semi-mobile investiva in immagini simboliche, ma su oggetti che potevano viaggiare con lui. Il passaggio all'arte su superficie fissa all'aperto richiedeva ancora qualcosa in più.

Il Dryas Recente (10.950--9.700 anni prima dell'anno zero) interrompe bruscamente questa colonizzazione. I siti di quota vengono abbandonati; il registro archeologico alle altitudini medie e alte si svuota. Non è solo un fatto climatico: è la prova che la colonizzazione delle Prealpi durante il Bølling-Allerød era contingente alle condizioni ambientali, non ancora definitiva. Il Mesolitico non sarà una continuazione lineare di quell'esperienza, ma una seconda partenza in un paesaggio trasformato.

\subsection{Mesolitico (11.700--7.500 anni prima dell'anno zero)}

La fine del Dryas Recente (9.700 anni prima dell'anno zero) segna l'inizio di una riorganizzazione ecologica che cambia profondamente le condizioni di vita dei cacciatori-raccoglitori veneti. Quella riorganizzazione non è una scelta: è una risposta forzata. Le foreste avanzano rapidamente verso le quote medie e alte, chiudendo la pianura aperta che era stata il teatro principale della vita paleolitica. Con la foresta scompaiono i grandi erbivori di ambiente aperto, cavallo e renna, che avevano costituito le prede principali per millenni. Chi rimane nel territorio veneto trova le sue prede principali nelle nuove condizioni alpine: stambecco e camoscio in quota, cervo e capriolo ai margini della foresta. Per raggiungerli sistematicamente, i cacciatori devono colonizzare le quote, e questa volta in modo definitivo.

Questa seconda colonizzazione delle Prealpi, che apre il Mesolitico, è documentata da una moltiplicazione di siti a quote progressivamente più elevate. Sull'Altopiano dei Sette Comuni (Asiago, Vicenza), il Riparo Battaglia testimonia l'occupazione sauveterriana del plateau già nelle prime fasi dell'Olocene: lo strumentario microlitico, con le sue punte a dorso e i suoi segmenti geometrici, è pensato per armi composite come la freccia, adatte alla caccia in ambienti boschivi e rupestri. Val Lastaro, sempre sull'altopiano asiaghese, ha restituito materiali sia epigravettiani sia mesolitici, confermando che queste quote erano frequentate con continuità stagionale. I cacciatori mesolitici non scoprivano per la prima volta questi luoghi: in parte li conoscevano dalle generazioni dell'Allerød. Li ritrovavano in un paesaggio trasformato e li riconfiguravano per una diversa economia.

I siti si moltiplicano alle quote di 1.500--2.000 metri e oltre. Le Prealpi vicentine e bellunesi, la Lessinia e le propaggini dell'Altopiano di Asiago sono costellate di stazioni mesolitiche di superficie, spesso individuate dai ritrovamenti di microliti geometrici caratteristici. In molti casi si tratta di posizioni con ampia visibilità sulle valli sottostanti, scelte per il controllo del territorio di caccia più che per il riparo fisico. Il modello insediativo che emerge è quello di una mobilità verticale stagionale: quote basse in inverno, quote alte in estate, con spostamenti lungo corridoi vallivi ben definiti. L'Altopiano dei Sette Comuni, con la sua posizione di cerniera tra la pianura vicentina e le valli dolomitiche, era uno di questi nodi nella rete di movimento stagionale.

Mondeval de Sora (San Vito di Cadore, Belluno), a 2.150 metri di quota nelle Dolomiti bellunesi, è il sito mesolitico veneto più noto al grande pubblico. Una sepoltura individuale con corredo ricco: utensili di osso e selce, pezzi di ornamento, ocra rossa, e conchiglie di specie marine atlantiche. Queste ultime implicano contatti a lunga distanza con le coste dell'Adriatico o del Tirreno, o scambi con gruppi in contatto con le coste. I cacciatori mesolitici dell'alta quota veneta non erano gruppi isolati: erano inseriti in reti di relazione che coprivano centinaia di chilometri. La datazione di Mondeval si colloca intorno a 8.000 anni prima dell'anno zero, nella fase castelnoviana del Mesolitico.

L'arte mesolitica all'aperto rimane la grande questione aperta del Veneto preistorico. Non esistono pannelli incisi o dipinti attribuiti con certezza al Mesolitico nella regione. Le ragioni possibili sono molteplici e non si escludono: degrado differenziale delle superfici rocciose nel corso di 9.000 anni di clima alpino; ricerca ancora insufficiente, soprattutto sulle quote medie e alte; e la possibilità, da non scartare, che i cacciatori mesolitici non praticassero l'arte rupestre su superfici fisse su larga scala, preferendo forme di espressione simbolica su supporti deperibili o portatili, come suggerisce l'esempio del Dalmeri. La questione è discussa nel capitolo~\ref{cap:tafonomia}.

\subsection{Neolitico ed Eneolitico (7.500--4.250 anni prima dell'anno zero)}

Il Neolitico veneto, datato in Italia settentrionale tra circa 7.500 e 5.250 anni prima dell'anno zero, non arriva come una sostituzione netta di popolazione. Le evidenze genetiche e culturali disponibili suggeriscono una trasformazione progressiva dei gruppi mesolitici locali, influenzati da comunità agricole che avanzavano da ovest e da est, piuttosto che una colonizzazione che cancellasse i predecessori. Il contesto climatico è favorevole: l'optimum dell'Olocene medio (7.000--6.000 anni prima dell'anno zero), con temperature leggermente superiori alle attuali e precipitazioni abbondanti, crea condizioni ottimali per l'insediamento stabile presso i bacini lacustri che in quel momento erano pieni e produttivi.

È in questa fase che il ruolo di Euganei e Berici nella storia dell'occupazione umana veneta si esprime con maggiore chiarezza. Il Lago di Fimon, nei Colli Berici a pochi chilometri da Vicenza, è uno dei siti palafitticoli neolitici meglio documentati del Veneto: un villaggio su palafitta ai margini di un bacino lacustre oggi in gran parte interrato, con ceramiche, resti botanici e faunistici che testimoniano un'economia mista di agricoltura, allevamento e caccia. Negli Euganei, il Lago della Costa presso Arquà Petrarca (Padova) ha restituito evidenze di frequentazione preistorica nelle fasi neolitica ed eneolitica: un piccolo bacino lacustre ai piedi delle colline trachitiche, sfruttato probabilmente sia per la pesca sia come punto di approdo e scambio. La continuità di questi luoghi come attrattori umani, dalle fasi paleolitiche di rifugio agli insediamenti neolitici stabili, non è casuale.

La pianura veneta neolitica non era però solo un sistema di laghi e palafitte. Era soprattutto un paesaggio di fiumi e paludi in cui i pochi punti stabili erano gli unici luoghi sicuri per insediarsi. I dossi fluviali, creste sabbiose o ghiaiose lasciate dai meandri abbandonati, erano le piattaforme naturali degli insediamenti di pianura: leggermente sopraelevati rispetto alla pianura alluvionale, con accesso diretto all'acqua, stabili in caso di piena. Queste piccole alture naturali, chiamate ancora oggi \textit{motte} nel lessico geografico veneto, erano punti di riferimento visibili nel paesaggio piatto, segnali nel territorio: non a caso molte di esse hanno restituito materiali preistorici di età diverse, a conferma di una frequentazione che attraversa i millenni. Alcune motte sono strutture artificiali, cumuli di terra costruiti intenzionalmente in epoche successive, ma poggiano spesso su rialzi naturali già usati in preistoria.

La tradizione di segnare il paesaggio con strutture permanenti si esprime nel Neolitico veneto anche attraverso alcune evidenze di strutture megalitiche, documentate in modo ancora frammentario nella fascia prealpina vicentina e sull'Altopiano dei Sette Comuni: menhir isolati, allineamenti di massi, strutture in pietra la cui attribuzione cronologica è dibattuta. La tradizione megaliitica veneta è meno sviluppata e meno studiata di quella presente in Sardegna, Puglia o in Francia, ma la sua esistenza parziale è significativa: i megaliti sono, al pari dell'arte rupestre all'aperto, una forma di marcatura permanente e visibile del paesaggio. La loro presenza periferica nel Veneto eneolitico e neolitico si inserisce nello stesso interrogativo che questo libro pone per l'arte rupestre: con quale intensità e con quale forma le comunità preistoriche venete comunicavano simbolicamente attraverso il paesaggio fisico?

L'Eneolitico (5.250--4.250 anni prima dell'anno zero) porta i primi manufatti in rame, anche se il Veneto rimane in posizione periferica rispetto ai principali circuiti di metallurgia del tempo. I confronti sono inevitabili con le regioni limitrofe: la Lombardia e il Trentino producono in questo periodo statue-stele e le incisioni della Valcamonica sono già in piena attività. In Veneto la traccia più rilevante dell'Eneolitico è funeraria: sepolture collettive nelle grotte prealpine della Lessinia e nei Berici, e sepolture in grotticella nella fascia prealpina. Il sito di Sovizzo (Vicenza) è l'unico complesso eneolitico veneto di carattere funerario-sacrale con evidenze di un certo rilievo. La perifericità del Veneto rispetto ai grandi circuiti di elaborazione culturale del tardo Eneolitico alpino è un dato di fatto: anche in questo caso, come per la successiva arte rupestre del Bronzo, la posizione geografica del Veneto orientale ai margini delle rotte di transumanza e di scambio che percorrevano il versante alpino occidentale sembra aver avuto un ruolo.

\subsection{Età del Bronzo e del Ferro (4.250--0 anni prima dell'anno zero)}

L'età del Bronzo (4.250--2.750 anni prima dell'anno zero) è il periodo meglio documentato della preistoria veneta, e il primo in cui le incisioni rupestri divengono un elemento attestato e databile con sicurezza. Il cambiamento fondamentale rispetto ai secoli precedenti è la sedentarizzazione: le comunità del Bronzo medio e recente costruiscono insediamenti stabili, usano lo stesso luogo per generazioni, e hanno interesse a segnare il territorio in modo permanente. È questa stabilità insediativa, più che un cambiamento nelle capacità tecniche o simboliche, a spiegare la comparsa dell'arte rupestre documentata nel Veneto.

I siti palafitticoli del Lago di Garda sono l'espressione più conosciuta di questa sedentarizzazione. Inclusi dal 2011 nel patrimonio mondiale dell'UNESCO come parte del sito seriale delle palafitte preistoriche circumalpine, i villaggi gardesani mostrano insediamenti stabili e tecnicamente elaborati, con bronzi lavorati, ceramiche di alta qualità e ornamenti. Sulle rocce del Garda e della Valle dell'Asse compaiono coppelle, figure geometriche e zoomorfe che trovano confronti nell'arte rupestre alpina di questa fase su entrambi i versanti della catena. La Val d'Assa, sull'Altopiano di Asiago, è la zona rupestre più rilevante del Veneto per questo periodo: un sistema di incisioni che include figure antropomorfe, scene di caccia e motivi geometrici, in un corridoio naturale che collegava l'altopiano alla pianura vicentina.

Sull'Altopiano dei Sette Comuni, il Bostel di Rotzo (Rotzo, Vicenza) è uno dei castellieri meglio conservati del Veneto: un insediamento d'altura del Bronzo finale e dell'età del Ferro, con strutture difensive in pietra, che domina visivamente una vasta porzione del territorio prealpino. I castellieri veneti, di cui il Bostel è uno degli esempi più studiati, non erano solo insediamenti: erano punti di controllo visivo del territorio, segnali di presenza permanente nel paesaggio. La loro posizione sulle creste e sugli spalloni dei pianori prealpini rispecchia una logica di organizzazione territoriale completamente diversa da quella dei cacciatori mesolitici: non mobilità verticale stagionale, ma presidio fisso di punti strategici.

La pianura veneta del Bronzo era percorsa da un sistema di vie d'acqua che integrava le rotte di terra. I fiumi erano strade: il Brenta, il Bacchiglione, l'Adige e i loro affluenti connettevano le Prealpi alla pianura e al delta del Po. Le piroghe monossili rinvenute nell'area del Bacchiglione e in altri contesti fluviali veneti attestano una navigazione fluviale organizzata, con imbarcazioni scavate da tronchi interi, alcune delle quali di dimensioni considerevoli. La mobilità del Bronzo veneto non era dunque solo terrestre: il territorio era letto e percorso anche lungo l'asse delle acque, e i siti di pianura si distribuiscono in modo coerente con questa logica, preferendo i dossi fluviali stabili vicini ai corsi d'acqua navigabili. Anche le palafitte lacustri dei bacini dei Berici ed Euganei vanno lette in questo quadro: non sono un fenomeno isolato, ma parte di una rete di insediamenti acquatici che sfruttava ogni specchio d'acqua disponibile, dai grandi laghi prealpini ai piccoli bacini collinari.

L'età del Ferro veneta (da circa 2.750 anni prima dell'anno zero) è dominata dalla cultura dei Veneti antichi, una delle grandi civiltà dell'Italia preromana. I Veneti svilupparono un sistema di scrittura (l'alfabeto venetico, derivato dall'etrusco), centri urbani complessi (Padova, Este, Altino, Vicenza), e una tradizione artistica ricca che culmina nelle situle figurate: grandi contenitori di bronzo decorati con scene narrative di banchetto, gara e processione, in uno stile narrativo senza paralleli esatti nell'Italia coeva. Le incisioni rupestri della fase venetica, dove attestate, mostrano figure di cavalieri, motivi geometrici e scene che ricorrono anche nell'arte mobiliare. La conquista romana, completata nel I secolo prima dell'anno zero, portò a una progressiva romanizzazione del paesaggio e della cultura materiale; le tradizioni religiose locali sopravvissero a lungo, come attestano i santuari di tipo \textit{fonte sacra} diffusi nell'area veneta, dove le offerte votive continuarono per secoli.

% ─────────────────────────────────────────────────────────────────
\section{Perché questo quadro conta per la domanda di questo libro}

Tre osservazioni emergono da questa rassegna e sono rilevanti per i capitoli successivi.

La prima: il Veneto preistorico non era una periferia vuota. Dalla frequentazione di \textit{Homo heidelbergensis} sui Lessini ai Veneti antichi, questo territorio è stato abitato con continuità per mezzo milione di anni. La scarsità di arte rupestre nota non rispecchia una scarsità di presenza umana: rispecchia, almeno in parte, un tipo di occupazione (mobile, stagionale, non ancora sedentaria) che non produceva marcatura permanente del paesaggio su larga scala.

La seconda: c'è una correlazione tra tipo di occupazione e tipo di produzione simbolica. I cacciatori paleolitici producono arte in grotta (Fumane) o su supporto portatile (Dalmeri). I cacciatori mesolitici probabilmente non producono arte rupestre su superficie fissa, o la producono in misura non ancora rilevata. Le prime incisioni rupestri all'aperto datate con certezza nel Veneto appartengono all'età del Bronzo, quando le comunità sono sedentarie e hanno interesse a segnare il territorio in modo permanente. Questa gradazione non è casuale: è predetta dalla logica del tipo di insediamento.

La terza: la discontinuità culturale è reale e profonda. Tra i cacciatori mesolitici del Riparo Battaglia e le comunità dell'età del Bronzo sulle rive del Garda ci sono quattro o cinque millenni di trasformazioni culturali, economiche e demografiche. Non c'è ragione di aspettarsi continuità nelle tradizioni simboliche attraverso quelle trasformazioni. Cercare arte rupestre mesolitica e arte rupestre della tarda età del Bronzo con gli stessi criteri, negli stessi luoghi, significa ignorare che erano prodotte da popolazioni con economia, mobilità e rapporto con il territorio profondamente diversi.


% ============================================================
% PARTE II — Il quadro veneto
% ============================================================
\part{Il quadro veneto}

% Cap. 3 — PLACEHOLDER/DA ESPANDERE: screening provincia per provincia
% Veneto geografico + zone contigue (Folgaria, Primiero, sponda trentina
% e lombarda del Garda, morene gardesane, Cansiglio, bassa friulana)
% Siti paleo/meso/neolitici in montagna e pianura; arte rupestre nota;
% grotte, ripari, siti di superficie, palafitte, dossi fluviali
\chapter{Cosa si conosce davvero, in Veneto}
\label{cap:cosa_si_conosce}

Prima di chiedersi cosa manca, conviene fare ordine su cosa esiste, con la maggiore precisione possibile. Il territorio veneto, inteso qui nel senso geografico che tiene conto dei confini naturali piuttosto che amministrativi, non è affatto privo di testimonianze della frequentazione preistorica: possiede siti documentati in ogni periodo, dall'ultima glaciazione fino all'età del ferro, distribuiti in modo disomogeneo tra montagna e pianura, tra aree intensamente indagate e aree che per ragioni storiche e logistiche non hanno ancora ricevuto un'attenzione sistematica comparabile. Questo capitolo li ripercorre provincia per provincia, risalendo per quanto possibile ai lavori originali, con le cifre, le date e le riserve che quei lavori stessi dichiarano.

Una precisazione metodologica preliminare: la ripartizione per provincia moderna è uno strumento pratico, non una categoria preistorica. Le popolazioni paleolitiche e mesolitiche non seguivano confini amministrativi, e molte delle aree più interessanti si trovano nelle fasce di confine tra province o in territori che la geografia fisica renderebbe meglio descritti in relazione ai bacini fluviali. Il concetto di \emph{Veneto geografico esteso} usato in questo libro include, per le stesse ragioni, alcune zone oggi trentine o lombarde ma storicamente legate allo stesso sistema di valli e passi: la sponda settentrionale del Garda, il Primiero, la Valsugana, Folgaria.

% ============================================================
\section{Verona: le Basse Colline, il Garda, la Valpolicella}
% ============================================================

La provincia di Verona ospita due dei tre siti di arte rupestre o arte preistorica documentati con maggiore solidità nell'intero Veneto, e uno dei siti paleolitici più antichi d'Europa. Per questa ragione, e perché la storia della loro scoperta è essa stessa parte del problema che questo libro discute, vengono trattati con il maggior dettaglio consentito dalle fonti disponibili.

\subsection{La grotta di Fumane: quarantamila anni di arte figurativa}

La grotta di Fumane, in Valpolicella, non lontana da Verona, è stata scavata sotto la direzione di Alberto Broglio dell'Università di Ferrara a partire dagli anni novanta del Novecento. Dai livelli attribuiti alla fase più antica della presenza della nostra specie in Europa sono stati recuperati sei frammenti di roccia caduti dalla volta per crioclastismo, ritrovati capovolti sul terreno, dipinti con ocra rossa. Due frammenti sono i più noti e i più discussi in letteratura: una figura antropomorfa alta circa diciotto centimetri, vista frontalmente, che indossa quello che sembra un copricapo o una maschera con due protuberanze simili a corna, con in mano un oggetto interpretato come votivo; e un animale a quattro zampe visto di profilo, verosimilmente un mustelide o un felide, la cui identificazione precisa resta incerta \citep{broglio2003, broglio2005}.

La datazione radiometrica sui livelli associati si aggira attorno ai quarantunomila anni prima dell'anno zero: tra le più antiche mai ottenute per opera figurativa in Europa, e un dato che colloca Fumane in un piccolo gruppo di siti di riferimento mondiale per lo studio delle origini dell'arte figurativa umana.

Fumane non è, in senso stretto, arte rupestre all'aperto: è arte parietale in ambiente di grotta, un fenomeno diverso per almeno due ragioni. La prima è la condizione di conservazione: una grotta protegge dagli agenti atmosferici che, all'aperto, cancellano progressivamente un'incisione o una pittura. La seconda è il contesto di ritrovamento: Fumane è nota da uno scavo stratigrafico controllato, con datazione diretta dei livelli associati. Va tenuto fermo cosa Fumane dimostra senza possibilità di dubbio: che la pratica di produrre immagini su superfici rocciose era presente sul territorio veneto già quarantamila anni fa. Va altrettanto tenuto fermo cosa Fumane non dimostra: non è una prova dell'esistenza di incisioni all'aperto nelle fasi successive, separate da decine di millenni e da tradizioni culturali completamente diverse.

\subsection{Il Monte Luppia e le incisioni del Benaco}

Il complesso di arte rupestre all'aperto più importante del Veneto si trova sulle pendici collinari che si affacciano sul lago di Garda, in particolare sul Monte Luppia, tra Punta San Vigilio e Malcesine. Le prime rocce incise furono individuate nel 1964 da Mario Pasotti, insegnante, dopo una visita alle incisioni della Valcamonica che gli fece riconoscere segni simili sulle rocce vicino a casa. Il dettaglio vale la pena sottolinearlo per intero: Pasotti non era un archeologo di professione, e la sua scoperta non nasce da una prospezione pianificata a tavolino, ma dal riconoscimento di un pattern già visto altrove, applicato a un territorio che nessuno aveva ancora guardato con quello sguardo. Il lavoro di catalogazione fu proseguito e ampliato da Fabio Gaggia, che nel 1985 costituì un centro di studi dedicato al territorio benacense e ne rimane il principale studioso \citep{gaggia1982, gaggia2002}.

Il catalogo oggi comprende più di duecentocinquanta rocce incise, per un totale di almeno tremila figure: guerrieri, cavalieri, imbarcazioni, armi, simboli solari, coppelle isolate o in gruppo, e un numero consistente di segni geometrici di incerta interpretazione. La cronologia copre un arco lungo: le fasi più antiche vengono attribuite all'età del bronzo (la Pietra di Castelletto e la Pietra delle Griselle, collocate nel secondo millennio prima dell'anno zero), le fasi intermedie all'età del ferro (con nuclei a Senge di Marciaga, nel comune di Costermano, e a Crero, nel comune di Brenzone), e uno strato più recente arriva fino a incisioni chiaramente medievali o moderne. Il museo del castello scaligero di Torri del Benaco colloca esplicitamente questo complesso subito dopo la Valcamonica e il Monte Bego per importanza nel panorama alpino.

Il dato più importante per il tema di questo libro, però, non è la dimensione del catalogo: è che quel catalogo è iniziato nel 1964 e non prima. Le rocce del Monte Luppia erano lì, esposte, raggiungibili, da millenni. Quello che è cambiato in quell'anno non è il territorio, ma la disponibilità di uno sguardo capace di riconoscere un'incisione intenzionale: un caso da manuale del meccanismo del research bias, documentabile con un nome, una data e una motivazione dichiarata dallo stesso scopritore.

\subsection{Riparo Tagliente e la Lessinia}

La Lessinia, altopiano carsico a nord di Verona, ospita una serie di siti paleolitici documentati con scavi sistematici. Il più importante è Riparo Tagliente, nella Valle delle Falci, scavato a partire dagli anni sessanta e poi con continuità nei decenni successivi. Il sito ha restituito una lunga sequenza occupazionale che va dal Musteriano (Neandertal) fino all'Epigravettiano finale, con una presenza documentata attorno ai diciottomila anni prima dell'anno zero per le fasi più recenti. La fauna cacciata include cavallo, cervo, stambecco, camoscio e rinoceronte lanoso per le fasi più antiche. Non è stata trovata arte parietale \emph{in situ}, ma il sito ha restituito elementi di ornamento personale e strumenti ossei lavorati.

Nella stessa area: Riparo Mezzena (sequenza musteriana), Grotta della Ghiacciaia (resti faunistici del Pleistocene superiore), e la Grotta di San Bernardino (occupazione paleolitica con resti di Homo sapiens datati a circa ventottomila anni prima dell'anno zero, tra i più antichi della regione) \citep{bartolomei1992}.

% ============================================================
\section{Vicenza: l'Altopiano di Asiago e i Colli Berici}
% ============================================================

\subsection{La Val d'Assa e le incisioni dell'altopiano}

Sull'altopiano di Asiago, presso Canove e Roana, la Val d'Assa presenta un nucleo di incisioni di natura diversa dal Monte Luppia per contesto di scoperta, tecnica e grado di incertezza. Le incisioni furono esposte dalla grande piena del torrente Assa del 1966, che asportando terreno e vegetazione ne scoprì alcune e ne liberò parzialmente altre. La prima pubblicazione scientifica, firmata da Piero Leonardi con Giuseppe Rigoni e Aldo Allegranzi, risale al 1982 \citep{leonardi1982}. Il testo si apriva con una cautela significativa: si parlava di ``varie centinaia o forse addirittura migliaia (il rilevamento è appena iniziato)'' di incisioni, la maggior parte delle quali geometriche, eseguite a graffito, cioè per incisione leggera con uno strumento appuntito; una tecnica che produce segni meno profondi e quindi più difficili da datare.

Gli autori, riprendendo un'osservazione di Paolo Graziosi, concludevano che una parte consistente delle incisioni era probabilmente medievale, senza escludere che alcuni segni geometrici potessero risalire alla preistoria. La cautela non ha impedito la formazione di un'interpretazione più marcatamente preistorica del sito, sostenuta da Ausilio Priuli (1983), che ha mantenuto vivo un dibattito durato decenni, fino al convegno di Gallio e Canove del 1996 (atti 2001).

Un lavoro recente con tecniche digitali di rilievo tridimensionale e luce strutturata ha permesso una lettura più oggettiva delle superfici. La conclusione è che un'origine preistorica appare ``molto dubbia, o eventualmente limitata a segni geometrici che sono essenzialmente privi di tempo'', cioè non databili per loro stessa natura \citep{bettineschi2022}. Gran parte delle figure è di epoca medievale o moderna, comprese croci, raffigurazioni di chiese e diverse date leggibili, la più antica delle quali riporta l'anno 1506. Il sito è attualmente interdetto al pubblico per gravi dissesti del versante.

\subsection{I Colli Berici e la pianura vicentina}

I Colli Berici hanno restituito tracce di frequentazione risalenti al Paleolitico, in particolare nella Grotta di Broion e nella Grotta dei Ladroni, scavate da Luigi Capitan e poi da altri ricercatori nel corso del Novecento. La Grotta di Broion ha restituito industria musteriana associata a fauna tipica delle fasi fredde: rinoceronte lanoso, orso delle caverne, iena delle caverne \citep{bartolomei1980}. La frequentazione epigravettiana dei Berici è più frammentaria, con presenze attestate ma non ancora sistematicamente catalogate in anni recenti.

Il lago di Fimon, nella pianura vicentina, è noto per la presenza di un sito palafitticolo neolitico riportato alla luce da ricerche condotte tra gli anni sessanta e ottanta: ceramica della Cultura dei Vasi a Bocca Quadrata, resti vegetali carbonizzati e fauna domestica. È uno dei pochi siti lacustri del Veneto interno con una sequenza neolitica documentata con scavo.

% ============================================================
\section{Padova: gli Euganei, la pianura e il distretto di Este}
% ============================================================

La provincia di Padova presenta un quadro preistorico in cui montagna e pianura svolgono ruoli nettamente diversi, condizionati dalla geomorfologia: i Colli Euganei, gruppo di rilievi vulcanici isolati nella pianura veneta, hanno offerto risorse minerarie e punti di riferimento visivo eccezionali, mentre la pianura è stata abitata nelle sue fasi più recenti (bronzo, ferro) in modo intenso ma con una tafonomia complicata dalle alluvioni ripetute e dall'agricoltura.

\subsection{Paleolitico e Mesolitico: una presenza difficile da documentare}

Non risultano siti paleolitici con stratigrafia conservata nella provincia di Padova: i depositi pleistocenici della pianura sono stati ampiamente rimaneggiati dal reticolo fluviale e dalle bonifiche. La litica sporadica segnalata in letteratura per i versanti degli Euganei (Baone, Cinto Euganeo) non ha ancora prodotto contesti stratigrafici chiusi che permettano attribuzioni sicure.

La situazione mesolitica è analoga: l'assenza di siti con contesto è la regola, non l'eccezione. Un'eccezione parziale è il sito di Bosco dei Fontanassi a Piombino Dese, sul margine nordorientale della provincia, dove microliti geometrici riferibili al Mesolitico sono stati segnalati in raccolta di superficie in terreni agricoli; ma anche in questo caso la mancanza di contesto stratigrafico rende problematica qualsiasi interpretazione funzionale del sito.

\subsection{Neolitico: i villaggi della pianura}

Il Neolitico padovano è documentato principalmente dalla Cultura dei Vasi a Bocca Quadrata (VBQ), nella sua espressione pianura veneta, con villaggi attestati a Vigodarzere, Mejaniga e Campodarsego: insediamenti di pianura con strutture abitative, ceramica decorata, fauna domestica e industria litica levigata. Questi siti appartengono alla seconda metà del quarto millennio prima dell'anno zero e si collocano in un areale che dalla pianura padana si estende verso la costa adriatica.

Abano Terme, con le sue sorgenti termali calde, è considerata un possibile polo di attrazione già in età preistorica, ma le testimonianze dirette di frequentazione neolitica nell'area termale sono frammentarie, sovrapposte a fasi storiche successive e difficili da isolare senza scavi profondi che le infrastrutture moderne hanno reso impraticabili.

\subsection{Eneolitico: le risorse degli Euganei}

I Colli Euganei in età eneolitica (terzo millennio prima dell'anno zero) costituivano una fonte di materie prime di interesse sovraregionale: la trachite euganea era estratta e lavorata per la produzione di macine, pestelli e strumenti pesanti. Sono state segnalate presenze della Cultura del Vaso Campaniforme in contesti di altura sui Colli, ma il quadro resta poco sistematizzato, con materiali in gran parte da superficie o da scavi vecchi non documentati con standard moderni.

In questo contesto si inserisce una questione aperta: sulle superfici di trachite dei Colli Euganei sono state segnalate, in modo non sistematico, figure a coppella che potrebbero essere di età preistorica, per analogia con coppelle documentate su altri substrati vulcanici nell'arco alpino. Al momento in cui si scrive, non risulta una campagna di rilevamento sistematico che abbia affrontato questa questione con metodi comparabili a quelli usati per l'arte rupestre della Valcamonica o del Monte Luppia.

\subsection{Bronzo e Ferro: l'intensità della pianura}

L'età del bronzo padovana è documentata con una densità che in pianura non ha equivalenti nel Veneto montano. I dossi fluviali residui (paleoisole di sedimento emergenti dall'antica pianura alluvionale) hanno restituito materiali dell'età del bronzo a Padova città, Casalserugo, Legnaro e Este. Il Monte Cero a Baone (Colli Euganei) e la Rocca di Monselice sono siti di altura con frequentazione del bronzo medio e recente, con materiali difensivi e di stoccaggio. Il ritrovamento di piroghe monossili nel letto del Bacchiglione e del Brenta testimonia l'uso sistematico dei fiumi come vie di comunicazione già nel secondo millennio.

Il distretto di Este (Ateste) è il sito di riferimento per l'età del ferro del Veneto meridionale: necropoli con centinaia di tombe, situle decorate a sbalzo con scene narrative, e il santuario della dea Reitia con migliaia di offerte votive sono documentati con scavi sistematici dall'Ottocento \citep{chieco_bianchi1988}. Padova preromana, con il suo impianto urbano databile almeno al V secolo prima dell'anno zero, completa il quadro di una provincia con una densità di frequentazione storico-preistorica molto elevata nelle fasi più recenti.

Scoperte recentissime (2025) hanno portato alla luce una necropoli venetica di trecento tombe in via Campagnola a Padova, e un santuario venetico a Ponso: ulteriori conferme di un distretto densamente abitato durante l'età del ferro, con una produzione simbolica ricca di cui le situle sono il prodotto più noto, ma che non include arte rupestre documentata.

\subsection{Arte rupestre: assente come corpus}

Non risulta arte rupestre documentata nella provincia di Padova. Le coppelle sugli Euganei, già citate, rappresentano un punto interrogativo aperto piuttosto che un corpus attestato. Questa assenza, come si è detto a proposito della tafonomia nel capitolo precedente, non equivale a una prova di assenza: equivale a documentare lo stato della ricerca, che su questo aspetto specifico non ha ancora prodotto un inventario sistematico.

% ============================================================
\section{Rovigo: il Polesine e la sua preistoria sommersa}
% ============================================================

La provincia di Rovigo presenta una situazione preistorica strutturalmente diversa da tutte le altre province venete, condizionata in modo determinante dalla sua geomorfologia: il Polesine è una pianura alluvionale depressa, compresa tra Po e Adige, che nel corso della preistoria era un paesaggio di paludi, dossi fluviali e lagune, in continua trasformazione a causa delle variazioni del livello marino e delle divagazioni fluviali. Questa geomorfologia ha una conseguenza diretta per l'archeologia: i siti più antichi, quando esistono, sono sepolti a profondità variabili sotto i sedimenti alluvionali e non sono accessibili con ricognizioni di superficie.

\subsection{Paleolitico e Mesolitico: assenza strutturale}

Non risultano siti paleolitici o mesolitici documentati con contesto stratigrafico nella provincia di Rovigo. Questa non è un'eccezione: è la norma attesa per un territorio che, nelle fasi più antiche del Pleistocene superiore, era il delta di un sistema fluviale maggiore, e nelle fasi del Mesolitico era probabilmente ancora in gran parte sommerso o lagunare, con la linea di costa a una posizione notevolmente più interna rispetto all'attuale. I depositi che potrebbero contenere testimonianze di frequentazione paleolitica si trovano a profondità non indagabili con archeologia convenzionale.

\subsection{Neolitico ed Eneolitico: la prima presenza documentata}

Le prime testimonianze di frequentazione documentata nel Polesine risalgono al Neolitico, con ceramica raccolta in superficie sui dossi fluviali residui (le uniche quote positive del territorio). La rivista \emph{Padusa}, pubblicata dal Centro Polesano di Studi Storici Archeologici ed Etnografici di Rovigo, costituisce il principale strumento di censimento della letteratura locale, con segnalazioni che si estendono dal Neolitico fino all'età medievale \citep{padusa}. La qualità della documentazione è però variabile: molte segnalazioni derivano da rinvenimenti fortuiti in lavori agricoli o canalizzazioni, senza documentazione stratigrafica.

\subsection{Età del bronzo: Frattesina e il distretto metallurigico}

L'età del bronzo recente e finale (XII-X secolo prima dell'anno zero) ha lasciato nel Polesine tracce di un'intensità eccezionale nel panorama dell'Italia settentrionale. Il sito di Frattesina di Fratta Polesine, scavato sistematicamente dall'Università di Ferrara, copre una superficie stimata in oltre venticinque ettari e ha restituito più di mille tombe, officine per la lavorazione dell'avorio, del vetro e dell'ambra, e la prima fornace vetraria documentata in Europa occidentale \citep{bietti_sestieri1996}. Il sito di Campestrin a Grignano Polesine ha restituito oggetti in ambra baltica, prova di contatti commerciali con un'area distante migliaia di chilometri. Il sito di Canàr a Castelnovo Bariano era un insediamento su palafitta lungo il Po, con tipologie ceramiche di contatto tra il mondo dell'Italia settentrionale e quello transalpino. Villamarzana ha restituito una necropoli con cremazioni della stessa fase.

\subsection{Età del ferro: San Basilio di Ariano e Adria}

San Basilio di Ariano Polesine è un emporio dell'età del ferro (VII-V secolo prima dell'anno zero) con testimonianze di ceramica greca, etrusca e venetica: un punto nodale nella rete commerciale adriatica \citep{bonomi1991}. Adria, che ha dato il nome all'Adriatico, è una città portuale con una lunga frequentazione preistorica e storica; i materiali più antichi dal suo territorio risalgono all'età del bronzo, ma è in età arcaica e classica che il sito raggiunge la massima importanza come porto di transito.

\subsection{Arte rupestre: impossibile per ragioni geomorfologiche}

Nel Polesine l'arte rupestre è strutturalmente impossibile: mancano rocce affioranti. Il substrato è ovunque coperto da sedimenti alluvionali profondi, e nessuna superficie rocciosa è accessibile in modo naturale in tutta la provincia. Questa non è una lacuna della ricerca: è una conseguenza della geologia. La ricchissima preistoria del Polesine si esprime attraverso insediamenti, necropoli e manufatti mobili, non attraverso arte su roccia.

% ============================================================
\section{Venezia: la laguna come risultato, non come sfondo}
% ============================================================

La provincia di Venezia richiede una premessa geomorfologica senza la quale i dati archeologici non si capiscono. La laguna che oggi identifica il territorio non è un paesaggio immutabile: è il risultato finale di un processo durato diecimila anni. Nel Tardoglaciale, attorno ai diciottomila anni prima dell'anno zero, il livello del mare era circa centoventi metri più basso dell'attuale e tutta l'area dell'Alto Adriatico era una vasta pianura emersa, percorsa dal Brenta, dall'Adige e da altri fiumi che defluivano verso un litorale molto più a est. La trasgressione Flandriana, tra i diecimila e i cinquemilacinquecento anni prima dell'anno zero, ha progressivamente sommerso quella pianura. Circa cinquemilacinquecento anni fa la linea di costa si trovava trenta chilometri più a ovest dell'attuale, e i dossi fluviali del Brenta (Mirano, Spinea, Noale, Martellago) erano già gli unici rilievi sabbiosi emergenti nell'area perilagunare. La laguna si è strutturata nella sua forma attuale solo tra quattromila e tremila anni prima dell'anno zero, in età eneolitica e del bronzo.

Questa cronologia ha una conseguenza diretta per l'archeologia: i siti più antichi, se sono stati occupati nelle fasi in cui il territorio era ancora pianura emersa, si trovano oggi sepolti sotto metri di sedimento lagunare o marino, non accessibili con metodi di ricerca convenzionali.

\subsection{Paleolitico e Mesolitico: assenza strutturale e tracce di superficie}

Non risultano siti paleolitici documentati nella provincia di Venezia; la spiegazione è la stessa del Polesine: i depositi pleistocenici sono sepolti sotto le alluvioni successive e non affiorano. Per il Mesolitico la situazione è parzialmente diversa: prospezioni di superficie a partire dal 1975 nella terraferma perilagunare, nell'area di Mestre, hanno individuato materiale litico in selce riferibile al Mesolitico finale (attorno ai seimilacinquecento anni prima dell'anno zero). Il limite di questi rinvenimenti è lo stesso che si incontra altrove nella pianura veneta: si tratta di materiale raccolto in superficie in terreni agricoli, senza contesto stratigrafico chiuso, e la sua interpretazione funzionale rimane problematica.

\subsection{Neolitico: Concordia Sagittaria}

Il sito di Concordia Sagittaria, in particolare la località Loncon, ha restituito nel 2012 un insediamento tardo-neolitico databile alla metà del quarto millennio prima dell'anno zero, con fosse domestiche, vasche e canali, seguito da una fase del bronzo antico. La pubblicazione in \emph{Notizie di Archeologia del Veneto} è uno dei pochi casi in cui un sito della provincia veneziana è documentato con scavo stratigrafico e cronologia assoluta. Per il resto del territorio neolitico, i dossi del Brenta (Mirano-Spinea in particolare) sono morfologicamente favorevoli ma non ancora indagati sistematicamente.

\subsection{Età del bronzo: le piroghe del Brenta meridionale}

Un ritrovamento del 2014 nel canale del Cornio a Campagna Lupia ha portato alla luce scafi monossili (piroghe in quercia) nell'antico alveo del Brenta meridionale, pubblicati sulla rivista del Museo Civico di Storia Naturale di Venezia \citep{asta2014}. Il contesto, con resti di castoro che indicano un ambiente forestale-ripario oggi scomparso, permette di fissare la frequentazione del corso d'acqua in un'epoca in cui la laguna non esisteva ancora nella sua forma attuale e i fiumi defluivano liberamente verso il mare.

\subsection{Età del ferro: Altino e i Veneti antichi}

Il sito più importante della provincia per la preistoria recente è Altino (Quarto d'Altino), che ha restituito una sequenza occupazionale continua dal bronzo recente (XIII-XII secolo prima dell'anno zero) fino alla romanizzazione. Tra il IX e l'VIII secolo prima dell'anno zero era già uno scalo fluviale venetico; tra il VII e il VI secolo aveva assunto una fisionomia urbana con un insediamento di circa quaranta ettari. Il Museo Nazionale di Altino, con la sua collezione di materiali venetici, è il riferimento principale per chi voglia capire cosa significa, in termini di produzione simbolica, la fase che precede immediatamente la romanizzazione in quest'area. Concordia Sagittaria presenta fasi venetiche analoghe, più a est.

\subsection{Arte rupestre: assente per ragioni strutturali}

Come nel Polesine, nell'intera provincia di Venezia non esistono affioramenti rocciosi naturali. L'arte rupestre è strutturalmente impossibile, non una lacuna della ricerca. La preistoria veneziana si esprime attraverso insediamenti, piroghe, ceramica e necropoli: non attraverso segni sulla roccia.

% ============================================================
\section{Treviso: la Marca e le Prealpi}
% ============================================================

% SEGNAPOSTO

\emph{Sezione in preparazione. La provincia di Treviso è inclusa nella pianura trevigiana e nelle Prealpi Trevigiane (Cansiglio, Cesen, Nevegal), con frequentazione mesolitica documentata nell'area del Cansiglio.}

% ============================================================
\section{Belluno: le Dolomiti e il Cadore}
% ============================================================

% SEGNAPOSTO: la ricerca sistematica del Cadore merita trattazione estesa

\emph{Sezione in preparazione. La provincia di Belluno, con il Cadore e le Dolomiti, è trattata parzialmente nel capitolo~\ref{cap:dove_guardare} a proposito dei casi di successo recente: l'Uomo di Mondeval, i siti mesolitici di Val Visdende e i ritrovamenti di Pian dei Buoi. Questa sezione integrerà il quadro con i siti del Bellunese orientale (Alpago, Zoldo, Agordino) e le grotte del massiccio delle Vette.}

% ============================================================
\section{Il Veneto geografico esteso: Garda trentino e lombardo, Primiero, Folgaria}
% ============================================================

% SEGNAPOSTO

\emph{Sezione in preparazione. La sponda trentina (Riva del Garda, Arco, Torbole) e la sponda lombarda settentrionale del Garda hanno restituito arte rupestre dell'età del bronzo in continuità con il complesso del Monte Luppia. Primiero e Folgaria sono inclusi per la continuità geografica con i valichi alpini frequentati in preistoria, non per ragioni amministrative.}

% ============================================================
\section{Un confronto numerico}
% ============================================================

I siti e i complessi descritti in questo capitolo acquistano significato pieno solo se messi accanto ai numeri delle aree più studiate d'Italia e d'Europa.

\begin{table}[htbp]
\centering
\small
\begin{tabular}{@{}p{2.6cm}p{1.6cm}p{1.8cm}p{1.6cm}p{2.4cm}@{}}
\toprule
\textbf{Sito} & \textbf{Rocce} & \textbf{Figure} & \textbf{Inizio ricerca} & \textbf{Cronologia proposta} \\
\midrule
Valcamonica (BS) & $\sim$2.400 & $>$140.000 & anni 1930 (sistematica dal 1964) & dal Neolitico al periodo storico \\
Monte Bego (Francia) & migliaia & decine di migliaia & primi '900 & età del bronzo \\
Garda/Monte Luppia (VR) & $>$250 & $\geq$3.000 & 1964 & età del bronzo -- epoca moderna \\
Val d'Assa (VI) & non quantificato & centinaia/migliaia (stima 1982) & 1966/1982 & prevalentemente medievale/moderna \\
Alto Lario (CO) & $\sim$300 & non quantificato & anni 1970 & età del ferro \\
\bottomrule
\end{tabular}
\caption{Confronto tra i principali complessi di arte rupestre alpina citati in questo libro. I dati sul Garda derivano da Gaggia (1982, 2002); quelli sulla Val d'Assa da Leonardi, Rigoni e Allegranzi (1982) e dalla revisione di Bettineschi e collaboratori (2022); quelli sulla Valcamonica dalla documentazione presentata all'UNESCO nel 1979 e da lavori successivi.}
\label{tab:confronto_siti}
\end{table}

Alcuni aspetti di questa tabella meritano un commento esplicito. La colonna delle figure per la Val d'Assa non è comparabile con le altre: mentre per Garda e Valcamonica il numero deriva da un catalogo sistematico di figure attribuite a un'epoca, per la Val d'Assa è una stima complessiva di tutti i segni visibili, indipendentemente dalla loro epoca; la maggior parte di quei segni è oggi ritenuta medievale o moderna. La colonna dell'inizio della ricerca mostra un divario che nessun'altra colonna può spiegare: la Valcamonica ha alle spalle quasi un secolo di attenzione scientifica; il Garda ha sessant'anni, iniziati da un'intuizione personale innescata dal contatto con la Valcamonica; la Val d'Assa ha un rilevamento iniziato per caso, in seguito a un'alluvione, e mai completato con la stessa sistematicità.

% ============================================================
\section{Cosa questo quadro dimostra, e cosa no}
% ============================================================

Va detto con la massima chiarezza possibile cosa questo capitolo non dimostra. Non dimostra che il Veneto fosse, in antichità, meno frequentato della Lombardia o delle Alpi occidentali: potrebbe esserlo stato, ma il dato numerico da solo non lo prova. Quello che il confronto mostra con certezza è una differenza nella storia della ricerca, misurabile e non solo impressionistica: alcune aree (Lessinia per il Paleolitico, distretto di Este per l'età del ferro, Polesine per l'età del bronzo) hanno ricevuto un'attenzione sistematica di livello europeo, e i risultati sono stati proporzionali all'investimento. Altre aree, in particolare la fascia alpina e prealpina centrale, hanno ricevuto ricerche discontinue o fortuite, e il quadro che ne emerge riflette questa discontinuità più che la realtà preistorica.

La Valcamonica ha alle spalle una scuola di studi continuativa, nata negli anni sessanta del Novecento attorno al lavoro di Emmanuel Anati e mai interrotta, con generazioni di ricercatori dedicati esclusivamente a quel territorio. Il Veneto non ha mai avuto un equivalente per l'arte rupestre. Le sue scoperte più importanti in questo campo, come si è visto, sono nate da episodi in larga parte fortuiti: una visita che accende un'intuizione, una piena che scopre una parete. Questa differenza, da sola, non spiega necessariamente tutto il divario numerico osservato. Ma è sufficiente per rendere prudente qualunque conclusione che dia per scontata l'assenza di ciò che, semplicemente, non è ancora stato cercato con la stessa intensità.


% Cap. 4 — Cosa avrebbero potuto lasciare
% Repertorio comparativo per periodo: cosa è attestato in contesti analoghi
% nel mondo (stesso clima, stessa organizzazione sociale, stesso tipo di
% ambiente), filtrato per plausibilità locale. Poi: condizioni di
% sopravvivenza e vettori di perdita (naturali e antropici).
\chapter{Cosa avrebbero potuto lasciare}
\label{cap:repertorio_simbolico}

Il capitolo precedente ha censito ciò che è stato trovato. Questo capitolo affronta una domanda diversa: a prescindere da quello che è stato trovato, cosa avrebbero potuto produrre le popolazioni che abitavano le Alpi venete, in base a quello che sappiamo di come vivevano, di dove vivevano, e di cosa hanno prodotto popolazioni analoghe in luoghi analoghi del mondo? La risposta non può essere una lista universale: deve essere filtrata caso per caso, periodo per periodo, applicando tre criteri distinti. Il primo è l'analogia: cosa hanno lasciato popolazioni con lo stesso tipo di organizzazione economica e sociale in altri contesti ben documentati? Il secondo è la plausibilità locale: delle cose prodotte altrove, quali erano compatibili con il clima, il substrato geologico e il paesaggio delle Alpi venete? Il terzo è la conservabilità: di ciò che era plausibile produrre, cosa ha una probabilità non nulla di essere sopravvissuto fino a oggi?

Applicare questi tre filtri in sequenza non produce una lista di siti da trovare: produce una lista di aspettative ragionevoli, che è già qualcosa di più utile dell'assenza totale di aspettative con cui di fatto si affronta la ricerca in molte delle aree descritte nel capitolo precedente.

% ============================================================
\section{Il Paleolitico superiore: grotte, ripari e oggetti portatili}
% ============================================================

\subsection{Cosa è attestato in contesti analoghi}

Il Paleolitico superiore europeo (tra 45.000 e 12.000 anni prima dell'anno zero, grosso modo coincidente con le fasi fredde descritte nel quarto capitolo) è il periodo per cui il confronto mondiale è più ricco e più preciso. In Francia e Spagna, la tradizione franco-cantabrica ha prodotto centinaia di grotte con pitture parietali complesse: Chauvet (circa 36.000 anni prima dell'anno zero), Lascaux (circa 17.000), Altamira (circa 18.500). In Germania meridionale, le grotte dello Jura svevo (Hohle Fels, Vogelherd) hanno restituito le figure intagliate in avorio di mammut più antiche del mondo, tra 40.000 e 33.000 anni prima dell'anno zero. In Italia, oltre a Fumane (già discussa), la Grotta Paglicci in Puglia e la Grotta del Romito in Calabria documentano arte parietale figurativa. In Portogallo, la Valle del Côa è l'esempio più importante di arte rupestre paleolitica all'aperto in Europa: figure di animali incise su roccia scistosa lungo il fiume, datate tra 22.000 e 10.000 anni prima dell'anno zero.

In tutti questi contesti, accanto all'arte parietale nelle grotte, c'è una produzione portatile costante: figurine, bastoni traforati, elementi di ornamento in conchiglia e dente di animale, ciottoli decorati. Il Riparo Dalmeri nella Valsugana orientale, già citato nel capitolo precedente, appartiene esattamente a questa categoria: sessanta ciottoli con pitture in ocra, in contesto epigravettiano a 13.000 anni prima dell'anno zero, sono la prova che questa produzione non si fermava ai Pirenei.

\subsection{Cosa è plausibile per le Alpi venete}

Il confronto con Fumane, Riparo Tagliente (arte mobilière), Riparo Villabruna (ciottolo inciso a 14.000 anni) e Riparo Dalmeri (pietre dipinte) dimostra che tutte e tre le principali forme di produzione simbolica paleolitica erano presenti nell'area: arte parietale in grotta, arte mobilière su osso e pietra, e pittura su ciottoli. La domanda non è se queste cose venissero prodotte, ma dove e in quale quantità.

Per l'arte parietale nelle grotte, i candidati geologicamente plausibili nelle Alpi venete sono le grotte calcaree della Lessinia (dove già si scava da decenni senza trovare pitture, ma le superfici non scavate sono ancora molte) e i Colli Berici. La tradizione franco-cantabrica non si estendeva fino alle Alpi orientali con la stessa intensità del versante atlantico: le pitture di Fumane sono eccezionali proprio perché la densità di siti figurativi decresce verso est. Ma Fumane esiste, il che significa che la pratica c'era.

Le incisioni all'aperto in stile paleolitico, come quelle del Côa, sono molto meno probabili: il Côa è un caso quasi unico in Europa, e la tradizione di incidere rocce all'aperto non sembra essere stata quella dominante per i cacciatori epigravettiani delle Alpi orientali. Non è impossibile, ma non è nemmeno l'aspettativa più ragionevole per questa fase.

Quello che con più certezza veniva prodotto, e di cui quasi certamente non resterà nulla, è la decorazione su materiali organici: legno, pelle, tendine, piume. Il confronto etnografico con i cacciatori-raccoglitori storici documentati in tutto il mondo suggerisce che la maggior parte della produzione simbolica delle società di caccia sia su supporti deperibili, e che l'arte su pietra e osso rappresenti solo la punta visibile di un sistema simbolico molto più ampio.

\subsection{Ambiente e stagionalità}

I siti paleolitici delle Alpi venete mostrano un pattern ricorrente: i ripari sotto roccia e le grotte vengono occupati in inverno o nelle stagioni di transizione, mentre i siti d'alta quota (come suggerisce Dalmeri, a quota relativamente bassa ma in posizione di transito) vengono usati in estate per seguire la fauna. L'arte parietale nelle grotte è associata ai siti di occupazione stabile invernale; l'arte mobilière e i ciottoli dipinti seguono il gruppo anche in quota. Se questo schema vale per le Alpi venete come vale per il Trentino confinante, allora le grotte da cercare per l'arte parietale sono in fondovalle o a quota intermedia, non sulle creste.

% ============================================================
\section{Il Mesolitico: mobilità, quota e supporti deperibili}
% ============================================================

\subsection{Cosa è attestato in contesti analoghi}

Il Mesolitico europeo (tra 11.500 e 6.500 anni prima dell'anno zero circa, con variazioni regionali significative) è il periodo meno ricco di arte conservata in tutto l'arco alpino, e probabilmente non per caso. Le società mesolitiche erano piccoli gruppi mobili di cacciatori-raccoglitori che seguivano la fauna in cicli stagionali verticali: le basse quote in inverno, i pianori e i passi d'alta quota in estate. Questa mobilità favorisce la produzione su supporti portatili e sfavorisce la produzione su roccia fissa.

Le eccezioni documentate sono istruttive proprio perché rare. In Scandinavia, i petroglifi di Alta (Norvegia settentrionale) includono una fase mesolitica con scene di caccia. Nella penisola iberica, alcune delle incisioni della tradizione levantina sono attribuite al Mesolitico. In Italia, i ciottoli dipinti aziliani sono stati trovati in contesti del Mesolitico antico in diverse grotte del nord: sono piccoli, portatili, e hanno lasciato tracce perché si trovavano in ripari protetti. Riparo Dalmeri, con i suoi ciottoli epigravettiani, è il precedente cronologicamente più vicino a questa tradizione.

\subsection{Cosa è plausibile per le Alpi venete}

Mondeval de Sora (sepoltura con corredo a 2.150 metri) e i siti d'alta quota del Cadore e del Cansiglio documentano la presenza mesolitica intensa sulle Alpi venete. A Plan de Frea in Val Gardena, i denti infantili in un campamento d'alta quota dimostrano che non si trattava di spedizioni di soli adulti maschi ma di movimenti di gruppi familiari. Questi gruppi producevano oggetti simbolici: il corredo di Mondeval include elementi di ornamento e strumenti decorati.

Quello che è meno probabile, ma non impossibile, è che abbiano lasciato incisioni su roccia in luoghi fissi. Non perché non ne fossero capaci, ma perché la logica della mobilità non favorisce l'investimento in segni permanenti su superficie fissa. Quando lo fanno, come in Scandinavia, tendono a scegliere punti geograficamente significativi: capi lacustri, guadi fluviali, valichi. Nelle Alpi venete, i passi dolomitici frequentati dai cacciatori mesolitici (Giau, Falzarego, Valparola) sono i candidati più ragionevoli se si cercano segni di questo tipo, anche se non è ciò che ci si aspetta di trovare con più probabilità.

% ============================================================
\section{Neolitico ed Eneolitico: i segni che durano}
% ============================================================

\subsection{Cosa è attestato in contesti analoghi}

Il Neolitico europeo (approssimativamente tra 7.500 e 5.000 anni prima dell'anno zero nelle Alpi) è il periodo in cui la produzione di arte su roccia diventa sistematica e diffusa in tutto l'arco alpino. Le coppelle (cavità emisferiche poco profonde picchiettate sulla roccia) sono la forma più ubiqua: si trovano su ogni tipo di roccia, ad ogni quota, in ogni ambiente. In Valle d'Aosta ve ne sono decine di migliaia. In Trentino-Alto Adige sono documentate ovunque ci sia calcare o cristallino affiorante. In Liguria e in Piemonte costellano i ripiani rocciosi dei valichi.

Accanto alle coppelle, il Neolitico alpino ha prodotto stele antropomorfe: le statue-menhir del Trentino meridionale (Arco, Riva del Garda, Cles), le stele della Val Venosta (Lagundo, San Martino al Monte), le statue della Lunigiana, le statue sarde. Alcune di queste si trovano a pochi chilometri dal confine veneto. La Valcamonica inizia la sua produzione documentata in questo periodo. Il Monte Bego, in Liguria, ha la sua prima fase.

\subsection{Cosa è plausibile per le Alpi venete}

Le coppelle sono praticamente certe, e non averle ancora censite sistematicamente è una lacuna della ricerca, non un'assenza delle coppelle. Ogni sistema montuoso calcareo del Veneto che è stato guardato con attenzione le ha trovate: la segnalazione non sistematica sugli Euganei (trachite) e sull'Altopiano di Asiago (calcare) non è la prova che non ci siano, è la prova che nessuno le ha ancora cercate in modo comparabile a come le ha cercate qualcuno in Valle d'Aosta o in Provenza.

Le stele antropomorfe sono meno probabili come oggetti conservati (il calcare si degrada, le stele vengono riutilizzate come materiale edilizio), ma potrebbero esistere frammenti non riconosciuti in contesti di reimpiego. Le stele trentine più vicine al confine veneto (Riva del Garda, Arco) mostrano che la tradizione era presente nell'area immediatamente adiacente.

L'eneolitico aggiunge la Cultura del Vaso Campaniforme, che ha una produzione simbolica ricca (decorazioni ceramiche, pugnali in selce con foderi decorati, ornamenti in rame e oro) ma quasi nessuna arte rupestre specifica attribuita: non è una fase in cui ci si aspetta una grande espansione delle incisioni.

% ============================================================
\section{Età del bronzo: il periodo più probabile}
% ============================================================

\subsection{Cosa sappiamo già}

Il Monte Luppia dimostra direttamente che nell'età del bronzo le popolazioni delle Alpi venete producevano arte rupestre complessa: guerrieri, armi, imbarcazioni, simboli solari, coppelle. Non è una questione aperta se lo facessero: è una questione aperta dove lo facessero oltre al Garda. La Valcamonica, il Monte Bego e le incisioni di Ala documentano che la produzione rupestre dell'età del bronzo copriva tutto l'arco alpino con una densità che non ha equivalenti nei periodi precedenti.

\subsection{Cosa è plausibile altrove nel Veneto}

Ogni massiccio calcareo delle Prealpi venete con superfici rocciose esposte e frequentazione dell'età del bronzo documentata è un candidato ragionevole. La Lessinia ha frequentazione dell'età del bronzo (Sant'Anna d'Alfaedo), superfici calcaree abbondanti, e nessuna campagna di rilevamento dell'arte rupestre comparabile a quella condotta sul Monte Luppia. L'Altopiano di Asiago ha frequentazione dell'età del bronzo e superfici calcaree: stessa situazione. I versanti del Grappa verso il Brenta, con la loro morfologia di creste e ripiani, non sono stati ancora indagati sistematicamente per questo periodo.

Il dato più rilevante per orientare la ricerca è che nell'età del bronzo l'arte rupestre alpina è associata a luoghi specifici con caratteristiche ricorrenti: superfici inclinate verso il cielo o verso un bacino visivo ampio, vicine a vie di transito, in prossimità di punti d'acqua, spesso a quota intermedia tra i fondivalle abitati e le creste di pascolo estivo. Questi criteri, applicati alla morfologia delle Prealpi venete, producono una lista di candidati più ristretta e più utile di un'ispezione casuale del territorio.

% ============================================================
\section{Cosa sopravvive e cosa no}
% ============================================================

\subsection{I fattori naturali}

La conservazione di un segno su roccia dipende da tre variabili fisiche che si combinano in modo non lineare: il tipo di roccia, l'esposizione alla meteorizzazione, e il tempo trascorso.

Il calcare è il substrato più comune nelle Prealpi venete e uno dei peggiori per la conservazione delle incisioni nel lungo periodo: è solubile nell'acqua piovana acida, si desquama per effetto del gelo e disgelo ripetuto (crioclastismo), e sviluppa patine biologiche (licheni, muschi, alghe) che possono cancellare i segni più superficiali. Le incisioni a martellina (picchiettatura profonda, come al Monte Luppia e in Valcamonica) resistono meglio di quelle a graffito (incisione leggera, come alla Val d'Assa). Dopo quattromila anni, un'incisione a martellina su calcare esposto è ancora riconoscibile; un graffito sottile potrebbe essere scomparso o ridotto a una traccia illeggibile.

La dolomia è chimicamente simile al calcare ma meccanicamente più fragile ai cicli termici: i versanti dolomitici sono dominati dalle frane e dai crolli, che periodicamente rinnovano la superficie e distruggono quanto inciso. Questo non significa che non esistano incisioni dolomitiche (esistono), ma che la probabilità di conservazione è inferiore a quella su calcare compatto.

Il porfido del Lagorai è una roccia magmatica, chimicamente molto più resistente e praticamente insolubile. Se qualcuno ha inciso il porfido, quella traccia sopravvive molto meglio del calcare. Il problema è che nessuno ha ancora cercato sistematicamente incisioni sul porfido del Lagorai: è una lacuna attiva, non una risposta.

L'esposizione conta quanto il tipo di roccia. Una superficie verticale riparata da un aggetto (un'overhang) è protetta dalla pioggia diretta e dall'accumulo di acqua; si conserva incomparabilmente meglio di una superficie orizzontale esposta, su cui l'acqua stagna, il gelo agisce e la vegetazione radica. Le incisioni di Valcamonica su superfici inclinate con aggetto naturale hanno resistito cinquemila anni. Le incisioni a graffito della Val d'Assa su superfici esposte rischiano di sparire nel corso di qualche secolo.

\subsection{I fattori umani antichi}

Prima dell'età moderna, la roccia incisa veniva riutilizzata. Una superficie con coppelle dell'età del bronzo poteva ricevere nuove incisioni nell'età del ferro, poi croci medievali, poi iniziali di pastori moderni: le stratificazioni sono la norma, non l'eccezione, in tutti i complessi rupestri alpini. Questo non cancella necessariamente i segni più antichi (a Valcamonica le sovrapposizioni vengono decifrate con tecniche di rilievo tridimensionale), ma rende più difficile la lettura e talvolta fisicamente copre segni precedenti.

La costruzione di strade, terrazzamenti agricoli, edifici ha rimosso o incorporato rocce incise in tutta la pianura veneta e lungo le valli principali. Nel caso della Lessinia, la coltivazione medievale intensiva delle superfici calcaree ha probabilmente cancellato o incorporato in murature gran parte di quanto poteva esistere a quote basse.

\subsection{I fattori umani moderni}

L'estrattivismo moderno (cave di calcare, di trachite, di porfido) ha rimosso fisicamente porzioni intere di paesaggio roccioso. Le cave di trachite euganea, attive per secoli, hanno eliminato superfici che non potranno mai essere indagate. Lo stesso vale per le cave calcaree della Lessinia e per le cave di porfido della Valsugana.

La Prima Guerra Mondiale merita un paragrafo a sé, che questo libro tratta in dettaglio nel capitolo successivo e in quello sulla tafonomia: è un vettore di distruzione con caratteristiche specifiche (concentrato su determinate creste, combinato con la difficoltà di accesso successiva) che non ha equivalenti negli altri fattori considerati finora.

Il turismo e le infrastrutture turistiche moderne hanno alterato o reso inaccessibili (paradossalmente, per eccessiva frequentazione) superfici che erano rimaste intatte per millenni. Il sentiero che consolida un pianoro d'alta quota con massicciata cementizia sigilla ciò che potrebbe trovarsi sotto.

\subsection{Cosa resta, in sintesi}

La combinazione dei fattori appena descritti produce un quadro in cui le condizioni di conservazione ottimale si concentrano in un insieme di situazioni specifiche: superfici verticali calcaree con aggetto, a quota intermedia (approssimativamente tra 600 e 1.500 metri sul livello del mare), in aree non coinvolte nell'agricoltura medievale intensiva, non raggiunte dalle opere militari della Prima Guerra Mondiale, non estratte da cave moderne. Nelle Alpi venete, questo profilo esclude quasi tutta la pianura (troppo antropicizzata), le creste prealpine occidentali (guerra), e i fondivalle principali (infrastrutture). Lascia come candidati più ragionevoli: i versanti laterali della Lessinia e del Monte Baldo a quote intermedie, alcune aree dell'Altopiano di Asiago non coinvolte negli scontri del Carso interno, la fascia prealpina della provincia di Treviso (Cesen, Col Visentin), e le zone dolomitiche bellunesi già indagate con successo per la frequentazione mesolitica.

Non è una lista di promesse. È una lista di dove il rapporto tra probabilità di conservazione e probabilità che qualcosa ci fosse è più favorevole della media. La differenza tra questa lista e il silenzio delle carte che descrivono il Veneto come privo di arte rupestre è la stessa differenza che il capitolo successivo chiamerà research bias.


% Cap. 5 — DA ESPANDERE: research bias + antropicizzazione + WWI
% Bias aggiuntivi rispetto alla versione attuale:
%   - bias da antropicizzazione moderna (pianura, turismo, cemento)
%   - WWI come caso speciale: trincee su creste prealpine, Grappa,
%     Pasubio, Marmolada, Altopiano di Asiago (alcuni reperti sopravvissuti)
\chapter{Una mappa logica della ricerca preistorica in Veneto}
\label{cap:research_bias}

La mappa della preistoria veneta che il capitolo precedente ha costruito provincia per provincia non è una mappa distorta: è una mappa parziale. La differenza è importante. Una mappa distorta riflette la realtà ma con proporzioni alterate; una mappa parziale ha zone illuminate con buona precisione e zone semplicemente buie, in cui il silenzio non dice nulla sulla realtà sottostante. Nelle zone illuminate, dove la ricerca è stata sistematica e continuativa, i risultati sono proporzionali all'investimento e sostanzialmente affidabili. Nelle zone buie, l'assenza di dati non è un dato: è l'assenza di chi ha cercato.

Questo capitolo non parla di distorsioni né di errori della ricerca passata. Parla della logica con cui la ricerca si è distribuita sul territorio veneto nel corso degli ultimi cento anni: perché alcune aree sono state illuminate, chi le ha illuminate, cosa le ha rese candidate alla ricerca sistematica; e perché altre aree sono rimaste nell'ombra, per ragioni che di solito non hanno nulla a che vedere con la loro importanza preistorica e molto a che vedere con la storia delle istituzioni, con la geografia logistica, e con eventi che hanno reindirizzato l'attenzione collettiva.

% ============================================================
\section{Dove la luce è caduta, e perché}
% ============================================================

\subsection{La scuola ferrarese e la Lessinia}

Il fatto più importante per capire la distribuzione della ricerca preistorica nel Veneto non è geografico: è istituzionale. L'Università di Ferrara, attraverso il suo dipartimento di preistoria, è la presenza dominante nella storia della ricerca preistorica veneta per quasi settant'anni. Alberto Broglio ha costruito a partire dagli anni sessanta un programma di ricerca sistematico attorno ai ripari della Lessinia che ha prodotto la sequenza paleolitica più completa dell'Italia nordorientale. Riparo Tagliente fu il sito-guida: una volta documentata la sua straordinaria continuità occupazionale dal Musteriano all'Epigravettiano, attirò finanziamenti, dottorati, collaborazioni internazionali, e spinse verso prospezioni nei siti vicini (Mezzena, Ghiacciaia, San Bernardino). Ogni grande scoperta in un'area aumenta la probabilità che vengano cercate le aree adiacenti: è un meccanismo noto in tutte le scienze empiriche, e in archeologia lo si vede funzionare con precisione nella Lessinia degli anni settanta e ottanta.

Lo stesso programma, esteso e rinnovato da Marco Peresani (allievo di Broglio), si è spostato al Cansiglio nel 1993 e ha prodotto in trent'anni di lavoro sistematico la sequenza epigravettiana-mesolitica più completa delle Prealpi venete: oltre 4.000 manufatti, archivi pollinici, siti di riferimento metodologico per tutta l'archeologia prealpina italiana. Al Cansiglio era già andata l'équipe di Broglio negli anni precedenti per le prime prospezioni. La continuità istituzionale è il filo: stessa scuola, stesso approccio metodologico, risultati comparabili per rigore e densità.

È la stessa scuola ferrarese, con Attilio Guerreschi prima e Federica Fontana poi, che ha scavato Mondeval de Sora a partire dal 1986 e ha scoperto la seconda sepoltura mesolitica al Passo Giau nel 2019. È ancora Broglio che ha diretto gli scavi di Fumane a partire dagli anni novanta, portando alla luce le figure aurignaziane che oggi collocano la grotta tra i siti di riferimento mondiale per l'origine dell'arte figurativa. Il dato ricorrente è questo: dove la scuola ferrarese ha portato il suo programma sistematico, ha trovato risultati di livello internazionale. Dove non è arrivata, la mappa rimane bianca.

\subsection{La Grotta di Fumane: una scoperta che ne genera altre}

Fumane merita un paragrafo separato perché illustra un meccanismo distinto da quello della Lessinia. Non era un'area già nota per i suoi siti paleolitici: la grotta era conosciuta ma non scavata sistematicamente. Fu il programma di Broglio, nel quadro più ampio delle indagini sulla Valpolicella, a portarla alla luce nella sua straordinaria rilevanza. Una volta pubblicati i frammenti dipinti aurignaziani, la grotta è entrata nel circuito internazionale dei siti di riferimento per l'arte paleolitica: le datazioni del 2025 con metodi bayesiani sono il prodotto diretto di quella visibilità, che ha portato collaborazioni con laboratori francesi e spagnoli impossibili senza il profilo internazionale acquisito.

Il meccanismo è circolare ma non vizioso: un sito importante attrae risorse che permettono di documentarlo meglio, che aumentano la sua importanza, che attraggono più risorse. Il problema è che questo circolo virtuoso si innesca solo se qualcuno ha già fatto la prima mossa, cioè ha dedicato tempo e denaro a uno scavo in un'area senza garanzie di risultato. Nel caso di Fumane, quella mossa è stata fatta da un programma universitario già finanziato per altri motivi. Senza quel programma, la grotta sarebbe probabilmente rimasta dove era.

\subsection{Il Monte Luppia: un uomo, una visita, un catalogo}

La storia del Monte Luppia è diversa da tutte le precedenti, e per questo è la più utile per capire le zone buie della mappa. Nel 1964 Mario Pasotti, insegnante, visitò le incisioni della Valcamonica. Al ritorno, riconobbe sulle rocce del Monte Luppia, vicino a casa sua, gli stessi segni che aveva appena visto altrove. Non era un ricercatore professionale, non aveva un programma istituzionale alle spalle, non disponeva di fondi universitari. Aveva uno sguardo addestrato dall'esperienza diretta con un sito già noto.

Fabio Gaggia, che ha poi costruito il catalogo sistematico del complesso e nel 1985 ha fondato un centro di studi dedicato, ha continuato questo lavoro per decenni, producendo il corpus di oltre duecentocinquanta rocce e tremila figure che è oggi il riferimento principale per il sito. Ma il centro di studi di Torri del Benaco non ha mai avuto la struttura, il personale e i finanziamenti di un dipartimento universitario. Il suo raggio di azione è rimasto essenzialmente il territorio benacense. Le Prealpi vicentine, i versanti del Grappa, la Lessinia calcaree al di fuori dei siti già noti di Tagliente e Mezzena: non sono mai stati oggetto di un programma di rilevamento dell'arte rupestre comparabile per sistematicità a quello che, negli stessi decenni, veniva condotto in Valcamonica o in Valle d'Aosta.

La domanda implicita è semplice: se il Monte Luppia è stato scoperto perché qualcuno con lo sguardo giusto ci ha passato, quante superfici analoghe nel Veneto non sono mai state guardate con lo stesso sguardo?

\subsection{Este, Frattesina, Altino: la ricchezza che si trova da sola}

Le scoperte dell'età del ferro e del bronzo nel Veneto meridionale seguono una logica diversa ancora. Este è nota da secoli: la ricchezza dei materiali dalle necropoli (situle, bronzi figurati, corredi funerari elaborati) l'ha resa oggetto di attenzione antiquaria già nell'Ottocento, molto prima che il concetto di scavo sistematico si affermasse. Le prime collezioni sistematiche risalgono agli anni 1870-1880, e il museo di Este è uno dei più antichi del Veneto. Frattesina è stata scoperta in modo più recente ma per ragioni analoghe: una densità di materiali da superficie così eccezionale (avorio, vetro, ambra baltica) da rendere impossibile ignorarla. Altino è emersa progressivamente da decenni di costruzione e bonifica.

Queste scoperte non richiedevano uno sguardo addestrato all'arte rupestre: richiedevano di guardare quello che emergeva dal terreno durante i lavori agricoli, e di riconoscerne il valore. Sono le scoperte che si trovano quasi da sole, perché i materiali sono abbondanti, visibili e inequivocabili. Il contrasto con l'arte rupestre è totale: un'incisione su calcare richiede condizioni di luce specifica, un occhio allenato, spesso strumentazione dedicata. Una situla di bronzo non richiede niente di tutto questo.

% ============================================================
\section{Le zone buie, e perché sono buie}
% ============================================================

\subsection{L'assenza dell'ancora istituzionale}

La variabile più predittiva della densità di siti noti in un'area è l'esistenza di un programma di ricerca continuativo con un'istituzione alle spalle. Dove c'è stata un'università, un museo con un direttore attivo, o anche un singolo ricercatore con la tenacia di Gaggia, la mappa si illumina. Dove non c'è stata nessuna di queste cose, la mappa è bianca.

Il Massiccio del Grappa è il caso più evidente. Morfologicamente è un candidato interessante per la frequentazione preistorica: creste, selle, versanti a quota variabile, vicinanza alla pianura veneta, con le sue necropoli venetiche a pochi chilometri. Ma non è mai stato oggetto di una campagna di prospezione preistorica sistematica comparabile a quelle condotte sulla Lessinia o nel Cadore. Nessuna università ha mai avuto un programma centrato sul Grappa per la preistoria. Nessun Gaggia locale ha mai fondato un centro di studi sulle rocce del Grappa. Il risultato è il silenzio, che non deve essere confuso con l'assenza.

Lo stesso vale per la fascia prealpina trevigiana (Col Visentin, Cesen), per gran parte del Bellunese fuori dal Cadore e dal Cansiglio, per quasi tutto l'arco prealpino vicentino fuori dalla Val d'Assa. Sono aree senza un'ancora istituzionale, dove la ricerca è rimasta episodica, affidata a segnalazioni occasionali o a ricognizioni di superficie mai sistematizzate.

\subsection{Le zone di guerra: un problema doppio}

Il Pasubio, il Carega, Tonezza del Cimone, l'Altopiano di Asiago, la zona del Grappa prealpino, la Marmolada: la Prima Guerra Mondiale ha trasformato fisicamente questi massicci in modo irreversibile, e ha prodotto un effetto secondario altrettanto importante per la ricerca preistorica: ha catturato l'attenzione collettiva e istituzionale per decenni.

Il problema fisico è reale e viene trattato in dettaglio nel capitolo successivo: trincee scavate nelle creste calcaree, strade militari costruite su versanti vergini, minature e crolli controllati, depositi di materiale bellico che rendono tuttora pericolose alcune zone. Tutto questo ha alterato superfici che erano rimaste intatte per millenni.

Ma il problema culturale è forse altrettanto grave. Dopo il 1918, questi massicci sono entrati nell'immaginario collettivo come luoghi della Grande Guerra, non come luoghi della preistoria. I musei dedicati all'Altipiano di Asiago sono musei della Prima Guerra Mondiale. Le guide turistiche del Pasubio descrivono le gallerie e le trincee. I fondi pubblici per la ricerca storica in queste aree sono andati a quella storia, non alla preistoria. L'attenzione istituzionale è stata assorbita dall'evento più recente e più drammatico, e ha lasciato la preistoria in un secondo piano da cui non è mai uscita.

L'unica finestra affidabile sullo stato di queste superfici prima della guerra sono le fonti anteriori al 1915: i bollettini del Club Alpino Italiano, le relazioni delle Società di Scienze Naturali, le guide alpine degli anni 1870-1914. Sono fonti che non cercavano arte rupestre in senso moderno, ma che descrivevano morfologie, ripari, creste, e talvolta segni curiosi su rocce, con la precisione del naturalista d'epoca. Nessuno le ha ancora sistematicamente rilette con gli occhi di chi cerca arte preistorica.

\subsection{L'alta Carnia e le aree di confine}

Sappada, la Val Visdende, l'alta Carnia (Sauris, Ampezzo, Forni) condividono con le zone di guerra un carattere diverso ma ugualmente limitante per la ricerca: sono aree di confine politico, rimaste a lungo in una zona grigia tra competenze istituzionali diverse. La preistoria della Carnia è studiata principalmente nell'ambito della preistoria friulana; quella del Cadore bellunese nell'ambito di quella veneta; quella della Pusteria nell'ambito di quella altoatesina. Le valli che collegano questi sistemi, e che in preistoria erano probabilmente le più frequentate proprio perché connettevano i diversi versanti, ricadono nelle giunture tra le competenze istituzionali e tendono a essere studiate meno di quanto la loro posizione geografica giustificherebbe. Il Lesachtal austriaco, documentato solo in letteratura di lingua tedesca non ancora sistematicamente recepita in italiano, è il caso estremo di questa frammentazione.

\subsection{Il problema specifico dell'arte rupestre}

L'arte rupestre richiede metodi di ricerca diversi da quelli usati per trovare un accampamento o una necropoli, e questo crea una lacuna specifica che si sovrappone a tutte le altre.

Trovare un sito con industria litica richiede di percorrere il terreno guardando il suolo: è un'abilità che ogni archeologo addestrato possiede. Trovare un'incisione su calcare richiede condizioni di luce molto specifiche (luce radente, di preferenza bassa sul versante, mai zenitale), occhio allenato a distinguere il segno intenzionale dall'erosione naturale, e spesso strumentazione dedicata: la riflettografia a trasformata d'immagine (RTI), il rilievo fotogrammetrico tridimensionale, la documentazione con luce strutturata sono oggi le tecniche standard per documentare superfici incise in modo che resista alla revisione critica. Nessuna di queste tecniche è standard nell'arsenale dell'archeologo generalista.

In Italia, i centri con competenza specifica nell'arte rupestre si contano su poche mani: il Centro Camuno di Studi Preistorici a Capo di Ponte, il MUSE di Trento, il gruppo di ricerca coordinato da Angelo Fossati al Politecnico di Milano, poche altre équipe. La Valcamonica ha un'istituzione dedicata che lavora dal 1964. Il Veneto non ha mai avuto niente di equivalente. Gaggia a Torri del Benaco è il riferimento per il Garda, ma il suo lavoro non è mai diventato un programma istituzionale esteso al territorio veneto nella sua interezza.

Il risultato concreto è che non sappiamo se l'arte rupestre sia assente nel Veneto prealpino o semplicemente non sia stata cercata con i metodi giusti. Sono due ipotesi empiricamente indistinguibili nello stato attuale della ricerca, e questo è il silenzio più significativo che la mappa ci consegna.

% ============================================================
\section{La correlazione che conta}
% ============================================================

Mettendo insieme la mappa di dove si è cercato e la mappa di cosa si è trovato, emerge una correlazione che sarebbe sbagliato ignorare: le aree con ricerca sistematica hanno prodotto risultati di rilievo internazionale (Lessinia, Fumane, Cansiglio, Cadore, Monte Luppia, Este, Frattesina); le aree senza ricerca sistematica producono silenzio. La correlazione è quasi perfetta.

Questo non significa che sotto ogni zona buia ci sia necessariamente un sito importante. Alcune zone sono buie per ragioni che hanno anche a che fare con la loro storia preistorica reale: il Polesine non ha arte rupestre non per assenza di ricerca ma per assenza di rocce affioranti; le zone di guerra hanno subito distruzioni fisiche reali; alcune aree alpine possono davvero essere state meno frequentate di altre per ragioni climatiche o ecologiche. La correlazione non è una prova, è un segnale.

Ma è un segnale abbastanza forte da giustificare un'inversione nell'onere della prova. Nel ragionamento corrente, chi afferma che il Veneto potrebbe contenere arte rupestre non ancora trovata deve giustificare questa ipotesi contro l'evidenza del silenzio. Nel ragionamento corretto, dopo aver visto il capitolo precedente e questo, è il contrario: chi conclude che il silenzio equivale all'assenza deve spiegare perché la correlazione tra ricerca e risultati, perfetta ovunque tranne che nelle zone buie, dovrebbe interrompersi proprio là.

\begin{table}[htbp]
\centering
\small
\begin{tabular}{@{}p{3.2cm}p{2.2cm}p{2.8cm}p{2.2cm}@{}}
\toprule
\textbf{Area} & \textbf{Tipo di ricerca} & \textbf{Istituzione di riferimento} & \textbf{Risultato} \\
\midrule
Lessinia & sistematica, continuativa & Università di Ferrara & sequenza di riferimento EU \\
Fumane & sistematica, continuativa & Università di Ferrara & sito mondiale Aurignaziano \\
Cansiglio & sistematica, continuativa & Università di Ferrara & 4.000+ manufatti, 30 anni \\
Cadore/Mondeval & sistematica, poi continuativa & Università di Ferrara & sepolture mesolitiche \\
Monte Luppia & sistematica locale & Centro studi Gaggia & 250 rocce, 3.000 figure \\
Val d'Assa & episodica, incompleta & nessuna & corpus incerto, non finito \\
Este/Frattesina & sistematica & Università di Ferrara & siti di livello europeo \\
Grappa & assente & nessuna & silenzio \\
Prealpi vicentine & assente/episodica & nessuna & silenzio \\
Prealpi trevigiane & assente & nessuna & silenzio \\
Zoldano & assente & nessuna & silenzio (dichiarato) \\
Zone di guerra & assente per preistoria & musei WWI & silenzio \\
\bottomrule
\end{tabular}
\caption{Distribuzione della ricerca preistorica nel Veneto e nei territori contigui, con l'istituzione di riferimento e il risultato prodotto. La correlazione tra tipo di ricerca e risultato è il dato più rilevante della tabella.}
\label{tab:mappa_ricerca}
\end{table}

% ============================================================
\section{Quattro tipi di silenzio}
% ============================================================

Il silenzio della mappa non è uniforme. Vale la pena distinguere almeno quattro tipi diversi, perché hanno implicazioni diverse per la ricerca futura e per le conclusioni che è lecito trarne.

Il \emph{silenzio strutturale} è quello del Polesine o della laguna veneziana: non c'è arte rupestre perché non ci sono rocce affioranti, e non ci sarà mai, qualunque sia l'intensità della ricerca. È il tipo di silenzio più informativo: dice qualcosa di certo, e non deve essere cercato.

Il \emph{silenzio tafonomico} è quello delle creste di guerra, delle cave, delle aree agricolizzate in profondità: poteva esserci qualcosa, ma le condizioni di conservazione erano così sfavorevoli da rendere la sopravvivenza improbabile. Non è certo come il silenzio strutturale, ma nemmeno è colmabile con più ricerca: ciò che è stato distrutto non può essere recuperato. Il capitolo che segue questo dedica un'analisi specifica a questo tipo di silenzio.

Il \emph{silenzio da assenza di ricerca} è quello del Grappa, dello Zoldano, della fascia prealpina trevigiana: nessuno ha cercato sistematicamente, e il silenzio non dice nulla sulla realtà sottostante. È il tipo di silenzio più interessante per la ricerca futura, perché è l'unico in cui nuova ricerca può fare la differenza.

Il \emph{silenzio da metodo inadeguato} è specifico dell'arte rupestre: aree dove si è cercata la frequentazione preistorica in senso generale (con survey di superficie per la litica, con sondaggi per le stratigrafie) ma senza le tecniche specifiche per il rilevamento dell'arte rupestre. In queste aree, il silenzio sull'arte non dice nulla perché è stato usato lo strumento sbagliato per la domanda giusta.

La differenza pratica tra questi quattro tipi è netta. Il silenzio strutturale chiude la domanda. Il silenzio tafonomico la ridimensiona. Il silenzio da assenza di ricerca la mantiene aperta. Il silenzio da metodo inadeguato la rimanda a nuove indagini con strumenti diversi.

% ============================================================
\section{Il confronto che illumina la lacuna}
% ============================================================

Due casi documentati permettono di calibrare meglio quanto pesi la differenza tra ricerca assente e ricerca sistematica, senza limitarsi alla logica interna al Veneto.

Il primo è il progetto Silvretta Historica nelle Alpi svizzere centrali e sudorientali, documentato da Marcel Cornelissen e Thomas Reitmaier \citep{cornelissen2016}. Prima del progetto, le Alpi svizzere centrali erano descritte nella letteratura come un'area con presenza mesolitica rada o assente. In otto anni di ricerca sistematica, il progetto ha censito oltre 230 siti sopra i 2.000 metri in un'area di 550 chilometri quadrati, con una presenza mesolitica che gli autori definiscono ``diffusa e quasi continua''. Non è un ribaltamento della realtà: è la realtà che emerge quando viene cercata con i metodi giusti per il tempo giusto.

Il secondo è lo studio di Davide Visentin e Francesco Carrer sulle Dolomiti venete, pubblicato nel 2017 \citep{visentin2017}, che ha analizzato con modelli statistici di punto quanto della distribuzione nota dei siti mesolitici dipenda dal comportamento insediativo antico e quanto dallo sforzo di ricerca. Il risultato è più sfumato del primo caso: i due fattori, ambientale e da sforzo di ricerca, spiegano in misura pressoché uguale la distribuzione osservata. Né tutto bias, né tutto distribuzione reale: una combinazione dei due, in proporzioni che solo la ricerca sistematica nelle zone ancora bianche potrà chiarire.

Questi due casi non dimostrano che il Veneto nasconda una densità di arte rupestre paragonabile alla Valcamonica. Dimostrano che la correlazione tra intensità della ricerca e risultati è empiricamente documentata, non solo teoricamente plausibile; e che l'inversione del rapporto causa-effetto (l'assenza di risultati come causa, e non come effetto, della scarsità di ricerca) è un errore documentato, non solo un rischio astratto.

% ============================================================
\section{Una distinzione che vale la pena tenere ferma}
% ============================================================

I due casi discussi nell'ultima sezione, come la maggior parte della ricerca citata nel capitolo precedente, documentano la frequentazione preistorica attraverso industria litica, fauna, resti ossei, strutture. Non documentano arte rupestre. Il processo che porta a trovare un accampamento mesolitico è diverso da quello che porta a trovare un pannello inciso: diversi strumenti, diversi criteri di riconoscimento, diverse condizioni ambientali favorevoli.

Questa distinzione non indebolisce il ragionamento di questo capitolo: lo precisa. Il silenzio sull'arte rupestre nel Veneto prealpino non è spiegato solo da ciò che non è stato trovato in termini di frequentazione, ma da una lacuna specifica nella ricerca dedicata all'arte rupestre in quanto tale. Sapere che la frequentazione mesolitica del Cadore era sottostimata non permette di dedurre che lo sia anche l'arte rupestre del Cadore: sono domande separate che richiedono risposte separate. Ma sapere che la ricerca sistematica in una zona tende a trovare quello che cerca, e che la ricerca sistematica specificamente dedicata all'arte rupestre non è mai stata condotta in modo ampio sul territorio veneto, è sufficiente per mantenere aperta la domanda di questo libro, senza presupporne la risposta in nessuna direzione.


% Cap. 5 — DA ESPANDERE: tafonomia naturale + tafonomia antropica
% Collega i vettori di cancellazione: tempo/clima (come ora) +
% azione umana (deforestazione medievale, cave, vie militari, WWI)
\chapter{Quello che il tempo cancella}
\label{cap:tafonomia}

I due capitoli precedenti hanno mostrato la distribuzione attuale dell'arte rupestre nel Veneto e le ragioni per sospettare che quella distribuzione rispecchi in parte la storia della ricerca. Ma c'è una terza variabile, distinta dal research bias e dalla reale frequentazione antica, che condiziona in modo indipendente ciò che oggi è possibile trovare: la conservazione nel tempo delle tracce stesse. Prima di chiedersi dove cercare, vale la pena chiedersi cosa sopravvive, e per quanto, e perché.

\section{Due storie di conservazione molto diverse}

Il confronto tra il deserto di Atacama e le Alpi venete non è solo un confronto tra territori più o meno esplorati. È anche, e prima ancora, un confronto tra due ambienti di conservazione radicalmente diversi.

Il deserto di Atacama è uno degli ambienti più aridi del pianeta. In larghe fasce della regione le precipitazioni annue sono inferiori a un millimetro; in alcune stazioni meteorologiche non si registra pioggia per decenni consecutivi. In queste condizioni, i principali agenti di degrado di un'incisione o di una pittura rupestre, che altrove operano senza sosta, qui sono quasi inattivi: i cicli di gelo e disgelo sono rari o assenti alle quote rilevanti; la vegetazione che colonizza le superfici rocciose è quasi inesistente; i licheni, che in ambienti umidi avanzano di qualche millimetro all'anno, qui crescono a velocità irrisorie. Una pittura rupestre in questo contesto non ha molti nemici. Può invecchiare senza essere cancellata.

Le Alpi venete, in confronto, sono un ambiente di degrado attivo. Le precipitazioni annue superano i 1.000 millimetri su gran parte del territorio montano; i cicli di gelo e disgelo sono frequenti in quota per diversi mesi all'anno; la vegetazione ricolonizza rapidamente ogni superficie esposta, compresi i licheni crostosi che aderiscono alla roccia e ne accelerano la disgregazione chimica; le pareti rocciose non protette sono soggette a stillicidio, scivolamento di depositi superficiali, distacco di frammenti per crioclastismo. Un'incisione poco profonda in questo contesto ha la vita breve, anche in assenza di qualsiasi intervento umano.

\section{L'asimmetria temporale}

C'è però un'asimmetria ancora più fondamentale tra i due confronti, che è facile trascurare se ci si concentra troppo sul confronto geografico: i due fenomeni non si collocano nello stesso periodo.

Le pitture e i petroglifi del deserto di Atacama, nella grande maggioranza dei casi studiati, risalgono ai millenni compresi tra 3.000 e 500 anni prima dell'anno zero. Le prime corrispondono all'incirca all'apogeo delle culture regionali pre-incaiche; le più recenti si sovrappongono al periodo di espansione dell'impero inca e ai primi secoli del contatto europeo. Anche i siti più antichi attribuiti alla tradizione rupestre atacamena non superano in genere i 4.000--5.000 anni prima dell'anno zero. E, aspetto cruciale, nell'intero deserto di Atacama non sono ancora state identificate con certezza pitture o incisioni databili al periodo in cui le Alpi venete ci interessano di più: il Tardoglaciale e il Mesolitico, tra 17.000 e 7.500 anni prima dell'anno zero.

Questo non significa che il territorio fosse disabitato nel Tardoglaciale. Le miniere di ocra nell'area di Taltal, come discusso nel primo capitolo, attestano una presenza umana che risale probabilmente a 12.000--10.000 anni prima dell'anno zero. Ma quella presenza, ammesso che abbia lasciato segni sulle rocce, non ha lasciato segni che oggi sappiamo riconoscere o attribuire con sicurezza a quel periodo. I motivi possono essere molti: quei segni esistevano e sono stati cancellati dal tempo; esistevano e sono stati sovrapposti da segni successivi; non sono stati ancora cercati con metodi adeguati; oppure, più semplicemente, la pratica di incidere e dipingere la roccia nel modo che oggi riconosciamo non era ancora diffusa in quella fase.

In altri termini, il confronto diretto tra la ricchezza di arte rupestre del deserto di Atacama e la scarsità di quella veneta non è del tutto simmetrico nemmeno sul piano temporale: la ricchezza atacamena si riferisce a un arco di 3.000--5.000 anni relativamente recenti e in un ambiente di conservazione eccezionale; la scarsità veneta si riferisce a un arco cronologico molto più lungo e in un ambiente di conservazione sfavorevole. Non sono la stessa cosa, e mescolarle in un unico confronto rischia di produrre un'impressione più netta di quanto i dati giustifichino.

\section{Tre modi di non trovare un sito}

Questa asimmetria porta a distinguere tre situazioni che in superficie sembrano identiche (il sito non c'è), ma che hanno cause e implicazioni molto diverse.

La prima è l'assenza reale: in questo luogo non è mai stato inciso o dipinto nulla di intenzionale, o comunque nulla che oggi rientrerebbe nella categoria di arte rupestre. Può dipendere da ragioni culturali (le comunità che frequentavano questo territorio non praticavano questa forma di espressione), da ragioni pratiche (le rocce disponibili non erano adatte alla tecnica usata) o semplicemente dal caso.

La seconda è l'assenza apparente: il segno c'era, ma è stato cancellato o reso illeggibile dai processi di degrado, dall'erosione, dalla copertura vegetale, dal lichenamento, da crolli o da riutilizzi successivi. In questo caso l'assenza non è nel passato, ma nella catena tra il passato e il presente.

La terza è l'assenza non ancora osservata: il segno c'è, ma non è ancora stato trovato, o perché nessuno ha cercato in modo sistematico in quel luogo, o perché le condizioni di visibilità al momento della ricerca erano sfavorevoli, o perché chi ha cercato usava metodi non adatti a quel tipo di traccia.

Il research bias discusso nel secondo capitolo di questo libro riguarda principalmente la terza categoria. Ma il problema della tafonomia rupestre, cioè del degrado e della cancellazione delle tracce nel tempo, riguarda la seconda. Le due categorie non si escludono: un'area può essere sia poco indagata (terza categoria) sia fisicamente sfavorevole alla conservazione (seconda categoria). Confonderle porta però a conclusioni diverse: nella seconda, anche la ricerca più intensa non restituirebbe i siti perduti, perché non ci sono più; nella terza, la ricerca sistematica può ancora cambiare il quadro.

\section{I processi di degrado specifici per l'arte rupestre alpina}

Per l'area alpina, i principali processi che minacciano la leggibilità di un'incisione o di una pittura rupestre all'aperto sono documentati, anche se raramente quantificati in modo sistematico.

Il crioclastismo, già discusso nel terzo capitolo a proposito della grotta di Fumane, è il processo più generale: l'acqua penetra nelle microfessure della roccia, ghiaccia, si espande, e stacca frammenti. Sulle pareti verticali esposte, questo processo può rendere illeggibile un'incisione in poche centinaia di anni, specialmente se poco profonda. Un segno inciso a quattro millimetri di profondità resiste molto meno di un segnato a dodici; ma anche i segni profondi, se la parete si frattura lungo il piano dell'incisione, spariscono con il distacco dello strato superiore.

La lichenizzazione è un secondo processo difficile da contrastare. I licheni crostosi, in particolare le specie del genere \textit{Rhizocarpon} comuni sulle rocce calcaree alpine, colonizzano le superfici a velocità lenta ma continua, crescendo all'interno degli strati superficiali della roccia e disgregandola chimicamente attraverso la produzione di acidi organici. Un pannello ricco di licheni diventa progressivamente illeggibile non perché sia coperto da uno strato opaco, ma perché la superficie stessa della roccia, con le incisioni su di essa, viene consumata e polverizzata.

L'erosione idrica agisce in modo diverso ma convergente: le precipitazioni abbondanti sciolguono la roccia calcarea, levigando le superfici e riducendo le differenze di rilievo che rendono leggibili le incisioni. Sulle superfici sub-orizzontali esposte all'acqua corrente il processo è particolarmente rapido; sulle pareti verticali è più lento, ma agisce comunque nel tempo.

Infine, la copertura da parte di sedimenti, detrito vegetale, muschi e suolo organico può rendere invisibile un pannello anche senza cancellarlo: se il piano di campagna si alza di pochi decimetri nel corso dei secoli, porzioni di parete prima visibili diventano sepolte, e i segni che ospitavano rimangono fisicamente intatti ma invisibili a chi cammina sulla superficie attuale.

\section{Cosa può ancora essere trovato, e cosa no}

Questa rassegna non intende concludere che la ricerca sia inutile. Intende chiarire, per onestà metodologica, che la ricerca può cambiare il quadro nella misura in cui le tracce sono fisicamente sopravvissute e fisicamente accessibili.

Per le fasi più antiche, quelle mesolitiche e tardoglaciali che costituiscono il periodo di maggiore interesse di questo libro, il problema della conservazione è più acuto che per le fasi più recenti. Un'ipotetica incisione risalente a 10.000 anni prima dell'anno zero ha avuto dieci millenni di cicli di gelo e disgelo, di lichenazione, di erosione idrica e di variazioni del piano di campagna per diventare illeggibile. La probabilità che sopravviva in forma leggibile è bassa, anche se non nulla: esistono casi, in altri contesti alpini, di incisioni attribuite al Mesolitico in base al contesto stratigrafico, non solo allo stile.

Per le fasi più recenti, dall'età del bronzo in poi, la finestra di sopravvivenza è più ampia. I siti noti del Garda e della Val d'Assa rientrano in questa categoria, e mostrano che la sopravvivenza, almeno parziale, è possibile.

Questa distinzione temporale ha una conseguenza pratica per il tipo di ricerca da fare: cercare arte rupestre mesolitica nelle Alpi venete richiede non solo un metodo di prospezione sistematico, come quello di Kompatscher e Hrozny Kompatscher discusso nella terza parte di questo libro, ma anche metodi di rilevamento adattati a tracce parzialmente degradate o parzialmente coperte, come l'imaging multispettrale, la fotogrammetria a luce radente o la rimozione controllata di sedimenti superficiali. E richiede di accettare in anticipo che una parte delle tracce, forse la maggior parte, non sarà più trovabile, indipendentemente dall'intensità della ricerca.

\section{Un'ipotesi alternativa che non va ignorata}

Fin qui questo capitolo ha trattato la tafonomia come il principale fattore che spiega l'assenza di arte rupestre mesolitica nelle Alpi venete: i segni c'erano, ma sono stati cancellati o sono irraggiungibili. È un'ipotesi ragionevole, ma non è l'unica possibile, e sarebbe intellettualmente disonesto non enunciarla chiaramente.

L'ipotesi alternativa è che i cacciatori mesolitici delle Alpi venete semplicemente non producessero arte rupestre all'aperto su larga scala, o che non la producessero affatto. Non perché fossero culturalmente inferiori o privi di capacità simboliche: Fumane dimostra il contrario per l'arte parietale in grotta; la sepoltura di Mondeval, con il suo corredo di ocre e conchiglie, dimostra pratiche rituali elaborate per il Mesolitico. Ma l'arte parietale in grotta e l'arte rupestre all'aperto sono pratiche diverse, che non si implicano necessariamente a vicenda.

In molte regioni d'Europa il Mesolitico è quasi del tutto privo di arte rupestre all'aperto, e questo non è spiegato solo dal degrado differenziale: alcune delle stesse rocce che ospiteranno incisioni dell'età del Bronzo erano già esposte e accessibili durante il Mesolitico. La ragione più probabile è che la tradizione di marcare le rocce all'aperto in modo visibile e permanente sia associata, nella maggior parte dei contesti europei, a società più sedentarie, che hanno interesse a segnalare in modo duraturo il proprio rapporto con un territorio specifico. Cacciatori altamente mobili possono usare altre forme di comunicazione simbolica, materiali deperibili (pitture su corteccia, incisioni su legno, segnali al suolo) che non lasciano tracce archeologiche.

Questa ipotesi non è falsificabile in senso stretto con i dati attuali: non si può dimostrare che qualcosa non sia mai esistito. Ma è corretto tenerla sul tavolo accanto all'ipotesi della tafonomia, perché le due hanno implicazioni di ricerca diverse. Se si pensa che l'arte c'era e sia sparita, ha senso investire in tecniche di rilevamento avanzate. Se si pensa che non ci fosse, ha più senso orientare la ricerca verso le fasi cronologiche successive, Neolitico ed età del Bronzo, dove la tradizione di arte rupestre all'aperto è attestata anche in contesti alpini comparabili. Il capitolo~\ref{cap:storia_territorio} di questo libro discute entrambe le fasi in prospettiva storica.


% ============================================================
% PARTE IV — Dove guardare
% ============================================================
\part{Dove guardare}

% Cap. 9 — Dove guardare: logica del movimento, metodo Kompatscher,
% prospezioni sistematiche alpine, analisi per periodo
% (accorpa e sostituisce i precedenti capp. 07–12)
\chapter{Dove guardare}
\label{cap:dove_guardare}

La domanda da cui è partito questo libro si è progressivamente trasformata. Non è più solo "perché il Veneto sembra privo di arte rupestre?" ma qualcosa di più preciso: quali tipi di luoghi, in questo territorio, non sono mai stati cercati con i metodi giusti? La risposta non è una lista di coordinate, ma una logica: capire come si muovevano gli animali cacciati, come si muovevano i cacciatori, come funzionano i metodi sistematici sviluppati altrove, e come cambia tutto questo al variare del periodo. Ogni epoca ha una firma diversa nei depositi rupestri, e ogni firma richiede strumenti di ricerca adattati.

% ============================================================
\section{La logica del movimento: animali e uomini}
% ============================================================

Il primo errore da evitare è ragionare per corridoi. L'immagine del cacciatore preistorico che segue la selvaggina lungo vallate continue, risalendo dalle quote invernali alle quote estive secondo percorsi prevedibili, è intuitiva ma empiricamente fragile, almeno per le specie che hanno dominato la fauna alpina veneta dalla fine del Pleistocene in poi.

Il camoscio (\emph{Rupicapra rupicapra}) e lo stambecco alpino (\emph{Capra ibex}) non sono animali di lunga migrazione. Studi radiotelemetrici su popolazioni alpine contemporanee indicano home range che raramente superano i venti chilometri quadrati per il camoscio e i trenta per lo stambecco; in condizioni di densità normale, questi animali non abbandonano la fascia altitudinale che conoscono, e i movimenti stagionali avvengono entro pochi chilometri verticali, non lungo centinaia di chilometri di versante. Questo non significa che le valli non venissero attraversate, ma che l'attraversamento non era il pattern dominante: l'arte rupestre legata alla caccia segue la distribuzione spaziale degli animali, non i desideri del cacciatore di muoversi velocemente.

Il secondo errore simmetrico è immaginare che il cacciatore preistorico si muovesse senza vincoli energetici. Il modello ARC (Accounting for Roughness and Cost) sviluppato da Herman Pontzer e collaboratori \citep{pontzer2016} quantifica il costo metabolico della locomozione su terreno irregolare per mammiferi di taglia molto diversa. Per un essere umano adulto, il costo energetico di un dislivello positivo è circa quattro volte quello dello spostamento in piano a parità di distanza orizzontale; il costo del dislivello negativo è circa il doppio. Su un percorso alpino tipico con cento metri di salita e ottanta di discesa per chilometro orizzontale, questo si traduce in un aumento del costo totale dell'ordine del quaranta per cento rispetto al piano. Le scelte di localizzazione dei siti non seguono la distanza lineare ma il costo del tragitto.

Uniti, questi due vincoli producono una predizione verificabile: i siti si addenseranno in luoghi che soddisfano contemporaneamente la distanza di costo accettabile dalla base (Kompatscher e Hrozny Kompatscher parlano di due ore di marcia, che su terreno alpino corrispondono a distanze assai variabili a seconda del dislivello), la visibilità sulle zone di pascolo degli animali cacciati, e la disponibilità d'acqua a meno di ottanta metri. La combinazione dei tre fattori restringe enormemente il numero di candidati su qualsiasi cartografia topografica.

% ============================================================
\section{Il metodo Kompatscher: nato dal camminare, non dal calcolare}
% ============================================================

Il metodo messo a punto da Karl Kompatscher e Margit Hrozny Kompatscher nel corso di decenni di lavoro sulle Alpi trentine e alto-atesine non nasce da un modello teorico, ma dalla sistematizzazione di un'esperienza accumulata sul campo \citep{kompatscher2007}. I due ricercatori hanno censito oltre trecento siti in Trentino-Alto Adige applicando una griglia di criteri ricavati a posteriori dall'analisi dei siti già trovati, e verificati poi su aree nuove.

Il metodo si regge su due parametri fondamentali. Il primo è il bacino di due ore: un sito di abitazione temporanea, per essere funzionale, deve essere raggiungibile in circa due ore di marcia da un punto d'acqua affidabile e deve permettere di raggiungere, nelle stesse due ore, le zone di caccia principali. L'area entro due ore di marcia su terreno alpino ha forma irregolare, dipendente dal dislivello, e in media copre una superficie di cinque-quindici chilometri quadrati attorno al sito. Il secondo parametro è il numero di direttrici: un sito è preferenziale se da esso partono almeno due percorsi distinti verso zone di caccia o verso valichi, perché riduce l'imprevedibilità dell'esito di una singola battuta.

Da questi due parametri derivano tre tipi di sito. Il campo base (campo di sosta prolungata, spesso con acqua a venti-ottanta metri, morfologia riparata, tracce di combustione) è il nodo della rete; il sito di caccia e avvistamento è un punto panoramico o un valico da cui si controllano i movimenti degli animali; il punto di sosta temporanea è un riparo minimo usato durante la battuta. L'arte rupestre, quando è presente, tende a concentrarsi nei campi base e nei siti di avvistamento, raramente nei punti di sosta breve.

La quota dell'acqua è un filtro molto selettivo. Su terreno calcareo carsificato, le sorgenti emergono a quote precise legate alla stratigrafia; su terreno cristallino, i ruscelli di fusione nivale funzionano stagionalmente. La distanza di venti-ottanta metri non è arbitraria: al di sotto di venti metri l'acqua è troppo vicina per garantire visibilità sufficiente verso le zone di caccia (il sito è in fondo valle), al di sopra di ottanta metri il costo di approvvigionamento diventa penalizzante nelle operazioni quotidiane di un campo stagionale.

Ciò che è utile di questo metodo per il Veneto non è la lista dei siti trentini trovati, ma la logica sottostante: cercare su cartografia i punti che soddisfano contemporaneamente più direttrici di movimento, acqua a distanza ragionevole, riparo morfologico. Questi punti si possono identificare prima di andare sul campo, riducendo il territorio da esplorare da migliaia a decine di ettari.

% ============================================================
\section{Le prospezioni sistematiche nelle Alpi: tre casi di riferimento}
% ============================================================

Il metodo Kompatscher non è un caso isolato. Tre programmi di ricerca condotti in contesti alpini diversi mostrano convergenze significative nella logica e divergenze istruttive nei risultati, e costituiscono insieme un quadro di riferimento per capire cosa può aspettarsi una prospezione sistematica in Veneto.

\subsection*{Silvretta Historica: 230 siti in 550 chilometri quadrati}

Il programma Silvretta Historica, coordinato da Martin Cornelissen e Tobias Reitmaier tra il 2009 e il 2016, ha coperto sistematicamente circa 550 chilometri quadrati di territorio alpino al confine tra Austria e Svizzera, a quote prevalentemente superiori ai duemila metri \citep{cornelissen2016}. Il risultato è stato il censimento di oltre 230 siti, il novanta per cento dei quali sconosciuto prima dell'avvio del programma. Nessuno di essi era stato trovato nelle prospezioni precedenti, nonostante l'area fosse attraversata da itinerari escursionistici e studiata da ricercatori locali.

La ragione non è che i siti fossero nascosti: è che le prospezioni precedenti non erano state progettate per trovarli. Un camminatore che percorre un sentiero vede le rocce a meno di tre metri dalla traccia; un ricercatore che progetta una prospezione sistematica a transetti paralleli di trenta-cinquanta metri copre l'intero territorio. La Silvretta ha prodotto non solo un catalogo, ma la dimostrazione empirica che un'area descritta come "priva di tracce" prima di una prospezione sistematica può rivelarsi densamente frequentata. Per il Veneto, questo è il riferimento metodologico più diretto: l'assenza dalla letteratura non è prova di assenza dal territorio.

\subsection*{Tirolo e Carinzia: le revisioni sistematiche}

In Alto Adige, Tirolo e Carinzia, una serie di ricercatori ha sviluppato programmi di revisione dell'esistente e di prospezione di nuove aree, con approcci parzialmente diversi da Kompatscher ma convergenti nella logica. Franz Mandl in Stiria \citep{mandl2002} e Walter Leitner in Tirolo \citep{leitner2013} hanno documentato centinaia di siti di quota (principalmente caccia, accampamenti stagionali, strutture pastorali di varie epoche) attraverso prospezioni che privilegiano la copertura sistematica rispetto alla verifica di segnalazioni pregresse.

Ciò che emerge da questi lavori è una distribuzione dei siti che conferma la logica di Kompatscher: i siti si addensano intorno a valichi, in spianate protette dal vento nelle vicinanze di sorgenti, e su punti panoramici con visibilità multipla. La variazione rispetto al modello Kompatscher riguarda principalmente l'adattamento a morfologie diverse (rocce cristalline invece di calcareo-dolomitiche, fondovalli più ampi). Per il Veneto, che condivide con queste aree sia morfologie carbonatiche (Prealpi, Dolomiti) sia aree cristalline (Lagorai, Cima d'Asta), entrambi i modelli sono pertinenti.

\subsection*{Scandinavia: RTI e metodi di rilevamento per l'arte rupestre}

La ricerca sull'arte rupestre scandinava, concentrata soprattutto in Norvegia (Alta, sito UNESCO) e Svezia (Nämforsen, Bohuslän), ha sviluppato negli ultimi vent'anni una metodologia di rilevamento che va oltre la semplice prospezione topografica. La tecnica RTI (Reflectance Transformation Imaging) \citep{drap2013}, applicata sistematicamente a pannelli rupestri, permette di rivelare incisioni superficiali invisibili in luce naturale diretta e non rilevabili con la pulizia meccanica. La stessa tecnica è stata applicata con risultati significativi a Creswell Crags (Inghilterra), dove incisioni paleolitiche sono state identificate su pareti che erano state studiate per decenni senza rilevarle.

La fotogrammetria strutturata (Structure from Motion, SfM) permette invece la costruzione di modelli tridimensionali delle superfici rocciose, utile sia per documentare siti già noti sia per identificare morfologie superficiali compatibili con incisioni profonde in condizioni di luce non favorevole. La luce radente artificiale (raking light), nella versione standardizzata sviluppata in ambito scandinavo, permette campagne notturne che rivelano segni non visibili di giorno.

Questi metodi non sostituiscono la prospezione pedestr, ma la integrano: una volta identificata tramite criteri Kompatscher un'area prioritaria, RTI e fotogrammetria permettono di esaminare sistematicamente le superfici rocciose senza dipendere dalle condizioni di illuminazione naturale.

% ============================================================
\section{Il Paleolitico superiore: grotte, rocce e clima}
% ============================================================

Il Paleolitico superiore nelle Alpi venete (approssimativamente da trentacinquemila a undicimilacinquecento anni prima dell'anno zero) è caratterizzato da condizioni ambientali radicalmente diverse da quelle attuali. Le fasi più fredde (Glaciale Massimo, intorno a ventiquattromila anni prima dell'anno zero) vedono il limite delle nevi perenni scendere fino a seicento-ottocento metri sulle Prealpi; la fauna include renna, rinoceronte lanoso, mammut nelle fasi più antiche, e camoscio, stambecco e cervo nelle fasi di interstadio. La tecnologia è epigravettiana (lamelle a dorso, armature microlitiche) nelle fasi tardo-pleistoceniche; la struttura sociale suggerisce gruppi di venti-quaranta individui con ampi range di spostamento nella fase fredda, tendenti a restringersi con il miglioramento climatico.

Per questo periodo, i siti di arte rupestre attesi in Veneto non si trovano in quota: le Alpi erano in gran parte inospitali o inaccessibili. I siti si concentrano nelle valli prealpine inferiori, nelle grotte e ripari tra cento e seicento metri, nella fascia di transizione tra pianura e montagna dove la selvaggina era disponibile e l'accesso relativamente sicuro. Riparo Tagliente, Riparo Mezzena, Grotta di Fumane (Lessinia) sono tutti in questa fascia; Riparo Dalmeri (Valsugana) a 1240 metri è un caso eccezionale legato a una fase di miglioramento climatico tardo-pleistocenico.

Il tipo di arte atteso per questo periodo è quello documentato a Dalmeri: pitture su supporto mobile (ciottoli dipinti), raramente arte parietale, e quando parietale, in cavità con pareti calcaree lavate o in ripari con fondo di accumulo sedimentario. Le grotte della Lessinia non risultano sistematicamente rilevate per arte parietale con metodi RTI; le valli prealpine della Valsugana e dell'Agordino non risultano oggetto di prospezioni dedicate per arte rupestre di questo periodo. I porfidi del Lagorai (quota inadatta per questo periodo) e le dolomie di vetta possono essere esclusi.

La metodologia appropriata per il Paleolitico è diversa da quella per i periodi successivi: non si cercano valichi e selle, ma cavità o ripari con pareti calcaree tra cento e seicento metri, con attenzione particolare alle superfici interne lavate dall'acqua, alle pareti verticali in penombra, e ai depositi sedimentari che possono coprire pannelli dipinti. RTI e fotogrammetria sono qui lo strumento primario di indagine delle superfici, non la prospezione a transetti.

% ============================================================
\section{Il Mesolitico: valichi, passi e logica di Kompatscher}
% ============================================================

Con il riscaldamento postgiaciale (da circa undicimilacinquecento a settemila anni prima dell'anno zero), la fauna cambia e con essa la strategia di caccia. Il cervo, lo stambecco e il camoscio diventano le specie dominanti alle quote montane; la renna scompare progressivamente dal territorio veneto. Le quote accessibili si alzano: intorno a ottomila anni prima dell'anno zero, le Alpi sopra i duemila metri tornano praticabili nella stagione estiva. La tecnologia è geometrica (trapezi, punte a dorso di piccole dimensioni); la struttura sociale è quella del gruppo di cacciatori-raccoglitori con mobilità stagionale, con siti di quota usati nella stagione calda e siti di valle nella stagione fredda.

Questo è il periodo a cui il metodo Kompatscher è più direttamente applicabile. I criteri di ricerca sono ben definiti: valichi e selle con almeno due direttrici di movimento verso zone di caccia, a quote tra millequattrocento e duemilatrecento metri (con il Castelnoviano che tende a preferire versanti esposti a sud per le temperature mattutine), con sorgente o torrentello a venti-ottanta metri, e morfologia riparata almeno su un lato. L'arte rupestre del Mesolitico, quando presente, si trova spesso in corrispondenza di queste strutture: non pitture elaborate, ma coppelle (bicchierini scolpiti), tratti incisi, impronte di mano, segni geomorfici su superfici calcaree subpianeggianti.

Nel Veneto, i candidati applicando questi criteri si distribuiscono lungo l'arco alpino a quote appropriate. I passi del gruppo dolomitico cadorino (Passo Duran, Passo Cereda, Passo di Giau, col dei passi sopra Alleghe) presentano morfologie adeguate e non risultano sistematicamente cercati per arte rupestre. Il Cansiglio, con le sue doline e i pianori carsici, rappresenta una variante morfologica del modello: non un valico, ma un bacino di confluenza di più percorsi con acqua abbondante (Pian Cansiglio ha sorgenti multiple). La ricerca di Ferrara ha documentato la frequentazione mesolitica in quest'area; l'arte rupestre specifica non è stata ancora il target principale.

% ============================================================
\section{Il Neolitico: coppelle, stele e pianori}
% ============================================================

Con il Neolitico (da circa settemila a cinquemila anni prima dell'anno zero) cambia non solo la tecnologia ma l'assetto insediativo. Le comunità agricole occupano la pianura e i fondovalle in modo più stabile; la montagna è frequentata stagionalmente per la pastorizia e la caccia, ma non con la stessa logica di mobilità mesolitica. L'arte rupestre neolitica ha caratteristiche proprie: le coppelle (cavità circolari o ovali incise per percussione sulla superficie rocciosa) diventano ubiquitarie e spesso si trovano su rocce subpianeggianti o debolmente inclinate, non sulle verticali preferite in altri periodi. Le stele menhir compaiono in questo periodo in vari contesti alpini (Trentino, Lombardia, Valle d'Aosta, Alto Adige) e vengono spesso riutilizzate in strutture medievali, dove possono essere recuperate come materiale di spoglio.

La metodologia per il Neolitico si divide quindi tra ricerca di coppelle e ricerca di stele. Le coppelle si cercano su pianori rocciosi calcarei a quote variabili (dalla pianura alle quote di mezza montagna), con attenzione ai punti panoramici e ai crocevia di sentieri storici; la densità attesa in contesti calcarei europei simili è di uno-cinque siti per dieci chilometri quadrati, ma questa è una stima molto rozza, dipendente dalla morfologia locale. Le stele si cercano in murature di edifici storici (chiese rurali, masi, fontane), nelle collezioni locali di materiale di reimpiego, e in archivi fotografici di lavori edilizi del Novecento. Il recupero di stele dalle murature è metodologicamente diverso dalla prospezione rupestre, ma può produrre risultati significativi: in Trentino, decine di stele sono state identificate proprio in questo modo.

In Veneto, le Prealpi calcaree (Lessinia, Piccole Dolomiti, Alpago) presentano morfologie idonee per le coppelle; i fondovalle storicamente abitati (Valpolicella, Vallagarina, Agordino) potrebbero nascondere stele in reimpiego. La ricerca specifica per questi target non risulta sistematicamente condotta.

% ============================================================
\section{L'Eneolitico: valichi come nodi di scambio}
% ============================================================

L'Eneolitico (da circa cinquemila a tremilacinquecento anni prima dell'anno zero) è il periodo delle prime metallurgie e degli scambi a lunga distanza. Ötzi, l'Uomo del Similaun, appartiene a questo periodo (3300-3100 anni prima dell'anno zero); la sua presenza su un passo alpino d'alta quota nel momento della morte illustra il ruolo che i valichi avevano nei movimenti di persone e merci. In questo periodo, le selle e i valichi diventano nodi di un sistema di scambio di rame, ossidiana, ambra e prodotti agricoli; l'arte rupestre eneolitica tende a concentrarsi in corrispondenza di questi nodi, come documentato in Val Camonica (incisioni sui bordi dei valichi laterali) e nel Tirolo meridionale.

Per il Veneto, i valichi delle Piccole Dolomiti (Passo Borcola, Passo Campogrosso), del Pasubio (Bocca di Campiglia, passo Xomo), del Brenta-Adamello (Passo Tremalzo) e delle Dolomiti bellunesi e cadorine (Passo Duran, Passo Cereda, Forcella Valbona) sono candidati naturali per arte rupestre eneolitica. La morfologia eneolitica preferita è diversa da quella mesolitica: non il campo base riparato di un cacciatore, ma la parete verticale o inclinata in corrispondenza di un passaggio, dove si incidono segni visibili da chi transita. Le superfici calcaree lisce, esposte in verticale o con inclinazione di quaranta-settanta gradi, nel raggio di cinquanta metri da un punto di passaggio obbligato, sono il target specifico per questo periodo.

% ============================================================
\section{L'età del bronzo: pareti verticali e bacini visivi}
% ============================================================

L'età del bronzo (da circa tremilacinquecento a duemilanovecento anni prima dell'anno zero) è il periodo di maggiore intensità documentata per l'arte rupestre in area alpina. Valcamonica e Valtellina hanno migliaia di figure di questo periodo; la Val d'Assa veneta è in larga parte bronzina; Monte Luppia (Garda) ha pannelli riferibili a questo periodo. La firma stilistica è riconoscibile: figure geometriche complesse (retticolati, spirali, cerchi concentrici), figure umane stilizzate (oranti con braccia alzate), armi (asce, pugnali, lance), animali (bovini, cervi) su superfici verticali o leggermente inclinate con bacino visivo ampio.

Il target morfologico per l'età del bronzo è specifico: superficie calcarea verticale o inclinata tra i trenta e i settanta gradi, con overhang (sporgenza soprastante) che protegga dall'erosione pluviale, a quote tra quattrocento e milleduecento metri, con ampia visibilità su almeno una vallata principale. La combinazione di inclinazione, protezione e quota restringe i candidati su qualsiasi cartografia topografica. La metodologia di rilevamento appropriata è la combinazione di fotogrammetria SfM (per il modello 3D della superficie), luce radente artificiale nelle ore notturne (per rivelare le incisioni più superficiali), e RTI su campioni selezionati.

In Veneto, le aree candidate per arte rupestre dell'età del bronzo non ancora sistematicamente cercate includono la Lessinia settentrionale (zona Sant'Anna d'Alfaedo, Vajo dell'Anguilla, pareti calcaree al limite del bosco tra seicento e novecento metri), l'Altopiano di Asiago nelle zone non interessate dai combattimenti della Prima guerra mondiale (parte meridionale, verso Roana e Foza), e il versante meridionale del Grappa verso Romano d'Ezzelino e Pove del Grappa. Tutte e tre le aree hanno morfologie calcaree con pareti verticali nella fascia di quota appropriata; nessuna delle tre risulta oggetto di prospezioni sistematiche dedicate all'arte rupestre.

% ============================================================
\section{L'età del ferro e il periodo venetico: continuità e aggiunte}
% ============================================================

Nell'età del ferro (da circa duemilanovecento a duemila anni prima dell'anno zero) e nel periodo venetico (fino alla romanizzazione), l'arte rupestre non cessa ma spesso si sovrappone ai pannelli bronzini già esistenti. Le aggiunte sono riconoscibili per stile (figure più schematiche, iscrizioni, simboli geometrici elementari) e tecnica (incisione più fine rispetto alla martellina del Bronzo). In Val Camonica, le aggiunte di età del ferro e periodo celtico o romano sono documentate su decine di pannelli già utilizzati in precedenza: il luogo continua ad essere frequentato e segnato, ma con un'intensità e un vocabolario diversi.

La metodologia per identificare arte rupestre di questo periodo nel Veneto passa quindi, principalmente, per la rilettura dei siti già noti. Monte Luppia e Val d'Assa sono i candidati ovvi per trovare aggiunte di età del ferro non ancora distinte stilisticamente dal bronzino circostante: una campagna RTI sistematica su pannelli già catalogati potrebbe rivelare stratificazioni non visibili a occhio nudo. Il periodo venetico, specifico del territorio veneto, ha una produzione epigrafica nota (stele funerarie, iscrizioni votive su bronzi), ma non risultano documentate iscrizioni venetiche su supporto rupestre: la loro esistenza è possibile ma non dimostrata.

% ============================================================
\section{I luoghi persistenti: quattro tipi di attrattore}
% ============================================================

L'analisi periodo per periodo rivela un dato trasversale: alcuni tipi di luogo ricorrono come attrattori di attività umana indipendentemente dall'epoca. Questi "luoghi persistenti" non sono una teoria astratta, ma un pattern documentato empiricamente in Valle Camonica, in Trentino-Alto Adige, in Scandinavia, e nelle Alpi svizzere e tirolesi. I siti rupestri tendono a concentrarsi in luoghi che soddisfano contemporaneamente criteri di più periodi, e questa sovrapposizione è di per sé un indizio di ricerca.

Il primo tipo di attrattore persistente è la sella con acqua e direttrici multiple: un passo o una forcella che mette in comunicazione due o più bacini, con sorgente a meno di ottanta metri, e da cui partono almeno due percorsi in direzioni diverse. Questo tipo di luogo è frequentato dal Mesolitico all'età del ferro senza discontinuità apparente. In Veneto, i passi delle Piccole Dolomiti (Borcola, Campogrosso) e alcuni valichi cadorini hanno questa morfologia; vanno esaminati prioritariamente.

Il secondo tipo è la parete calcarea verticale con overhang a quota intermedia: una superficie rocciosa protetta, tra quattrocento e milleduecento metri, con visibilità su una vallata. Questo tipo di luogo è prevalentemente bronzino, ma può avere aggiunte eneolitiche e dell'età del ferro. La Lessinia settentrionale ha decine di pareti di questo tipo; nessuna è stata sistematicamente cercata.

Il terzo tipo è la superficie calcarea orizzontale o subpianeggiante vicina a sentieri di quota: un pianoro o una cengia con roccia liscia, percorsa da un sentiero di mezza montagna, a quote tra ottocento e millecinquecento metri. Questo tipo di luogo ospita prevalentemente coppelle (Neolitico e Bronzo antico), ma può avere segni di tutti i periodi. I pianori carsici della Lessinia meridionale (Malga San Giorgio, altopiano di Campofontana) hanno morfologie di questo tipo.

Il quarto tipo è il crinale panoramico su due o più versanti, con morfologia a lama o a cresta arrotondata, a quote tra milleduecento e duemila metri. Questo tipo di luogo ha funzione prevalente di avvistamento e probabilmente non ospita arte rupestre in senso stretto (la roccia esposta ai venti di cresta si erode rapidamente), ma è un nodo logistico che connette altri siti. La sua importanza è metodologica: permette di identificare i percorsi tra siti, non i siti stessi.



% Cap. 10 — Bilancio dell'ignoto: sintesi tafonomia + metodologia,
% due tavole (tipi di arte × epoca; tipi di luogo × arte × epoche),
% santuari alpini (Ceruti, Reinhard, Tresero), lacune osservative
% su larga scala, conclusioni aperte come sezione finale.
\chapter{Un bilancio dell'ignoto}
\label{cap:bilancio}

I capitoli precedenti hanno costruito tre strumenti distinti per avvicinarsi alla domanda centrale di questo libro. Il primo è tafonomico: una mappa di ciò che il tempo, il clima e le azioni umane cancellano o conservano in modo selettivo e prevedibile. Il secondo è critico: un'analisi delle lacune della ricerca, delle zone bianche che non riflettono l'assenza di tracce ma l'assenza di chi le cercasse. Il terzo è metodologico: i criteri che, applicati alla morfologia del territorio, identificano i tipi di luoghi più promettenti per ciascun periodo preistorico. Questo capitolo li mette insieme in due tavole sintetiche, analizza cosa dicono quando letti in parallelo, e poi indica le lacune osservative su larga scala che nessuna delle tre analisi precedenti ha ancora esaurito.

% ============================================================
\section{Come leggere le due tavole}
% ============================================================

Le tavole che seguono usano tre simboli. Il simbolo $\bullet$ indica presenza documentata: l'arte rupestre o la costruzione di quel tipo è attestata in Veneto o in un'area immediatamente comparabile per geografia, substrato e periodo. Il simbolo $\circ$ indica presenza plausibile, non documentata in Veneto ma attesa per analogia con contesti alpini simili in cui è stata cercata con metodi adeguati. Il trattino $-$ indica assenza o improbabilità nelle condizioni del Veneto, per ragioni tafonomiche, climatiche o legate alle forme di occupazione del territorio nell'epoca considerata.

La colonna \emph{costruzioni} riunisce categorie materialmente eterogenee: megaliti e stele menhir nel Neolitico e nell'Eneolitico; dossi artificiali e terrapieni nelle pianure dell'età del Bronzo; castellieri e strutture difensive sulle creste dell'età del Bronzo e del Ferro; strutture con funzioni rituali o calendariali di cui restano tracce indirette. Le accomuna un'unica caratteristica: sono modificazioni permanenti e intenzionali del paesaggio fisico, distinte dall'arte rupestre in senso stretto ma espressione dello stesso impulso di marcare il territorio in modo durevole. In Veneto alcune di queste categorie sono documentate; altre sono quasi assenti dalla letteratura non perché manchino, ma perché non sono mai state sistematicamente cercate.

La seconda tavola include le epoche prevalenti per ciascun tipo di luogo, poiché la pertinenza di un sito non dipende solo dalla sua morfologia ma dal periodo in cui era frequentato con l'intensità sufficiente a produrre arte o strutture permanenti. Un valico a duemila metri è un candidato plausibile per il Mesolitico e per l'Eneolitico; lo è molto meno per il Paleolitico superiore, quando quella quota era ghiacciata o periglaciale.

% ============================================================
\section{Tavola 1: quali tipi di arte, in quale epoca}
% ============================================================

\begin{table}[h]
\small
\setlength{\extrarowheight}{3pt}
\caption{Presenza attesa per tipo di arte e per epoca nel territorio veneto e nelle aree immediatamente comparabili. Date in migliaia di anni prima dell'anno zero astronomico. $\bullet$~=~documentato;\enspace $\circ$~=~plausibile per analogia;\enspace $-$~=~assente o molto improbabile.}
\label{tab:bilancio_tipi}
\begin{tabularx}{\textwidth}{>{\raggedright\arraybackslash}X
                              >{\centering\arraybackslash}p{2.0cm}
                              >{\centering\arraybackslash}p{2.0cm}
                              >{\centering\arraybackslash}p{2.4cm}}
\toprule
\textbf{Epoca} & \textbf{Incisioni} & \textbf{Pitture} & \textbf{Costruzioni} \\
\midrule
\rowcolor{blue!8}
Paleolitico superiore (38.000--11.500 a.p.a.z.) & $\circ$ & $\bullet$ & $-$ \\
Mesolitico (11.500--7.500 a.p.a.z.) & $\circ$ & $\circ$ & $-$ \\
\rowcolor{blue!8}
Neolitico (7.500--5.250 a.p.a.z.) & $\circ$ & $-$ & $\circ$ \\
Eneolitico (5.250--4.250 a.p.a.z.) & $\circ$ & $-$ & $\bullet$ \\
\rowcolor{blue!8}
Età del Bronzo (4.250--2.750 a.p.a.z.) & $\bullet$ & $-$ & $\bullet$ \\
Età del Ferro e periodo venetico (2.750--0 a.p.a.z.) & $\bullet$ & $-$ & $\bullet$ \\
\bottomrule
\end{tabularx}
\end{table}

% ============================================================
\section{Tavola 2: quali luoghi, quali tipi di arte, quali epoche}
% ============================================================

\small
\setlength{\extrarowheight}{3pt}
\begin{longtable}{>{\raggedright\arraybackslash}p{3.0cm}
                   >{\centering\arraybackslash}p{1.5cm}
                   >{\centering\arraybackslash}p{1.5cm}
                   >{\centering\arraybackslash}p{1.7cm}
                   >{\raggedright\arraybackslash}p{2.8cm}}

\caption{Presenza attesa per tipo di luogo e tipo di arte, con le epoche prevalenti. I giudizi tengono conto della tafonomia, del clima e delle forme di occupazione del territorio per ciascun periodo. Stessi simboli della tavola precedente.\label{tab:bilancio_luoghi}} \\

\toprule
\textbf{Tipo di luogo} & \textbf{Incis.} & \textbf{Pitt.} & \textbf{Costr.} & \textbf{Epoche prevalenti} \\
\midrule
\endfirsthead

\multicolumn{5}{l}{\small\textit{Tavola~\ref{tab:bilancio_luoghi} -- continua dalla pagina precedente}} \\[4pt]
\toprule
\textbf{Tipo di luogo} & \textbf{Incis.} & \textbf{Pitt.} & \textbf{Costr.} & \textbf{Epoche prevalenti} \\
\midrule
\endhead

\midrule
\multicolumn{5}{r}{\small\textit{continua nella pagina seguente}} \\
\endfoot

\bottomrule
\endlastfoot

\rowcolor{blue!8}
Grotte (interni) &
  $\circ$ & $\bullet$ & $-$ &
  Paleolitico, Mesolitico \\[2pt]

Ripari rocciosi semi-aperti &
  $\circ$ & $\bullet$ & $-$ &
  Paleolitico, Mesolitico \\[2pt]

\rowcolor{blue!8}
Punti di avvistamento (speroni e rilievi con visibilità multipla) &
  $\circ$ & $-$ & $\circ$ &
  Mesolitico, Bronzo \\[2pt]

Selle e valichi &
  $\circ$ & $-$ & $\circ$ &
  Mesolitico, Eneolitico, Bronzo \\[2pt]

\rowcolor{blue!8}
Pareti verticali a quota intermedia (400--1.200 m, con overhang) &
  $\bullet$ & $-$ & $-$ &
  Bronzo, Ferro \\[2pt]

Radure d'alta quota non abitate &
  $\circ$ & $-$ & $\circ$ &
  Neolitico, Bronzo \\[2pt]

\rowcolor{blue!8}
Falde dei gruppi montagnosi (zona di transizione bosco-quota) &
  $\circ$ & $-$ & $\circ$ &
  Mesolitico, Neolitico, Bronzo \\[2pt]

Crode dolomitiche d'alta quota (oltre 1.800 m) &
  $\circ$ & $-$ & $\circ$ &
  Bronzo, Ferro (uso cultuale) \\[2pt]

\rowcolor{blue!8}
Fondovalle montani (Val Belluna, Zoldano, Agordino, Piave alta) &
  $\circ$ & $-$ & $\circ$ &
  Neolitico, Bronzo, Ferro \\[2pt]

Colli isolati di pianura (Euganei, Berici) &
  $\circ$ & $-$ & $\circ$ &
  Paleolitico, Mesolitico, Neolitico \\[2pt]

\rowcolor{blue!8}
Dossi e terrazze fluviali di pianura &
  $-$ & $-$ & $\bullet$ &
  Neolitico, Bronzo, Ferro \\[2pt]

\end{longtable}
\normalsize

% ============================================================
\section{Cosa dicono le tavole}
% ============================================================

Lette insieme, le due tavole producono alcune osservazioni che non emergono separatamente dall'analisi tafonomica o da quella metodologica.

La prima: le pitture sono il tipo di arte rupestre più fragile all'aperto e al tempo stesso il meglio documentato in Veneto per il Paleolitico superiore. Fumane e Dalmeri sono pitture, non incisioni. Non è casuale: le pitture su parete interna di grotta, in penombra e protette dall'umidità costante, sopravvivono meglio di qualsiasi altra categoria in ambiente calcareo. La loro assenza dal mondo esterno è quasi totale per ragioni tafonomiche prevedibili; la loro presenza nelle grotte della Lessinia e del Veneto orientale è sottodocumentata perché non è mai stata sistematicamente cercata con metodi ottici avanzati.

La seconda: le incisioni all'aperto sono documentate solo per il Bronzo e il Ferro, ma plausibili per quasi tutte le epoche precedenti. Il divario non è nella produzione originaria, ma nella sopravvivenza. Un'incisione su calcare esposto alle precipitazioni perde leggibilità in poche centinaia di anni; un'incisione su calcare protetto da un overhang può durare millenni. La differenza tra un $\bullet$ e un $\circ$ nella colonna delle incisioni riflette quasi sempre la presenza o l'assenza di protezione morfologica, non la presenza o l'assenza dell'attività umana.

La terza: le costruzioni sono il tipo di traccia più resistente al tempo quando sono in pietra o in terra compattata, ma la più invisibile quando erano in legno o in materiale organico. I castellieri dell'età del Bronzo sopravvivono perché erano strutture in pietra su creste già rocciose. I probabili punti di sosta mesolitici non lasciano nulla di visibile in superficie. Il divario tra i $\bullet$ delle costruzioni documentate e i $\circ$ di quelle plausibili nelle epoche più antiche non misura l'intensità dell'occupazione umana ma la sua modalità: stabile e sedentaria contro mobile e stagionale.

La quarta, e forse la più importante per il ragionamento di questo libro: le righe con più $\circ$ nella seconda tavola non sono le meno promettenti. Sono le meno cercate. Un $\circ$ significa che l'analogia con aree comparabili dice che ci dovrebbe essere qualcosa, e che la ragione per cui non è documentato è che nessuno si è fermato a guardare con metodi adeguati.

% ============================================================
\section{Le lacune osservative su larga scala}
% ============================================================

Il capitolo sul research bias ha già analizzato le lacune legate ai siti noti: i meccanismi che producono mappe della conoscenza disomogenee, dove la presenza di un museo, di un naturalista ottocentesco o di una cava produce concentrazioni di segnalazioni che nulla hanno a che fare con la distribuzione reale dell'arte preistorica. Quello che rimane fuori da quella analisi sono le lacune geografiche strutturali: intere aree del territorio veneto e delle zone contigue che non figurano in alcun programma di ricerca sistematica per arte rupestre preistorica. Non si tratta di singole pareti o di piccole creste, ma di masse montuose, fasce altitudinali e bacini vallivi di scala tale da costituire, nella loro assenza dalla letteratura, un vuoto obiettivo.

\subsection*{Le crode del Cadore e del Comelico}

I principali gruppi dolomitici del Cadore e del Comelico, la Marmolada, le Pale di San Martino, le Pale di San Lucano, la Civetta, l'Antelao, il Pelmo, l'Averau, il Nuvolau, le Tofane, il Sorapiss, le Marmarole, il Monte Schiara, le Vette Feltrine e tutti i gruppi minori intermedi, non risultano oggetto di alcuna prospezione sistematica per arte rupestre preistorica. I pochi segnali disponibili confermano che la ricerca c'è stata appena abbozzata: Fabio Cavulli ha presentato nel 2025 a Pieve di Cadore i risultati preliminari di cerchi incisi al compasso su superfici dolomitiche, un indizio che queste pareti potrebbero riservare sorprese.

Ma c'è un secondo motivo per non escludere le alte quote, un motivo che va oltre la caccia e i campi base: il culto delle montagne. Constanza Ceruti, l'archeologa argentina che ha ritrovato mummie sacrificali sulle cime andine a oltre sei mila metri di quota, ha documentato in un recente lavoro sull'arco alpino come la sacralità delle cime alpine, comprese le Dolomiti italo-austriache, sia rimasta largamente ignorata dalla discussione accademica contemporanea \citep{ceruti2025}. Il riferimento andino non è decorativo: Ceruti e Reinhard hanno dimostrato che il culto delle montagne è documentato in modo indipendente su quattro continenti, dalle Ande all'Himalaya, dalle Alpi alle catene messicane \citep{reinhard2010}, il che suggerisce che risponda a qualcosa di radicato nell'esperienza umana del paesaggio d'alta quota. Non sarebbe la prima volta che le Alpi conservano tracce di questa relazione: nel 2024, quando il ghiacciaio del Tresero in Alta Valtellina si è ritirato abbastanza da esporre la roccia sottostante, sono comparse le incisioni rupestri più alte d'Europa, a tremila metri di quota, probabilmente parte di un santuario rupestre dell'età del Bronzo che il ghiaccio aveva coperto per millenni \citep{tresero2024}. Il ritiro glaciale accelerato degli ultimi decenni sta rendendo visibili superfici che erano inaccessibili durante tutta la storia della ricerca moderna. Le Dolomiti venete hanno avuto ghiacciai propri.

\subsection*{Le falde dei gruppi montagnosi}

I versanti inferiori e le fasce di transizione tra bosco e quota dei principali massicci veneti, Baldo, Carega, Pasubio, Recoaro, Tonezza, le zone dell'Altopiano di Asiago non interessate dai combattimenti e il Grappa, sono aree frequentate dalla vita umana moderna ma non sistematicamente cercate per arte preistorica. Esistono già ritrovamenti in alcune di queste zone: il Baldo, Asiago e Folgaria hanno restituito materiali che attestano frequentazione preistorica. Questi ritrovamenti non sono la conclusione dell'indagine ma il punto di partenza: dimostrano presenza umana in queste fasce altitudinali, il che rende plausibile la presenza di arte rupestre nei tipi di luogo descritti nel capitolo precedente. Le falde a quote tra i quattrocento e i milleduecento metri sono precisamente la fascia in cui il capitolo sulla metodologia individua le pareti verticali con overhang dell'età del Bronzo e i ripari semi-aperti del Mesolitico.

\subsection*{I fondovalle montani non indagati}

La Val Belluna, la valle del Piave nella sua sezione alta, lo Zoldano e l'Agordino sono pressoché assenti dalla letteratura sull'arte rupestre preistorica veneta. Questi fondovalle non sono ambienti sfavorevoli: sono corridoi naturali di transito tra la pianura e le quote alpine, con substrati calcarei nelle pareti laterali delle valli e microclimi stabili che avrebbero favorito sia la frequentazione sia la conservazione delle tracce. La loro assenza dalla letteratura è una lacuna di ricerca, non un risultato.

\subsection*{Le radure non abitate}

Le radure d'alta quota sono ambienti trasformati dall'azione umana, ma non necessariamente da un'azione recente. La maggior parte dei pianori erbosi delle Prealpi venete è stata disboscata a partire dal Medioevo, e le strutture pastorali visibili in superficie, malghe, tabià, muretti a secco, datano agli ultimi cinque-otto secoli. Sotto queste strutture e nei loro margini possono esistere tracce di frequentazione molto più antica: coppelle su pianori rocciosi calcarei, superfici sub-pianeggianti con incisioni, depositi di materiale in anfratti protetti. Le radure attualmente in uso per la pastorizia sono difficili da esplorare sistematicamente; quelle abbandonate, dove la vegetazione ha ripreso il sopravvento dopo la cessazione dell'uso, sono quasi totalmente ignorate dalla ricerca. Sono anche le meno disturbate dalle attività recenti.

\subsection*{I Colli Euganei}

I Colli Euganei hanno una geologia anomala rispetto al resto del Veneto: sono un massiccio di rocce magmatiche intrusive ed effusive, trachiti, fonoliti e rioliti di età eocenica, che emergono isolati dalla pianura fino a 601 metri. Come già discusso nel capitolo sulla storia del territorio, questi colli hanno funzionato per tutta la preistoria come rifugi ecologici con microclima mite, acqua affidabile e substrato duro favorevole alla scheggiatura. Ciò che non è mai stato sistematicamente cercato è l'arte rupestre su roccia vulcanica. Gli affioramenti trachitici hanno superfici dure e resistenti, molto più favorevoli all'incisione duratura rispetto al calcare. La presenza di minerali cupriferi e ferrosi nei Colli Euganei rende geologicamente plausibile che comunità eneolitiche e dell'età del Bronzo abbiano riconosciuto e frequentato questi colli non solo come rifugio ma come fonte di materie prime: e dove c'era estrazione e frequentazione intensiva, c'era anche la condizione che altrove ha prodotto arte rupestre. Non esistono ancora prove, ma esistono i presupposti perché valga la pena cercare.

% ============================================================
\section{Quello che potremmo trovare}
% ============================================================

Mettendo insieme la tafonomia, le lacune della ricerca e la metodologia di prospezione, si può dire qualcosa di ragionevole sulla natura di ciò che potrebbe emergere da una ricerca sistematica del territorio veneto, senza cedere né all'ottimismo facile né al pessimismo ingiustificato.

I ripari e le grotte sono il tipo di contesto dove la probabilità di trovare qualcosa è concretamente più alta. Le pitture del Paleolitico superiore nelle grotte della Lessinia e dei gruppi prealpini non sono mai state cercate con metodi ottici avanzati. I ripari rocciosi delle falde prealpine, tra i quattrocento e gli ottocento metri, potrebbero conservare incisioni su pareti protette dall'overhang che la vegetazione e la patina superficiale rendono invisibili alla ricognizione ordinaria. Questi contesti hanno la caratteristica tafonomica più favorevole: sono chiusi o semi-chiusi, protetti dalle precipitazioni, e hanno subito variazioni di temperatura e umidità meno estreme di qualsiasi superficie all'aperto.

L'arte rupestre su pareti aperte è meno probabile, per le ragioni discusse nel capitolo sulla tafonomia: il calcare delle Prealpi venete, esposto alle precipitazioni e ai cicli di gelo e disgelo per diecimila anni, ha perso la maggior parte delle superfici che in origine potrebbero aver portato incisioni superficiali. L'eccezione è rappresentata dalle pareti con overhang pronunciato, dove la protezione è stata sufficiente a rallentare l'erosione: e queste pareti sono precisamente il target metodologico del capitolo precedente.

Le crode d'alta quota restano l'incognita più aperta, e forse la più affascinante. Non per la probabilità statistica di un ritrovamento, che è bassa, ma perché un ritrovamento in quel contesto avrebbe un significato diverso da tutti gli altri: non un sito di caccia, non un campo base stagionale, ma qualcosa che richiederebbe una spiegazione diversa. Il ghiacciaio del Tresero ha mostrato che questa categoria esiste nelle Alpi italiane. Il ritiro glaciale in corso sta aprendo superfici che nessun ricercatore ha mai potuto vedere. Le Dolomiti venete hanno ghiacciai che si ritirano. Non sappiamo cosa nascondono.

Basta trovarne uno. Un riparo con pitture mai documentate, una grotta con incisioni che nessuno aveva cercato, un segno su una croda a quota elevata che non ha spiegazioni immediate. Cambierebbe tutto ciò che ora riteniamo di sapere sul silenzio di queste montagne.


\backmatter

% Appendice: Repertorio dei siti preistorici veneti nel catalogo MiC
% Abbreviazione usata in tutto il capitolo:
%   a.p.z.a. = anni prima dell'anno zero astronomico

\chapter{Repertorio dei siti preistorici veneti nel catalogo MiC}
\label{cap:appendice_siti}

\aprecapitolo{I}

\section*{Nota introduttiva}

I dati raccolti in questa appendice provengono interamente dal catalogo ufficiale del Ministero della Cultura italiano, accessibile al pubblico attraverso il portale \textit{catalogo.beniculturali.it} gestito dall'Istituto Centrale per il Catalogo e la Documentazione (ICCD). Le schede sono state consultate e rielaborate nell'estate del 2026 ed erano quelle disponibili a quella data.

Il catalogo MiC non è un inventario sistematico di tutto ciò che esiste nel sottosuolo veneto: è il registro di ciò che è stato formalmente catalogato, cioè scavato, segnalato e inserito nel sistema nazionale da parte di soprintendenze, università, musei civici e privati cittadini con le debite autorizzazioni. La sua incompletezza è dunque strutturale, ma anche istruttiva: le zone con densità alta di schede corrispondono ad aree esplorate; le zone vuote ci dicono dove la ricerca non ha ancora guardato.

Abbiamo incluso tutti i 136 siti con almeno una fase databile a età preistorica (dal Paleolitico all'età del Ferro), escludendo i siti esclusivamente romani o medievali (una trentina di schede). I siti sono ordinati per epoca, dalla più antica alla più recente, e all'interno di ciascuna epoca per zona geografica. Per i siti che attraversano più epoche, la descrizione completa si trova alla prima occorrenza; le sezioni successive riportano solo un rimando. Per le date numeriche usiamo l'abbreviazione \textbf{a.p.z.a.} (anni prima dell'anno zero astronomico), equivalente a ciò che la letteratura scientifica indica comunemente con ``a.C.'', ma senza l'ambiguità religiosa di quella formula.

% ============================================================
\section{Paleolitico}
\label{app:paleolitico}
% 13 siti
% ============================================================

Il Paleolitico copre un intervallo enorme, dai primi manufatti musteriani attribuibili all'uomo di Neanderthal (circa 300.000 a.p.z.a.) fino all'Epigravettiano finale (circa 10.000 a.p.z.a.), quando i cacciatori-raccoglitori del Paleolitico superiore si adattano al progressivo scioglimento dei ghiacciai. Nel Veneto questo periodo è documentato principalmente nella Lessinia veronese e, in misura minore, nell'Altopiano di Asiago e in alcune grotte dei Colli Berici.

\subsection{Altopiano Vicentino}

\textbf{Roana (VI).} Una grotta nel territorio di Roana ha restituito un deposito litico musteriano, senza ulteriori specificazioni nel catalogo. La scheda è essenziale: il sito è noto ma non scavato sistematicamente.

\textbf{Asiago (VI).} Un riparo sotto roccia in territorio di Asiago ha fornito evidenze di frequentazione epigravettiana finale, interpretate come sosta stagionale di cacciatori nell'area altopianare, probabilmente legata ai movimenti della fauna in quota.

\subsection{Colli Euganei e Berici}

\textbf{Longare (VI).} La grotta nel territorio di Longare mostra una sequenza che abbraccia Musteriano, Aurignaziano e Gravettiano, suggerendo frequentazioni ripetute in un arco temporale molto lungo. L'interpretazione MiC la inquadra come rifugio stagionale; si tratta di uno dei pochi siti dei Berici con attestazioni paleolitiche multiple.

\textbf{Cinto Euganeo (PD).} Nei pressi di Valnogaredo, i ritrovamenti superficiali sono interpretati come tracce di frequentazione musteriana. Non sono state individuate strutture; la localizzazione ai margini dei Colli Euganei è geograficamente interessante per comprendere le rotte di caccia in area collinare.

\subsection{Lessinia}

\textbf{Dolcè (VR).} Un riparo sotto roccia nel fondovalle dell'Adige, scavato tra il 1983 e il 1988, conserva una sequenza che va dal Paleolitico superiore al Mesolitico (culture Sauveterriana e Castelnoviana). Nel 1983 fu segnalata anche una sepoltura a inumazione di probabile età olocenica. Il sito è quindi multifase: vedi anche la sezione Mesolitico~(\ref{app:mesolitico}).

\textbf{Fumane (VR) -- Grotta di Fumane.} \label{app:fumane_paleo} Uno dei siti paleolitici più importanti d'Europa, secondo la valutazione stessa del catalogo MiC. La grotta, alla base della Valle dei Progni, fu intravista nel tardo Ottocento durante lavori stradali e saccheggiata ripetutamente prima che, a partire dal 1988, un gruppo di ricerca articolato tra Soprintendenza, Università di Ferrara, Università di Milano e Museo di Scienze Naturali di Verona avviasse scavi sistematici. La sequenza stratigrafica comprende livelli riferibili al Paleolitico medio (Musteriano), al Paleolitico superiore iniziale (Uluzziano, poi Aurignaziano) e al Gravettiano. La presenza di manufatti su osso, ornamenti personali e possibili tracce di comportamento simbolico colloca questa grotta tra i contesti di riferimento per la transizione tra Neanderthal e uomo anatomicamente moderno in Italia nord-orientale. Per la fase neolitica ed eneolitica vedi anche la scheda nella sezione Neolitico~(\ref{app:fumane_neo}).

\textbf{Grezzana (VR).} Un sito classificato come ``uno dei più importanti d'Italia'' per ampiezza stratigrafica, quantità di reperti e qualità della documentazione, con sequenze che vanno dal Paleolitico medio musteriano all'Aurignaziano e all'Epigravettiano. Non si hanno indicazioni di apertura turistica o di scavi recenti nel catalogo.

\textbf{Sant'Anna d'Alfaedo (VR).} Scavi condotti a partire dal 1922, sia nella caverna che nelle aree circostanti, hanno documentato una continuità di presenza umana dal Paleolitico inferiore al Neolitico. Il catalogo segnala come particolarmente sviluppata la lavorazione della selce, attiva sino all'età del rame. Vedi la sezione Neolitico~(\ref{app:alfaedo_neo}).

\textbf{Sant'Ambrogio di Valpolicella (VR).} Tracce di frequentazione musteriana, senza ulteriori dettagli nella scheda. La posizione all'imbocco della Valpolicella suggerisce un possibile punto di caccia o di transito lungo il bordo meridionale della Lessinia.

\subsection{Pianura Veronese e Bassa}

\textbf{Roncà (VR).} \label{app:ronca_paleo} Sito pluristratificato con fase musteriana, seguita da fasi eneolitica, bronzo recente-finale e tarda età del ferro, con tracce di romanizzazione. La scheda è sintetica (``sito pluristratificato'') senza descrizione delle strutture. Vedi anche le sezioni Bronzo~(\ref{app:ronca_bronzo}) e Ferro~(\ref{app:ronca_ferro}).

\subsection{Pianura e Delta}

\textbf{Vo (PD).} Frequentazione musteriana individuata nel Padovano. Tracce di superficie senza strutture accertate; la scheda interpreta il rinvenimento come frequentazione stagionale dell'area da parte di gruppi neanderthaliani.

\subsection{Prealpi Trevigiane e Pedemontana}

\textbf{Tarzo (TV).} \label{app:tarzo_paleo} Sito pluristratificato con tracce dell'uomo di Neanderthal (Paleolitico medio), frequentazione mesolitica castelnoviana e tracce di età del Bronzo. Il sito si trova nel versante meridionale delle Prealpi orientali; la successione di popolamenti ne fa un punto di riferimento per la preistoria trevigiana. Vedi anche le sezioni Mesolitico~(\ref{app:mesolitico}) e Bronzo~(\ref{app:tarzo_bronzo}).

\subsection*{Nota di sintesi -- Paleolitico}

I tredici siti paleolitici veneti conosciuti sono distribuiti in modo tutt'altro che casuale: la Lessinia veronese concentra i siti più ricchi e meglio documentati, grazie alla disponibilità di grotte e ripari in una roccia calcarea abbondante di selce di buona qualità. La pianura padana e il delta sono quasi assenti, non perché fossero disabitate ma perché migliaia di anni di depositi alluvionali ricoprono i livelli paleolitici a profondità difficilmente raggiungibili dagli scavi ordinari. Le Prealpi trevigiane e l'Altopiano di Asiago mostrano frequentazioni marginali, probabilmente legate a percorsi di caccia stagionale verso le quote medie. Il Bellunese e il Cadore sono del tutto assenti da questa lista, fatto che riflette la scarsità di prospezioni sistematiche in quota, non necessariamente l'assenza di occupazione umana.

% ============================================================
\section{Mesolitico}
\label{app:mesolitico}
% 3 siti
% ============================================================

Il Mesolitico (grossomodo tra 10.000 e 6.000 a.p.z.a.) è il periodo di adattamento alle condizioni postglaciali: i boschi avanzano, la fauna cambia, le strategie di sussistenza si moltiplicano. Nel Veneto il Mesolitico è documentato quasi esclusivamente in zone di transizione tra pianura e montagna, lungo fiumi e laghi.

\subsection{Lessinia}

\textbf{Dolcè (VR).} Vedi la descrizione completa alla sezione Paleolitico~(\ref{app:paleolitico}). La fase mesolitica comprende culture Sauveterriana e Castelnoviana, indicando una continuità di occupazione del riparo dal Paleolitico superiore attraverso l'intero Mesolitico.

\subsection{Prealpi Trevigiane e Pedemontana}

\textbf{Tarzo (TV).} Vedi la descrizione completa alla sezione Paleolitico~(\ref{app:tarzo_paleo}). La componente mesolitica è riferibile alla cultura Castelnoviana, ben attestata nel versante meridionale delle Prealpi orientali.

\subsection{Prealpi Vicentine e Piccole Dolomiti}

\textbf{Valdagno (VI).} Un riparo sotto roccia nel territorio di Valdagno ha restituito evidenze di frequentazione mesolitica. I saggi di scavo non hanno individuato strutture permanenti; si ipotizza un utilizzo stagionale del riparo, probabilmente connesso a caccia o pascolo nelle valli prealpine vicentine.

\subsection*{Nota di sintesi -- Mesolitico}

Tre sole schede per il Mesolitico veneto sono un numero che riflette soprattutto la storia della ricerca: il Mesolitico è difficile da identificare in superficie, i suoi siti sono spesso piccoli e con scarsa densità di materiali, e le prospezioni sistematiche in quota sono rare. I siti noti si concentrano lungo il bordo prealpino e nelle valli di accesso alla montagna, coerentemente con le strategie di mobilità stagionale dei cacciatori mesolitici. La rete di siti mesolitici documentata nel Trentino e nell'Alto Adige contigui (Riparo Dalmeri, Mondeval de Sora, piani di Eraclea) suggerisce che anche il Veneto montano fosse frequentato durante questo periodo, ma le tracce non sono ancora state trovate o catalogate.

% ============================================================
\section{Neolitico}
\label{app:neolitico}
% 22 siti
% ============================================================

Il Neolitico veneto (circa 6.000-3.500 a.p.z.a.) è caratterizzato dall'arrivo dell'agricoltura e dell'allevamento, con comunità che si stabilizzano in insediamenti più duraturi. La cultura dei ``Vasi a Bocca Quadrata'' è il marcatore culturale dominante nel Neolitico medio-recente della pianura padano-veneta.

\subsection{Altopiano Vicentino}

\textbf{Lusiana (VI).} Insediamento d'altura con elementi strutturali legati a un sistema di fortificazione, classificato come ``insediamento fortificato'' nell'ambito proto-veneto neolitico. La posizione elevata suggerisce una funzione di controllo del territorio, forse di frontiera tra aree di influenza culturale diversa.

\subsection{Bellunese e Feltrino}

\textbf{Fonzaso (BL) -- Malga Campon di Monte Avena.} \label{app:avena_fonzaso} Il sito di Malga Campon, da considerarsi parte di un unico contesto con il sito di Pedavena (vedi sotto), si estende sulla cima del Monte Avena a circa 1.400~m di quota. Le evidenze più antiche risalgono al Paleolitico (Musteriano e Aurignaziano), legate allo sfruttamento della selce locale e alla caccia. La fase eneolitica, con una probabile prosecuzione dello sfruttamento della selce, è seguita da un'occupazione medievale legata all'alpeggio. Il catalogo classifica il sito anche come neolitico, verosimilmente per la continuità con la fase eneolitica. Vedi anche la sezione Eneolitico~(\ref{app:eneolitico}).

\textbf{Pedavena (BL) -- Casera Campet di Monte Avena.} Estensione orientale del sito di Fonzaso nello stesso comune limitrofo. Le caratteristiche e la cronologia sono identiche (vedi la scheda di Malga Campon, \ref{app:avena_fonzaso}): la separazione nel catalogo riflette solo il confine comunale, non una distinzione funzionale o cronologica.

\subsection{Cadore e Comelico}

\textbf{Selva di Cadore (BL) -- Riparo Mandriz.} \label{app:mandriz} Un riparo sotto roccia frequentato come ricovero stagionale tra il Neolitico tardo e l'Eneolitico da comunità dedite alla pastorizia transumante e alla caccia. La quota e il contesto montano indicano che anche in quest'epoca le valli cadorine erano percorse da gruppi mobili che sfruttavano i pascoli estivi. Vedi anche la sezione Eneolitico~(\ref{app:eneolitico}).

\subsection{Colli Euganei e Berici}

\textbf{Arcugnano (VI) -- Lago di Fimon (tre schede).} \label{app:fimon} Il lago di Fimon, nelle colline Beriche, ha restituito tracce di almeno tre settori distinti di un grande abitato perilacustre neolitico (cultura dei Vasi a Bocca Quadrata, ambito proto-veneto). Le tre schede separate nel catalogo corrispondono a settori con funzioni diverse: abitativa, artigianale e funeraria. Il complesso nel suo insieme costituisce uno dei più importanti insediamenti neolitici dell'Italia nord-orientale, con una notevole durata di occupazione e una struttura spaziale differenziata.

\textbf{Longare (VI).} Materiali riconducibili al Neolitico (ambito proto-veneto) raccolti in superficie; in assenza di strutture individuate, si interpreta come frequentazione a carattere prevalentemente abitativo. Il sito si aggiunge al quadro della presenza neolitica nei Berici.

\textbf{Pojana Maggiore (VI).} Insediamento neolitico in pianura, attribuibile a un abitato sulla base della tipologia del materiale rinvenuto. La scheda è sintetica.

\textbf{Este (PD).} \label{app:este} Sito pluristratificato emerso presso lo scolo delle Monache, con fasi dal Neolitico finale all'età del Bronzo. L'insediamento aveva una zona in ambiente umido (probabilmente su palafitte) e una in asciutto su dosso. La scansione cronologica è difficile da precisare, ma l'occupazione inizia con una fase neo-eneolitica seguita da almeno tre fasi del Bronzo con diversa organizzazione spaziale. Vedi anche le sezioni Eneolitico~(\ref{app:eneolitico}) e Bronzo~(\ref{app:este_bronzo}).

\subsection{Lago di Garda e Adige}

\textbf{Garda (VR) -- Monte Luppia e dintorni.} \label{app:garda_neo} Sito pluristratificato di straordinaria continuità che abbraccia il Neolitico (cultura dei Vasi a Bocca Quadrata), l'Eneolitico, l'età del Ferro (Veneti antichi e seconda età del Ferro), l'età romano-tarda, il periodo altomedievale con sepolture longobarde, il Medioevo scaligero e l'epoca contemporanea con gallerie e trincee della prima guerra mondiale. La scheda non specifica la natura delle evidenze neolitiche ed eneolitiche, ma la stratificazione eccezionale del promontorio gardesano è ben nota dalla letteratura regionale. Per la fase più importante ai fini di questo libro, quella protostorica con i petroglifi, vedi la sezione Bronzo~(\ref{app:garda_petroglifi}). L'occupazione registrata si estende fino al periodo contemporaneo.

\subsection{Lessinia}

\textbf{Fumane (VR) -- Seconda scheda.} \label{app:fumane_neo} Seconda scheda MiC per Fumane, relativa a un sito diverso dalla Grotta di Fumane paleolitica (vedi \ref{app:fumane_paleo}): insediamento e area funeraria con fasi neolitiche (cultura dei Vasi a Bocca Quadrata) ed eneolitiche, scavate a partire dal 1876 da A.~Goiran. Scavi successivi di De Stefani, De Mortillet, Tedeschi e Battaglia hanno via via sconvolto i depositi originali, rendendo difficile la ricostruzione stratigrafica. Il catalogo segnala anche una frequentazione nel Paleolitico superiore, che connette questa scheda al contesto più ampio della Valle dei Progni. Vedi anche la sezione Eneolitico~(\ref{app:eneolitico}).

\textbf{Negrar (VR).} \label{app:negrar} Abitato tra la fine del Neolitico medio e l'Eneolitico, con fasi culturali riferibili ai Vasi a Bocca Quadrata, alla cultura ``Fontbouisse'', allo stile campaniforme e allo ``White Ware''. La varietà degli aspetti culturali indica un luogo di incontro tra influenze diverse, coerente con la posizione della Valpolicella come zona di passaggio tra la pianura padana e le Alpi. Vedi anche la sezione Eneolitico~(\ref{app:eneolitico}).

\textbf{Rivoli Veronese (VR).} \label{app:rivoli} La Rocca di Rivoli, che domina la stretta dell'Adige, ha conservato tracce di insediamento e necropoli neolitici (Vasi a Bocca Quadrata e Lagozza), seguiti da occupazioni dell'antico e tardo Bronzo (cultura di Polada, fase Protovillanoviana) e poi da un abitato altomedievale e un castello bassomedievale. Le ricerche di G.~Pellegrini e L.H.~Barfield hanno identificato sei aree distinte sul sito. Vedi anche la sezione Bronzo~(\ref{app:rivoli_bronzo}).

\textbf{Sant'Anna d'Alfaedo (VR).} \label{app:alfaedo_neo} Vedi la descrizione completa alla sezione Paleolitico. La fase neolitica (tecnica Campignana) si innesta su una lunga sequenza che sale dal Paleolitico inferiore. Le attività di scheggiatura della selce continuano sino all'età del rame, rendendo questo sito uno dei più rappresentativi per la comprensione dell'economia litica della Lessinia.

\textbf{Sant'Ambrogio di Valpolicella (VR).} \label{app:santambrogio_neo} Sito pluristratificato con fase della cultura dei Vasi a Bocca Quadrata, cultura retica e periodo romano. La scheda è sintetica; la continuità dalla preistoria all'età romana è comunque significativa per comprendere la longevità dell'insediamento in Valpolicella.

\subsection{Pianura Veronese e Bassa}

\textbf{Erbè (VR).} \label{app:erbe} Sito pluristratificato con fasi neolitiche (Vasi a Bocca Quadrata), Bronzo antico e paleoveneto. La scheda è sommaria; il sito si colloca nella pianura irrigua veronese.

\textbf{Lavagno (VR).} \label{app:lavagno} Sito pluristratificato che copre un arco cronologico eccezionalmente lungo: le evidenze registrate partono dall'Eneolitico e si estendono attraverso il Bronzo tardo, l'età del Ferro, l'epoca romana, il Medioevo e il periodo austriaco. La collocazione in pianura veronese, in una zona di bassa collina, ha favorito la sedimentazione di cultura su cultura. Vedi anche le sezioni Eneolitico~(\ref{app:eneolitico}), Bronzo~(\ref{app:lavagno_bronzo}) e Ferro~(\ref{app:lavagno_ferro}).

\subsection{Prealpi Trevigiane e Pedemontana}

\textbf{Cornuda (TV).} Insediamento neolitico recente (cultura dei Vasi a Bocca Quadrata, fase più tarda con influenze Chassey-Lagozza) identificato dopo una frana che ha rimesso in luce materiali in giacitura secondaria. Il recupero è avvenuto in parte durante la rimozione di terre accumulate durante la seconda guerra mondiale. Alcune evidenze indicano una persistenza del sito nel Neolitico finale.

\textbf{Revine Lago (TV).} \label{app:revine} Villaggio su bonifica di tardo Neolitico e inizio Eneolitico (fine del quarto-inizi del terzo millennio a.p.z.a.) sul bordo dei Laghi di Revine, con frequentazione più sporadica nel corso dell'antico Bronzo (attorno al diciottesimo secolo a.p.z.a.). I primi materiali furono recuperati nel 1923 durante scavi per un canale; successivi scavi clandestini nel 1987 spinsero la Soprintendenza a organizzare campagne sistematiche nel 1992 e 1996. Il sito ha restituito materiali fittili, litici, resti faunistici e paleobotanici di grande interesse per la ricostruzione dell'economia e dell'ambiente preistorici. Dal fondo lacustre provengono anche armi in bronzo della media età del Bronzo (spade tipo Sauerbrunn, pugnale tipo Peschiera), che indicano una persistenza o una rifrequentazione del luogo. Vedi anche le sezioni Eneolitico~(\ref{app:eneolitico}) e Bronzo~(\ref{app:revine_bronzo}).

\subsection{Prealpi Vicentine e Piccole Dolomiti}

\textbf{Cornedo Vicentino (VI).} Riparo sotto roccia neolitico (ambito proto-veneto) probabilmente attribuibile a un abitato, senza strutture identificate. Il sito è uno dei pochi segnalati nelle Prealpi vicentine per questo periodo.

\subsection*{Nota di sintesi -- Neolitico}

Con ventidue schede il Neolitico è il secondo periodo per numerosità nel catalogo veneto (dopo il Bronzo). La distribuzione mostra una forte concentrazione nella pianura e nella bassa collina (Berici, bassa Lessinia, pianura veronese), che rispecchia sia la vocazione agricola delle comunità neolitiche sia la maggiore accessibilità di quelle aree per i ricercatori. I siti perilacustri (Fimon, Revine Lago) e quelli in ambiente umido (Este) indicano una preferenza per le risorse idriche. Le quote montane (Monte Avena, Mandriz, Cornedo Vicentino) suggeriscono che il territorio alpino fosse già percorso da pastori e cacciatori neolitici, ma le schede sono sparse e la ricerca di alta quota è quasi assente. Il grande complesso di Fimon, con le sue tre schede distinte, è di fatto il sito neolitico meglio documentato dell'intera regione.

% ============================================================
\section{Eneolitico}
\label{app:eneolitico}
% 9 siti (escludendo i rimandi)
% ============================================================

L'Eneolitico (età del Rame, circa 3.500-2.300 a.p.z.a.) è il periodo di diffusione della metallurgia del rame e dei primi metalli, con continuità nelle tradizioni neolitiche ma nuove reti di scambio a lunga distanza. Nel Veneto è spesso difficile distinguere Neolitico finale ed Eneolitico, e molte schede li trattano insieme.

\subsection{Bellunese e Feltrino}

\textbf{Fonzaso (BL) e Pedavena (BL) -- Monte Avena.} Vedi la descrizione completa nella sezione Neolitico~(\ref{app:avena_fonzaso}). La fase eneolitica di Monte Avena è interpretata come legata sia all'estrazione e lavorazione della selce sia alla caccia, in uno dei pochi contesti di alta quota catalogati per questo periodo nel Veneto.

\subsection{Cadore e Comelico}

\textbf{Selva di Cadore (BL) -- Riparo Mandriz.} Vedi la descrizione completa nella sezione Neolitico~(\ref{app:mandriz}). La frequentazione eneolitica si affianca a quella tardo-neolitica senza soluzione di continuità, indicando un utilizzo pluridecennale del riparo come base per pastori transumanti.

\subsection{Colli Euganei e Berici}

\textbf{Este (PD).} Vedi la descrizione nella sezione Neolitico~(\ref{app:este}). La fase eneolitica rientra nella lunga continuità insediativa del sito presso lo scolo delle Monache.

\subsection{Lago di Garda e Adige}

\textbf{Garda (VR).} Vedi la descrizione nella sezione Neolitico~(\ref{app:garda_neo}). La fase eneolitica fa parte della lunga stratificazione del promontorio gardesano.

\subsection{Lessinia}

\textbf{Fumane (VR).} Vedi la descrizione nella sezione Neolitico~(\ref{app:fumane_neo}). La fase eneolitica del sito di Fumane (seconda scheda) si affianca a quella neolitica.

\textbf{Negrar (VR).} Vedi la descrizione nella sezione Neolitico~(\ref{app:negrar}). L'abitato tra Neolitico recente ed Eneolitico è il nucleo principale del sito.

\subsection{Pianura Veronese e Bassa}

\textbf{Lavagno (VR).} Vedi la descrizione nella sezione Neolitico~(\ref{app:lavagno}). L'Eneolitico è la fase iniziale registrata nel catalogo per questo sito pluristratificato.

\subsection{Prealpi Trevigiane e Pedemontana}

\textbf{Revine Lago (TV).} Vedi la descrizione nella sezione Neolitico~(\ref{app:revine}). Il tardo Neolitico e l'Eneolitico costituiscono la fase principale del villaggio lacustre di Colmaggiore.

\subsection*{Nota di sintesi -- Eneolitico}

L'Eneolitico veneto nel catalogo MiC è in larga parte documentato come prolungamento o fase terminale dei siti neolitici già noti: quasi tutte le schede eneolitiche sono in realtà schede pluristratificate con fasi precedenti o successive. Il dato più interessante ai fini di questo libro è la presenza di siti in quota (Monte Avena, Riparo Mandriz) che indicano una mobilità verticale già strutturata durante l'età del Rame, con una frequentazione stagionale delle aree montane non dissimile da quella che potremmo aspettarci per i contesti cultuali alpini. La totale assenza di arte rupestre eneolitica catalogata in Veneto, a fronte della sua abbondanza in contesti analoghi (Valcamonica, Monte Bego), costituisce uno dei paradossi centrali di questo libro.

% ============================================================
\section{Età del Bronzo}
\label{app:bronzo}
% 42 siti
% ============================================================

L'età del Bronzo (circa 2.300-900 a.p.z.a.) è il periodo meglio rappresentato nel catalogo veneto, con quarantadue schede. È l'epoca dei castellieri, delle palafitte gardesane, degli insediamenti arginati di pianura e dei primi sistemi territoriali complessi. I metalli circolano su lunga distanza; l'allevamento transumante porta i pastori verso le quote alpine.

\subsection{Bellunese e Feltrino}

\textbf{Sedico (BL) -- Noàl.} \label{app:sedico} Castelliere molto attivo nell'età del Bronzo recente, inserito in una rete di siti che controllavano i versanti della valle del Piave. Attorno al nono-ottavo secolo a.p.z.a. Noàl faceva parte di un sistema di altura coordinato per il controllo delle rotte metallurgiche. La fase medievale (castello con torre) sfruttò la stessa posizione strategica, a controllo dell'imboccatura della valle del Cordevole.

\textbf{Limana (BL) -- San Pietro in Tuba.} \label{app:limana} Sito pluristratificato che copre la transizione Bronzo Recente-Finale e la prima età del Ferro, poi rioccupato in pieno Medioevo come castello di frontiera tra Trevigiano e Bellunese (distrutto nel 1366). Il dato più rilevante per il tema del libro è la segnalazione di un possibile contesto di deposizioni cultuali di manufatti in bronzo nei pressi di una sorgente, un tipo di evidenza che ricorre in molti santuari protostorici alpini. Vedi anche la sezione Ferro~(\ref{app:limana_ferro}).

\textbf{San Gregorio nelle Alpi (BL) -- Castel de Pedena.} \label{app:pedena} Sito d'altura con una prima fase del Bronzo medio-recente senza strutture difensive permanenti, a cui seguono, attorno al dodicesimo-undicesimo secolo a.p.z.a., strutture difensive vere e proprie che si mantengono fino alla prima età del Ferro. Dopo un lungo iato, il colle conosce una rioccupazione medievale strategica. Vedi anche la sezione Ferro~(\ref{app:pedena_ferro}).

\textbf{Belluno -- Col del Buson.} \label{app:colbuson} Abitato d'altura di età tardo-neolitica ed eneolitica (la scheda compare sotto il Bronzo nel catalogo per le fasi del Bronzo finale e dell'età del Ferro), integrato nell'ambiente circostante con pascoli e terreni agricoli. La documentazione di una struttura per la lavorazione della selce e di indizi di metallurgia in situ indica una comunità autosufficiente anche tecnologicamente. Vedi anche la sezione Ferro~(\ref{app:colbuson_ferro}).

\textbf{Belluno -- Sant'Anna di Pedecastello.} Tracce di frequentazione del Bronzo medio-recente ai margini della sommità del colle, con strutture principali attribuibili all'alto Medioevo e al Medioevo pieno. Le evidenze bronzee sono marginali rispetto alle fasi successive.

\subsection{Cadore e Comelico}

\textbf{San Vito di Cadore (BL) -- Sito VF1.} Definito di ``straordinaria importanza'' per la presenza di una sepoltura castelnoviana (Mesolitico), il sito ha visto anche frequentazioni dell'età del Bronzo e in epoca sub-attuale. La sovrapposizione cronologica su un arco di quasi diecimila anni ne fa un punto di riferimento per la comprensione della mobilità umana in Cadore.

\subsection{Colli Euganei e Berici}

\textbf{Lozzo Atestino (PD).} Necropoli a incinerazione del Bronzo finale; la scoperta di almeno due sepolture e di chiazze di terreno scuro nei campi vicini suggerisce l'esistenza di un'area funeraria più estesa nella zona di Vignalon.

\textbf{Este (PD).} \label{app:este_bronzo} Vedi la descrizione nella sezione Neolitico~(\ref{app:este}). Le almeno tre fasi bronzee del sito mostrano una riorganizzazione progressiva degli spazi insediativi.

\textbf{Montegrotto Terme (PD) -- Insediamento.} Frequentazione della media-tarda età del Bronzo (diciassettesimo-dodicesimo secolo a.p.z.a.) su un'area abbastanza estesa, con probabili attività agro-pastorali.

\textbf{Montegrotto Terme (PD) -- Villa ex Piacentini.} \label{app:montegrotto} Sito pluristratificato con prima frequentazione sporadica nell'età del Rame (prima metà del terzo millennio a.p.z.a.), insediamento stabile del Bronzo (quattordicesimo-dodicesimo secolo a.p.z.a.), villa romana monumentale di prima-quarta età imperiale, piccolo villaggio con necropoli altomedievale (ottavo-nono secolo) e insediamento produttivo più stabile tra il nono e il quattordicesimo secolo. La villa romana, interpretata come \textit{villa d'otium} per un committente di altissimo rango, è l'elemento più eccezionale della sequenza.

\subsection{Lago di Garda e Adige}

\textbf{Lazise (VR) -- Sito 1.} \label{app:lazise1} Insediamento palafitticolo con due fasi: una più antica tra il 1939 e il 1844 a.p.z.a. (analisi dendrocronologiche), l'altra della media età del Bronzo, con una terza frequentazione nell'età del Ferro. I materiali sono in parte frutto di scavi clandestini; un saggio sistematico è stato avviato nel 1986. Vedi anche la sezione Ferro~(\ref{app:lazise_ferro}).

\textbf{Garda (VR) -- Petroglifi di Monte Luppia.} \label{app:garda_petroglifi} La scheda MiC relativa ai petroglifi è quella di un ``sito pluristratificato'' del Bronzo e del Ferro, con la sintetica interpretazione: ``Petroglifi databili genericamente tra l'età del bronzo e quella del ferro.'' Monte Luppia è il sito di arte rupestre più noto del Veneto orientale, con incisioni di tipo geometrico e figure stilizzate su roccia calcarea affiorante. La catalogazione generica riflette le difficoltà di datazione delle incisioni in assenza di contesti stratigrafici chiusi; la letteratura specialistica propende per una datazione prevalentemente al Bronzo, con possibili continuazioni nel Ferro. Vedi anche la sezione Ferro~(\ref{app:garda_ferro}).

\textbf{Lazise (VR) -- Sito 2.} Secondo insediamento palafitticolo a Lazise, distinto dal precedente, datato tra la fase recente dell'antico Bronzo e la media età del Bronzo, con un rilievo topografico di 88 pali emergenti dal fondale del lago di Garda.

\textbf{Peschiera del Garda (VR) -- Sito 1.} Insediamento palafitticolo tra l'antico e il Bronzo recente, parzialmente sconvolto da scavi clandestini.

\textbf{Peschiera del Garda (VR) -- Sito 2.} Insediamento palafitticolo di grande estensione con due fasi principali: antico Bronzo (analisi dendrocronologiche: circa 2060 a.p.z.a.) e media età del Bronzo (circa 1617 a.p.z.a.). Seicento campioni di legno recuperati.

\subsection{Lessinia}

\textbf{Fumane (VR) -- Insediamento Bronzo recente.} Abitato probabilmente fortificato del Bronzo recente, con muretti a secco forse di epoca preistorica. Sito distinto dalla grotta paleolitica e dall'insediamento neolitico-eneolitico sopra descritti.

\textbf{Negrar (VR) -- Castelliere.} Insediamento, probabile castelliere, del Bronzo medio (cultura di Polada). Scheda sintetica.

\textbf{Rivoli Veronese (VR).} \label{app:rivoli_bronzo} Vedi la descrizione nella sezione Neolitico~(\ref{app:rivoli}). Le fasi bronzee includono il Bronzo antico (cultura di Polada) e il Bronzo finale (cultura Protovillanoviana).

\textbf{Sant'Anna d'Alfaedo (VR) -- Insediamento Bronzo.} Abitato probabilmente fortificato del Bronzo medio e recente, distinto dal sito paleolitico-neolitico della stessa area.

\subsection{Pianura Veronese e Bassa}

\textbf{Monteforte d'Alpone (VR).} \label{app:monteforte} Insediamento d'altura con sequenza dal Bronzo recente-finale all'età del Ferro, con fasi abitative documentate al dodicesimo-decimo secolo a.p.z.a., al nono, all'ottavo-settimo, al sesto-quinto e al secondo-primo secolo a.p.z.a. La continuità di occupazione per oltre un millennio è significativa. Vedi anche la sezione Ferro~(\ref{app:monteforte_ferro}).

\textbf{Nogarole Rocca (VR).} Insediamento tra la media età del Bronzo e il Bronzo finale. Nel 1977 fu recuperata una figurina fittile zoomorfa del Bronzo recente; nel corso dei sopralluoghi degli anni Settanta anche un coltello in bronzo tipo Vadena del Bronzo finale avanzato.

\textbf{Bovolone (VR).} Abitato del Bronzo con struttura abitativa a capanna, datato alla fase iniziale della media età del Bronzo.

\textbf{Colognola ai Colli (VR).} \label{app:colognola} Abitato su altura con fase del Bronzo medio e poi presenza della seconda età del Ferro (cultura retica e paleoveneta), seguita da un insediamento medievale. Vedi anche la sezione Ferro~(\ref{app:colognola_ferro}).

\textbf{Erbè (VR).} Vedi la descrizione nella sezione Neolitico~(\ref{app:erbe}). La fase bronzea si innesta sull'insediamento neolitico e precede quella paleoveneta.

\textbf{Legnago (VR).} Sito pluristratificato con Bronzo recente e cultura paleoveneta. Scheda sintetica.

\textbf{Roncà (VR).} \label{app:ronca_bronzo} Vedi la descrizione nella sezione Paleolitico~(\ref{app:ronca_paleo}). Le fasi del Bronzo recente-finale si inseriscono nella lunga stratificazione del sito.

\textbf{Villa Bartolomea (VR) -- Sito 1.} Sito pluristratificato con fase del Bronzo recente e periodo romano.

\textbf{Villa Bartolomea (VR) -- Sito 2.} Sito pluristratificato con fase del Bronzo medio-recente e periodo romano.

\textbf{Cerea (VR).} Sito palafitticolo in pianura veronese. Scheda sintetica (``sito palafitticolo'').

\textbf{Caldiero (VR).} Abitato d'altura protostorico (Bronzo medio-finale e cultura paleoveneta) con successivo insediamento medievale e post-medievale.

\textbf{Lavagno (VR).} \label{app:lavagno_bronzo} Vedi la descrizione nella sezione Neolitico~(\ref{app:lavagno}). La fase del Bronzo tardo si inserisce nella lunga stratificazione del sito.

\textbf{Nogara (VR).} Vasta zona a rischio archeologico con fase del Bronzo medio-recente, romana e medievale. La scheda è una ``zona di interesse'' senza strutture puntualmente identificate.

\subsection{Pianura e Delta}

\textbf{Dolo (VE).} Tracce di frequentazione del Bronzo recente 1 (seconda metà del quattordicesimo-inizio tredicesimo secolo a.p.z.a.) emerse durante i lavori di scavo per il Passante di Mestre tra il 2011 e il 2012. Il sito mostra attività agricole e di allevamento (canalette, pozzi, pozzetti, uno scarico di ceramica, una tomba di infante), poi sigillato da un evento alluvionale. Una frequentazione agricola sporadica dell'età del Ferro è successiva all'abbandono bronzeo.

\textbf{Arquà Petrarca (PD).} Insediamento palafitticolo perilacustre della cultura di Polada (Bronzo antico), ai piedi dei Colli Euganei.

\textbf{San Martino di Lupari (PD).} Insediamento di abitato arginato dell'ambito culturale Sub-appenninico (Bronzo recente-finale). La struttura ad argine trova confronti con gli insediamenti della pianura friulana.

\subsection{Prealpi Trevigiane e Pedemontana}

\textbf{Castello di Godego (TV).} Abitato fortificato del Bronzo medio-recente a forma quadrilatera con argine (``Le Motte''), che si estende in parte nel comune di San Martino di Lupari. Interpretato come vallo romano o medievale fino agli anni Settanta del Novecento, fu riconosciuto come preistorico solo dagli scavi del 1984-1985.

\textbf{Cordignano (TV).} \label{app:cordignano_bronzo} Le testimonianze per le fasi più antiche (Bronzo avanzato sub-appenninico e Bronzo finale protovillanoviano) sono frammentarie. Il sito diventa pienamente leggibile solo con la grande area sacra della prima età del Ferro e dell'età del Ferro veneto-antica (sesto-quarto secolo a.p.z.a.), attiva fino al quarto secolo. Vedi anche la sezione Ferro~(\ref{app:cordignano_ferro}).

\textbf{Montebelluna (TV).} \label{app:montebelluna_bronzo} Sito pluristratificato con frequentazioni documentate dal Bronzo medio-recente all'età moderna. Le fasi più significative nel catalogo sono quelle del Ferro e romane (complessi abitativi, impianti artigianali, necropoli), ma le radici bronzee del sito sono documentate. Vedi anche la sezione Ferro~(\ref{app:montebelluna_ferro}).

\textbf{Revine Lago (TV).} \label{app:revine_bronzo} Vedi la descrizione nella sezione Neolitico~(\ref{app:revine}). Le armi in bronzo recuperate dal fondo lacustre (spade tipo Sauerbrunn, pugnale tipo Peschiera) indicano frequentazioni anche nella media età del Bronzo.

\textbf{San Zenone degli Ezzelini (TV).} Insediamento d'altura a frequentazione stagionale di gruppi agro-pastorali, datato alla media età del Bronzo avanzata e al Bronzo medio-recente (quattordicesimo-tredicesimo secolo a.p.z.a.). I materiali provengono dalla sommità del colle e dalle pendici.

\textbf{Tarzo (TV).} \label{app:tarzo_bronzo} Vedi la descrizione nella sezione Paleolitico~(\ref{app:tarzo_paleo}). La fase bronzea aggiunge una terza frequentazione documentata per questo sito pluristratificato.

\textbf{Vedelago (TV).} Tracce di frequentazione di gruppi umani della cultura sub-appenninica nel Bronzo recente (quattordicesimo-tredicesimo secolo a.p.z.a.).

\subsection*{Nota di sintesi -- Età del Bronzo}

Con quarantadue schede, il Bronzo è il periodo più documentato nel catalogo veneto. La distribuzione geografica è molto più equilibrata rispetto alle epoche precedenti: palafitte sul Garda, castellieri nella pianura veronese e nelle Prealpi, insediamenti arginati nella pianura trevigiana, siti d'altura nel Bellunese. La sistematica presenza di controllo sui passi e sulle valli fluviali (Noàl sulla val Cordevole, Castel de Pedena, Rivoli sull'Adige, Monteforte sull'Alpone) indica un territorio organizzato in reti di sorveglianza e scambio. Il dato più rilevante per il tema di questo libro è la quasi totale assenza di arte rupestre catalogata per questo periodo nel Veneto, in contrasto con la ricchezza delle incisioni camune e liguri coeve: i petroglifi di Garda (Monte Luppia) sono l'unica eccezione esplicita nel catalogo, e anche lì la scheda è sintetica e la datazione ``generica''.

% ============================================================
\section{Età del Ferro}
\label{app:ferro}
% 31 siti
% ============================================================

L'età del Ferro (circa 900-0 a.p.z.a.) coincide nel Veneto con lo sviluppo della cultura paleoveneta (Este) e retica, con i suoi santuari delle acque, le necropoli a incinerazione e la metallurgia del ferro. È il periodo che precede la romanizzazione e che produce, nel catalogo MiC, la maggior varietà di tipologie di siti: insediamenti, necropoli e, notevolmente, santuari con depositi votivi.

\subsection{Bellunese e Feltrino}

\textbf{Mel (BL) -- Necropoli.} \label{app:mel_nec} Necropoli dell'età del Ferro a incinerazione con tombe a cassetta litica sormontate da tumuli sepolcrali, simile tipologicamente alle necropoli coeve di Este. Testimonianza dell'area di influenza paleoveneta nel Bellunese.

\textbf{Limana (BL).} \label{app:limana_ferro} Vedi la descrizione nella sezione Bronzo~(\ref{app:limana}). La fase del Ferro continua la vita del castelliere prima dell'abbandono e della successiva rioccupazione medievale. Il possibile deposito cultuale presso la sorgente rimane il dato più suggestivo del sito.

\textbf{San Gregorio nelle Alpi (BL).} \label{app:pedena_ferro} Vedi la descrizione nella sezione Bronzo~(\ref{app:pedena}). Il castelliere di Castel de Pedena è attivo fino alla prima età del Ferro.

\textbf{Belluno -- Col del Buson.} \label{app:colbuson_ferro} Vedi la descrizione nella sezione Bronzo~(\ref{app:colbuson}). La fase del Ferro completa la sequenza dell'abitato d'altura bellunese.

\textbf{Sedico (BL) -- Santuario di Libano, loc. Costiet.} Struttura in pietra a secco con materiali metallici, tra cui laminette in bronzo tipiche dei contesti cultuali venetici (tipo Lagole), una moneta romana in bronzo e oggetti in ferro. L'interpretazione è quella di un piccolo santuario frequentato durante l'età del Ferro e l'età romana. La mancanza di elementi stratigrafici chiari impedisce di stabilire se la frequentazione fosse continua o con interruzioni. La tipologia delle laminette collega questo sito alla rete di santuari delle acque venetici del Veneto orientale.

\textbf{Mel (BL) -- Abitato.} Abitato della seconda età del Ferro, con alcune strutture interpretabili come unità abitative (proprietà Calzavara e Rui-Sciulli) e numerose strutture lapidee probabilmente legate al drenaggio del substrato impermeabile. La documentazione disponibile consente di riconoscere un areale di diffusione antropica nell'attuale località Cioppa.

\subsection{Cadore e Comelico}

\textbf{Lozzo di Cadore (BL) -- Riva del Brodevin.} Necropoli con due fasi d'uso: la prima tra il settimo e il quarto secolo a.p.z.a., la seconda tra la fine del primo secolo a.p.z.a. e la seconda metà del quarto secolo. La continuità d'uso tra il Ferro e l'età tardo-romana indica la persistenza dell'abitato sottostante.

\textbf{Calalzo di Cadore (BL) -- Santuario di Lagole.} \label{app:lagole} Santuario territoriale venetico con continuità di frequentazione dalla metà del sesto secolo a.p.z.a. fino al quarto secolo. Il catalogo segnala una componente celtica nelle prime fasi e un culto legato alle acque curative, testimoniato da numerosi ex-voto anatomici. L'assenza di figurine femminili suggerisce che il culto fosse prevalentemente rivolto alla componente maschile della popolazione. Lagole è il sito rituale più importante del Cadore: la sua esistenza, insieme al santuario di Costiet (Sedico), indica che anche le vallate alpine erano attraversate da reti di sacralità riconoscibile.

\subsection{Colli Euganei e Berici}

\textbf{Barbarano Vicentino (VI).} Terrazzamenti artificiali attribuiti all'ambito veneto pre-protostorico. Scheda sintetica senza ulteriori dettagli.

\textbf{Montegrotto Terme (PD) -- Santuario.} Area ad ovest del Colle di S.~Pietro Montagnone con luogo di culto delle acque curative, attivo tra il settimo e il terzo secolo a.p.z.a. e testimoniato dagli ex-voto anatomici. Il collegamento tra sacralità delle acque termali e culto protostorico è significativo in un'area che diventerà poi famosa in età romana per i bagni di Aquae Patavinae.

\subsection{Lago di Garda e Adige}

\textbf{Lazise (VR).} \label{app:lazise_ferro} Vedi la descrizione nella sezione Bronzo~(\ref{app:lazise1}). La frequentazione del Ferro si aggiunge alle fasi palafitticole bronzee.

\textbf{Garda (VR) -- Petroglifi di Monte Luppia.} \label{app:garda_ferro} Vedi la descrizione nella sezione Bronzo~(\ref{app:garda_petroglifi}). I petroglifi sono databili genericamente tra l'età del Bronzo e quella del Ferro.

\textbf{Garda (VR) -- Sito pluristratificato.} Vedi la descrizione nella sezione Neolitico~(\ref{app:garda_neo}). La fase del Ferro (Veneti antichi e seconda età del Ferro) fa parte della lunga stratificazione del promontorio.

\subsection{Lessinia}

\textbf{Velo Veronese (VR).} Abitato d'altura probabilmente circondato da un vallo, dell'età del Ferro seconda fase. Due asce ritrovate nel 1854 durante la fondazione della chiesa locale; ceramiche eneolitiche e ossa recuperati nel 1948 sulla sommità; deposito di selci raccolto quasi sulla vetta del monte.

\subsection{Pianura Veronese e Bassa}

\textbf{Monteforte d'Alpone (VR).} \label{app:monteforte_ferro} Vedi la descrizione nella sezione Bronzo~(\ref{app:monteforte}). Le fasi del Ferro coprono l'ottavo-settimo e il sesto-quinto secolo a.p.z.a., fino alla fine dell'età del Ferro al secondo-primo secolo a.p.z.a.

\textbf{Cologna Veneta (VR).} Abitato e vasta necropoli paleoveneta; la scheda segnala anche una fase medievale. Sito rilevante per la rete funeraria del territorio paleoveneto nella pianura orientale veronese.

\textbf{Colognola ai Colli (VR).} \label{app:colognola_ferro} Vedi la descrizione nella sezione Bronzo~(\ref{app:colognola}). La seconda età del Ferro (cultura retica e paleoveneta) è la fase meglio documentata dell'abitato d'altura.

\textbf{Gazzo Veronese (VR) -- Abitato.} Abitato dell'età del Ferro. Scheda sintetica.

\textbf{Gazzo Veronese (VR) -- Necropoli.} Necropoli paleoveneta del periodo Atestino III, nella pianura veronese meridionale.

\textbf{Roncà (VR).} \label{app:ronca_ferro} Vedi la descrizione nella sezione Paleolitico~(\ref{app:ronca_paleo}). La tarda età del Ferro e la fase di romanizzazione completano la lunga sequenza del sito.

\textbf{Verona (VR).} Area pluristratificata nel centro storico di Verona, con fasi dall'età del Ferro (paleoveneto e retico) all'epoca moderna. La scheda catalogata testimonia la profondità stratigrafica del centro urbano.

\textbf{Lavagno (VR).} \label{app:lavagno_ferro} Vedi la descrizione nella sezione Neolitico~(\ref{app:lavagno}). La fase del Ferro si inserisce nella lunga stratificazione del sito pianeggiante.

\subsection{Pianura e Delta}

\textbf{Montagnana (PD).} Necropoli e abitato preromani (ambito veneto protostorico), con possibile villa rustica romana. La scheda segnala un'area funeraria del Ferro integrata nel paesaggio di pianura.

\textbf{Teolo (PD).} Materiali sporadici dell'ambito veneto protostorico, interpretati come tracce di frequentazione e/o di abitato. Scheda minima.

\subsection{Prealpi Trevigiane e Pedemontana}

\textbf{Cordignano (TV).} \label{app:cordignano_ferro} Vedi la descrizione nella sezione Bronzo~(\ref{app:cordignano_bronzo}). Il santuario territoriale di Cordignano diventa pienamente operativo a partire dalla fine del sesto secolo a.p.z.a. e rimane attivo fino al quarto secolo, quando gli editti di Teodosio vietano i culti pagani. Gli scavi dell'Università di Padova (1997-2006) hanno definito con precisione la grande area sacra, caratterizzata da una continuità devozionale di circa mille anni. Si tratta probabilmente di un santuario di frontiera, dedicato a divinità non ancora identificate, con componenti culturali veneto-antiche.

\textbf{Montebelluna (TV).} \label{app:montebelluna_ferro} Vedi la descrizione nella sezione Bronzo~(\ref{app:montebelluna_bronzo}). Il sito pluristratificato di Montebelluna mostra le sue fasi più ricche e meglio documentate nell'età del Ferro e in età romana (complessi abitativi, impianti artigianali, necropoli), con una continuità che va dal Bronzo recente fino all'epoca moderna.

\subsection{Prealpi Vicentine e Piccole Dolomiti}

\textbf{Rotzo (VI).} Sito d'altura di ambito veneto-retico, con posizione strategica all'imboccatura della valle dell'Astico. Scheda sintetica.

\textbf{Creazzo (VI).} Sito pluristratificato di ambito veneto-romano, collegato alla coltivazione dei campi e alla lavorazione e stoccaggio dei prodotti agricoli. Scheda sintetica.

\textbf{Montecchio Maggiore (VI).} Area di frequentazione preromana (ambito veneto) e insediamento rustico romano, con necropoli longobarda. Continuità insediativa dal Ferro all'alto Medioevo.

\textbf{Santorso (VI).} Sito pluristratificato di ambito veneto-romano, con strutture dell'abitato preromano e romano. Scheda sintetica.

\textbf{Isola Vicentina (VI).} Sito pluristratificato di ambito veneto-romano. I resti romani includono un insediamento rustico con un settore interpretato come \textit{fullonica} (impianto per la lavorazione della lana). Scheda sintetica.

\subsection*{Nota di sintesi -- Età del Ferro}

L'età del Ferro nel catalogo veneto rivela la mappa di una società complessa, organizzata in reti di insediamenti e santuari. I dati più preziosi per il tema di questo libro sono i santuari delle acque: Montegrotto Terme (Colli Euganei), Lagole di Calalzo (Cadore), Costiet di Sedico (Bellunese), il possibile deposito votivo di Limana. Questi siti mostrano che la sacralizzazione di punti fisici nel paesaggio, in particolare quelli legati all'acqua, era una pratica strutturata e geograficamente distribuita nell'intero territorio veneto, dalle pianure termali alle vallate alpine. L'assenza di santuari analoghi nell'arco prealpino centrale e nelle Piccole Dolomiti non significa che non esistessero: significa che non li abbiamo trovati, o che non li abbiamo riconosciuti. Il santuario di Cordignano (al confine con il Friuli), attivo per circa mille anni, dimostra la capacità del registro archeologico di conservare luoghi di culto quando le condizioni topografiche e stratigrafiche sono favorevoli.

% ============================================================
\section*{Conclusione dell'appendice}
% ============================================================

I 136 siti censiti in queste pagine non sono un'enciclopedia della preistoria veneta, ma sono l'inventario di ciò che il catalogo ufficiale riconosce come noto e formalmente catalogato. Le lacune sono strutturali: mancano quasi completamente i siti d'alta quota (sopra i 1.500~m), le aree di transizione Bellunese-Cadore non sono sistematicamente prospettate, e intere epoche (Mesolitico in particolare) sono quasi prive di schede pur essendo rappresentate nei siti limitrofi del Trentino e del Friuli.

L'utilità di questa raccolta, come strumento per il ragionamento proposto nel libro, è duplice. Da un lato fornisce una base dati concreta che permette di distinguere il silenzio dell'arte rupestre veneta dal silenzio dell'archeologia in generale: il Veneto non è vuoto di preistoria, anzi è ricco di siti, di culture, di occupazioni ripetute; è specificamente l'arte rupestre a mancare, e questo la dice lunga sui tipi di ricerca che sono stati condotti e sui tipi che mancano. Dall'altro lato, i santuari delle acque dell'età del Ferro (Lagole, Montegrotto, Costiet, Limana) e i siti d'altura del Bronzo (Monte Avena, San Zenone, Noàl) indicano che il territorio montano veneto era percorso, frequentato e sacralizzato: le condizioni per la produzione di arte rupestre esistevano. Aspettano solo di essere cercate.


\chapter*{Glossario}
\markboth{Glossario}{Glossario}
\addcontentsline{toc}{chapter}{Glossario}

Questo glossario raccoglie i termini tecnici usati nel libro. Le epoche preistoriche riportano sempre le date nella convenzione adottata nel testo: anni prima dell'anno zero astronomico, ottenute sottraendo circa 1.950 anni dai valori cal BP originali della letteratura.

\begin{description}

\item[Anni cal BP] Anni trascorsi rispetto al 1950, ricavati da datazioni al radiocarbonio e corretti (calibrati) per le variazioni storiche del tenore di $^{14}$C nell'atmosfera. La calibrazione è necessaria perché la produzione di $^{14}$C non è costante nel tempo: senza correzione, le date radiocarboniche risulterebbero sistematicamente più giovani del reale. La sigla BP sta per \textit{Before Present}, dove il ``presente'' è convenzionalmente fissato al 1950. In questo libro le date sono sempre espresse in anni prima dell'anno zero astronomico, sottraendo circa 1.950 anni dal valore cal BP.

\item[ARC model] Modello del costo energetico della locomozione terrestre in funzione della pendenza, proposto da Herman Pontzer (\textit{Biology Letters}, 2016). Il nome deriva dai due processi della fisiologia muscolare su cui si fonda: \textit{Activation-Relaxation} (attivazione e rilassamento delle fibre muscolari) e \textit{Cross-bridge cycling} (ciclo dei ponti actina-miosina). Il modello è ricavato dai meccanismi cellulari della contrazione, non da una curva adattata empiricamente, ed è stato validato su un intervallo di masse corporee da 0,78 grammi a 431 chilogrammi e su pendenze dalla discesa moderata all'arrampicata verticale, con buon accordo tra le previsioni e i dati misurati per specie molto diverse. Per le specie di interesse di questo libro (camoscio, stambecco, cervo, essere umano), tutte con massa corporea tra quaranta e novanta chilogrammi, il modello si trova nel centro del suo intervallo di validazione.

\item[Bølling-Allerød] Fase di miglioramento climatico rapido compresa tra circa 12.700 e 11.000 anni prima dell'anno zero, intercalata all'interno del Tardoglaciale. Le temperature medie estive sull'arco alpino aumentarono di diversi gradi in poche centinaia di anni; le foreste avanzarono rapidamente in quota e il limite del bosco si spostò verso l'alto di centinaia di metri. Il nome è composto da quello di due località del Nord Europa — Bølling in Danimarca e Allerød anch'essa danese — dove le sequenze polliniche che hanno permesso di identificare questa fase furono descritte per la prima volta nella prima metà del Novecento. La fase si conclude con il brusco raffreddamento del Dryas Recente.

\item[Cal BP] Vedi \textbf{Anni cal BP}.

\item[Castelnoviano] Cultura litica che contraddistingue la fase più recente del Mesolitico nell'Italia settentrionale, collocata approssimativamente tra 9.500 e 7.500 anni prima dell'anno zero. Prende il nome dal sito di Castellnovo d'Elsa (Siena), dove il materiale litico rappresentativo di questa fase fu descritto per la prima volta. Si distingue dal Sauveterriano precedente per la presenza di trapezi e lame più regolari nell'industria litica, e per criteri insediativi in parte diversi: i siti castelnoviani mostrano una maggiore preferenza per l'esposizione solare rispetto ai soli criteri di visibilità e controllo del territorio.

\item[Crioclastismo] Processo fisico di frattura della roccia causato dalla ripetuta alternanza di gelo e disgelo dell'acqua infiltrata nelle fessure. L'acqua, ghiacciando, aumenta di volume di circa il 9\% ed esercita pressioni che possono superare la resistenza alla trazione della roccia, provocandone la frantumazione e il distacco di frammenti. Il crioclastismo è particolarmente intenso durante i periodi di transizione climatica, quando il numero di cicli gelo-disgelo è più elevato. Alla Grotta di Fumane, in Lessinia, il distacco dei frammenti di parete recanti pitture paleolitiche è attribuito dagli scavatori a questo meccanismo.

\item[Criterio di informazione (AIC, BIC)] Misure statistiche per confrontare modelli concorrenti costruiti sugli stessi dati, che bilanciano la bontà di adattamento con la complessità del modello. Il criterio di Akaike (AIC, Hirotugu Akaike, 1974) penalizza ogni parametro aggiuntivo con un costo fisso pari a 2; il criterio bayesiano (BIC, Gideon Schwarz, 1978) applica una penalità che cresce con il logaritmo della dimensione del campione, risultando più severo con dataset grandi. Un modello con AIC o BIC più basso è preferibile a parità di tutte le altre condizioni. Entrambi i criteri sono usati per confrontare modelli spaziali nella letteratura sui siti mesolitici alpini.

\item[Dryas Recente] Fase di raffreddamento climatico brusca e marcata, compresa tra circa 11.000 e 9.700 anni prima dell'anno zero, che interruppe il miglioramento del Bølling-Allerød. Sull'arco alpino le temperature estive calarono di diversi gradi; in alcune zone i ghiacciai avanzarono nuovamente. La durata, circa 1.300 anni, la distingue dai raffreddamenti minori del Tardoglaciale. Il nome deriva da \textit{Dryas octopetala}, una pianta erbacea artica i cui pollini sono abbondanti nelle sequenze sedimentarie di questa fase e servono da indicatore paleoclimatico. Si conclude con il passaggio rapido all'Olocene, circa 9.700 anni prima dell'anno zero.

\item[Epigravettiano] Cultura litica diffusa in Italia e nelle aree adriatiche dell'Europa meridionale nella fase finale del Paleolitico superiore, approssimativamente tra 23.000 e 11.000 anni prima dell'anno zero. Il nome allude a una derivazione dal Gravettiano, cultura più antica. L'industria è caratterizzata da lamelle a dorso, punte a dorso e geometrici; nella fase più recente compaiono elementi che anticipano le industrie mesolitiche. Ai siti epigravettiani veneti appartengono, tra gli altri, Riparo Tagliente (fase superiore), Villabruna e il riparo Dalmeri, quest'ultimo datato con precisione al periodo del Bølling-Allerød.

\item[Home range] Area geografica abitualmente utilizzata da un individuo o da un gruppo per svolgere le proprie attività vitali (alimentazione, riproduzione, riposo), nell'arco di un anno o di una stagione. Non include spostamenti eccezionali né dispersioni. A differenza del territorio in senso etologico, l'home range non implica difesa attiva dai conspecifici. Per le specie alpine discusse in questo libro i valori tipici, misurati con radiotelemetria e GPS, sono: camoscio 200--600 ettari (stagionale), stambecco 300--800 ettari, cervo 1.000--5.000 ettari con variabilità legata al sesso e alla stagione.

\item[Laacher See Tephra] Strato di cenere vulcanica emessa dall'eruzione del vulcano Laacher See (odierna Renania-Palatinato, Germania) circa 11.050--10.950 anni prima dell'anno zero. L'eruzione, classificata come pliniana con indice di esplosività vulcanica 6, disperse materiale piroclastico su un'area di oltre un milione di chilometri quadrati. Lo strato di tephra, sottile e chimicamente omogeneo, si riconosce nelle sequenze sedimentarie lacustri e torbose dell'Europa centrale e settentrionale, incluse alcune località del nord Italia. Poiché si è depositato in un intervallo di anni o decenni su un'area vastissima, costituisce un livello isocronico usato dai geologi per sincronizzare stratigrafie di regioni diverse.

\item[Megaloceros giganteus] Cervide estinto, noto comunemente come cervo gigante o alce irlandese. Caratterizzato da palchi di dimensioni eccezionali, che potevano raggiungere i 3,7 metri di apertura, su un corpo di taglia paragonabile a quella di un alce moderno. Era diffuso dall'Irlanda all'Asia centrale durante il Pleistocene. Le ultime popolazioni si estinsero circa 7.700 anni prima dell'anno zero. Nei depositi del Paleolitico medio e superiore dei siti prealpini veneti i suoi resti compaiono sporadicamente, in misura nettamente inferiore rispetto ai cervidi di taglia più piccola.

\item[Mesolitico] Periodo culturale compreso tra la fine dell'ultima glaciazione e la comparsa dell'agricoltura in una data regione. Nelle Alpi orientali si colloca approssimativamente tra 12.700 e 7.500 anni prima dell'anno zero, corrispondendo al Tardoglaciale finale e all'Olocene antico. È caratterizzato da un'economia di caccia e raccolta con industria litica microlitica (armature geometriche, trapezi, microliti) e da una progressiva specializzazione nella caccia agli ungulati di montagna, in particolare stambecco e camoscio. Si suddivide convenzionalmente in Sauveterriano (fase più antica) e Castelnoviano (fase più recente).

\item[Modelli a pattern di punti (\textit{point-process models})] Famiglia di modelli statistici per l'analisi della distribuzione spaziale di eventi puntiformi (presenze di specie, siti archeologici, episodi di incendio) in un'area continua. Un modello di processo puntuale specifica come l'intensità attesa degli eventi varia nello spazio in funzione di covariate ambientali o di altra natura. Tra i modelli più usati in archeologia e paleoecologia figurano i processi di Poisson non omogenei e le loro estensioni. Permettono di confrontare, mediante criteri come AIC e BIC, l'importanza relativa di fattori concorrenti (morfologia, clima, sforzo di ricerca) nel determinare la distribuzione osservata dei siti noti.

\item[Neanderthal (\textit{Homo neanderthalensis})] Specie del genere \textit{Homo}, evolutasi in Europa e Asia occidentale a partire da circa 400.000 anni fa. Anatomicamente distinta da \textit{Homo sapiens} per il cranio basso e allungato, le arcate sopraorbitali prominenti, il prognatismo medio-facciale e la robustezza dello scheletro postcraniale. Capace di comportamenti culturalmente complessi: produzione di strumenti musteriani, uso del fuoco, sepolture intenzionali, e in alcuni contesti uso di pigmenti e ornamenti. Si estinse circa 38.000 anni prima dell'anno zero, in un periodo di sovrapposizione con \textit{Homo sapiens} nella penisola iberica e in Italia. I siti prealpini veneti attribuiti a Neanderthal --- tra cui il Riparo Tagliente nella sua fase più antica e i siti della Lessinia --- risalgono a un periodo compreso tra 58.000 e 38.000 anni prima dell'anno zero.

\item[Neolitico] Periodo culturale caratterizzato dalla comparsa e dalla diffusione dell'economia produttiva (agricoltura e allevamento), della ceramica e, in molte regioni, dell'insediamento stabile. In Europa centro-meridionale si diffonde a partire da circa 9.000--8.000 anni prima dell'anno zero, con un avanzamento dalle aree balcaniche e danubiane verso ovest e nord. In Veneto le prime evidenze agricole sono datate intorno a 8.000 anni prima dell'anno zero. Il Neolitico segna la fine del Mesolitico e l'inizio di una trasformazione progressiva del paesaggio per intervento umano diretto, inclusi il disboscamento, la pastorizia d'altura e le prime forme di arte rupestre in contesto alpino non più strettamente legato alla caccia.

\item[Paleolitico] Il periodo più lungo della preistoria umana, esteso da circa 3,3 milioni di anni fa fino alla fine dell'ultima glaciazione. È convenzionalmente suddiviso in Paleolitico inferiore (industrie ad amigdale e chopper), Paleolitico medio (industria musteriana, associata in Europa ai Neanderthal) e Paleolitico superiore (industrie lamellari, associata in Europa a \textit{Homo sapiens} a partire da circa 43.000 anni prima dell'anno zero). Nell'Italia nordorientale il Paleolitico medio è attestato nei siti della Lessinia fino a circa 38.000 anni prima dell'anno zero; il Paleolitico superiore, con la cultura epigravettiana, è presente nei siti prealpini veneti fino a circa 11.000 anni prima dell'anno zero.

\item[Paraglaciale] Aggettivo che indica i processi geomorfici instabili e accelerati che seguono il ritiro di un ghiacciaio: frane, colate detritiche, soliflusso, riassestamento delle pareti rocciose private del sostegno del ghiaccio. La causa principale è la perdita del supporto laterale e basale fornito dal ghiaccio, che lascia versanti sovra-ripiditi e materiali sciolti in condizioni di instabilità. La fase paraglaciale ha una durata variabile da decenni a millenni e si conclude quando il paesaggio raggiunge un nuovo equilibrio. Nelle Alpi, la fase paraglaciale principale si sviluppò tra circa 17.000 e 13.000 anni prima dell'anno zero, durante il ritiro dei ghiacciai tardoglaciali.

\item[Qhapaq Ñan] Rete stradale dell'impero inca (\textit{Tawantinsuyu}), che collegava tra loro le principali regioni dell'impero lungo la dorsale andina e i versanti verso il Pacifico e l'Amazzonia, per una lunghezza complessiva di oltre 30.000 chilometri. La costruzione, iniziata da stati predecessori e ampliata sistematicamente dagli Inca a partire dal XV secolo d.C., era sostenuta dal sistema di lavoro collettivo obbligatorio (\textit{mit'a}). La rete comprendeva ponti, terrapieni, gradinate e stazioni di sosta (\textit{tambo}) a intervalli regolari. Dal 2014 è riconosciuta dall'UNESCO come patrimonio mondiale dell'umanità, nella forma di un sito seriale che comprende sezioni in sei paesi sudamericani.

\item[Research bias] Distorsione sistematica che si produce quando la distribuzione geografica dei ritrovamenti scientifici rispecchia la distribuzione degli sforzi di ricerca piuttosto che la distribuzione reale del fenomeno studiato. In archeologia preistorica, si manifesta quando le aree con una tradizione di scavi sistematici mostrano densità di siti molto superiori ad aree comparabili ma meno indagate, non perché il territorio fosse effettivamente più frequentato, ma perché nessuno ha ancora cercato in modo sistematico. Il termine è usato anche in ecologia (distribuzione delle specie note), epidemiologia (incidenza di malattie diagnosticate) e altri campi in cui la conoscenza dipende dall'intensità dell'indagine.

\item[Sauveterriano] Cultura litica che contraddistingue la fase più antica del Mesolitico nell'Italia settentrionale e in buona parte dell'Europa occidentale, collocata approssimativamente tra 12.700 e 9.500 anni prima dell'anno zero. Prende il nome dal sito di Sauveterre-la-Lémance (Lot-et-Garonne, Francia). L'industria è caratterizzata da microliti geometrici di piccole dimensioni (segmenti, triangoli, punte a dorso curvo) ottenuti da supporti lamellari molto regolari. I siti sauveterriani alpini mostrano preferenza per posizioni di visibilità e controllo del territorio, in connessione con la caccia a stambecco e camoscio.

\item[Tambo] Stazione di sosta, rifornimento e controllo amministrativo disposta lungo il Qhapaq Ñan a intervalli di una giornata di cammino, generalmente ogni 20--30 chilometri. I tambo erano strutture permanenti in pietra o adobe, gestite dallo Stato inca, che fornivano alloggio ai funzionari, ai messaggeri (\textit{chasqui}) e ai contingenti militari in transito. Custodivano riserve di cibo, tessuti e altri beni ridistribuiti secondo il sistema della reciprocità asimmetrica inca. L'ubicazione dei tambo non rispondeva esclusivamente a criteri di minimizzazione del percorso, ma anche a necessità amministrative, rituali e strategiche.

\item[Tardoglaciale] Periodo di transizione tra l'Ultimo Massimo Glaciale e l'Olocene, compreso, per le Alpi, tra circa 17.000 e 9.700 anni prima dell'anno zero. È caratterizzato da un andamento climatico oscillante: una prima fase di riscaldamento (con l'oscillazione di Ragogna/Gschnitz tra 15.000 e 13.600 anni prima dell'anno zero), seguita dal miglioramento rapido del Bølling-Allerød (12.700--11.000 anni prima dell'anno zero), interrotto dal raffreddamento del Dryas Recente (11.000--9.700 anni prima dell'anno zero). Durante il Tardoglaciale il ritiro dei ghiacciai alpini rese progressivamente accessibili le quote medie e alte, permettendo la prima frequentazione stabile della fascia montana da parte di cacciatori-raccoglitori.

\item[Ultimo Massimo Glaciale (UMG)] Fase di massima estensione dei ghiacciai durante l'ultima glaciazione, collocata globalmente tra circa 26.000 e 19.000 anni prima dell'anno zero. Nelle Alpi il massimo fu raggiunto intorno a 21.000--19.000 anni prima dell'anno zero: i ghiacciai scendevano fino alle pianure pedemontane, occupando le principali valli alpine fino alle quote di uscita in pianura. Le temperature medie estive erano inferiori di 8--10°C rispetto a oggi; la linea degli equilibri glaciali si trovava circa 1.000--1.200 metri più in basso rispetto all'attuale. La fine dell'UMG, circa 19.000 anni prima dell'anno zero, segna l'inizio del ritiro glaciale e del Tardoglaciale.

\end{description}


\cleardoublepage
\addcontentsline{toc}{chapter}{Fonti}
\renewcommand{\bibname}{Fonti}
\bibliography{references}

\end{document}
